% concatenated from main.tex
\documentclass[a4paper,10pt]{article}
\usepackage{soul}
\usepackage[usenames,dvipsnames]{color}
\bibliographystyle{JHEp}
\usepackage{jheppub}
\usepackage{esvect}
\usepackage{comment}
\usepackage{amsmath, amssymb, slashed, epsf, color, graphicx, latexsym, bm, amsfonts}
\usepackage{mathrsfs}
\usepackage{tensor}
\usepackage{physics}
\usepackage{pdflscape}
\usepackage{tikz}
\usetikzlibrary{positioning, calc}
\usepackage{dsfont}
\usepackage{epsfig}
\usepackage{graphics}
\usepackage[T1]{fontenc}
\usepackage{mathtools}
\usepackage{fixme}
\fxsetup{status=draft}
%\usepackage{caption}
\usepackage{float}
\usepackage{bbm}
\usepackage{subcaption}
%\usepackage{enumitem}
\usepackage{enumerate}
\usepackage{verbatim}
\usepackage{blindtext}
\usepackage{tikz}
  \usepackage{amsmath}
\newlength\squareheight
  \setlength\squareheight{6.40pt}
  \newcommand\squareslash{\tikz{\draw (0,0) rectangle (\squareheight,\squareheight);\draw(0,\squareheight) -- (\squareheight,0)}}
  \DeclareMathOperator\squarediv{\squareslash}
\makeatletter
\renewcommand*\env@matrix[1][*\c@MaxMatrixCols c]{%
  \hskip -\arraycolsep
  \let\@ifnextchar\new@ifnextchar
  \array{#1}}
\makeatother
\DeclareFontEncoding{LS1}{}{}
\DeclareFontSubstitution{LS1}{stix}{m}{n}
\DeclareSymbolFont{symbols2}{LS1}{stixfrak}{m}{n}
\DeclareMathSymbol{\typecolon}{\mathbin}{symbols2}{"25}
\usepackage[mathlines]{lineno}
\newcommand{\C}{\mathds{C}}
\newcommand{\R}{\mathds{R}}
\newcommand{\N}{\mathds{N}}
\newcommand{\Z}{\mathds{Z}}
\newcommand{\Q}{\mathds{Q}}
\newcommand{\F}{\mathds{F}}
\newcommand{\bz}{\bar{z}}
\newcommand{\bx}{\bar{x}}
\newcommand{\s}[1]{\mathcal{#1}}
\newcommand{\ov}[1]{\overline{#1}}
\newcommand{\mr}[1]{\mathsf{#1}}
\newcommand{\tsf}[1]{\textsf{#1}}
\newcommand{\mf}[1]{\mathfrak{#1}}
\newcommand{\myRed}[1]{\textcolor{red}{#1}}
\newcommand{\qPoch}[2]{\left(#1; #2\right)_\infty}

\def\PE{\mathsf{PE}}
\def\CA{{\cal A}}
\def\CB{{\cal B}}
\def\CC {{\cal C}}
\def\CD {{\cal D}}
\def\CE {{\cal E}}
\def\CF {{\cal F}}
\def\CG {{\cal G}}
\def\CH {{\cal H}}
\def\CI {{\cal I}}
\def\CJ {{\cal J}}
\def\CK {{\cal K}}
\def\CL {{\cal L}}
\def\CM {{\cal M}}
\def\CN {{\cal N}}
\def\CO {{\cal O}}
\def\CP {{\cal P}}
\def\CR {{\cal R}}
\def\CV {{\cal V}}
\def\CW {{\cal W}}
\def\CX {{\cal X}}
\def\CO {{\cal O}}
\def\CZ {{\cal Z}}
\def\CE {{\cal E}}
\def\CG {{\cal G}}
\def\CH {{\cal H}}
\def\CI {{{\cal I}}}
\def\CB {{\cal B}}
\def\CQ {{\cal Q}}
\def\CS {{\cal S}}
\def\CT {{\cal T}}
\def\CU {{\cal U}}
\def\CX {{\cal X}}
\def\CY {{\cal Y}}
\def\CZ{{\cal Z}}

%%%%%%%%%%%%%%%%
%%%%%%%%%%%%%%%%
%%%%%%%%%%%%%%%%  Boldface blackboard font
%%%%%%%%%%%%%%%%
%%%%%%%%%%%%%%%%

\def\IA{\mathds{A}}
\def\IB{\mathds{B}}
\def\IC{\mathds{C}}
\def\ID{\mathds{D}}
\def\IE{\mathds{E}}
\def\IF{\mathds{F}}
\def\IG{\mathds{G}}
\def\IH{\mathds{H}}
\def\IJ{\mathds{J}}
\def\IK{\mathds{K}}
\def\IL{\mathds{L}}
\def\IM{\mathds{M}}
\def\IN{\mathds{N}}
\def\IO{\mathds{O}}
\def\IP{\mathds{P}}
\def\IQ{\mathds{Q}}
\def\IR{{\mathds{R}}}
\def\IS{{\mathds{S}}}
\def\IT{{\mathds{T}}}
\def\IV{{\mathds{V}}}
\def\IX{{\mathds{X}}}
\def\IZ{{\mathds{Z}}}
\newcommand{\tQ}{\widetilde{\mathcal{Q}}}
\newcommand{\tS}{\widetilde{\mathcal{S}}}

\newcommand{\dminus}{\dot{-}}
\newcommand{\dplus}{\dot{+}}
\newtheorem{assumption}{Assumption}
\begin{document}
\title{Argyres-Douglas Theories, Macdonald Indices And Arc Space Of Zhu Algebra}
\author[a]{George Andrews,}
\author[b]{Anindya Banerjee,}\author[c]{Chinmaya Bhargava,}\author[d]{Ranveer Kumar Singh,} \author[d]{Runkai Tao}
\affiliation[a]{Department of Mathematics, Penn State University, University Park, PA 16802}
\emailAdd{gea1@psu.edu}
\affiliation[b]{Department of Physics, University of Cincinnati, Cincinnati, OH 45211}
\emailAdd{ab1702@scarletmail.rutgers.edu}
\affiliation[c]{CTP and Department of Physics and Astronomy, Queen Mary University of London, London E1 4NS, UK}
\emailAdd{c.bhargava@qmul.ac.uk}
\affiliation[d]{NHETC and Department of Physics and Astronomy, Rutgers University, 126
Frelinghuysen Rd., Piscataway NJ 08855, USA}
\emailAdd{rks158@scarletmail.rutgers.edu}
\emailAdd{runkai.tao@physics.rutgers.edu}
%\emailAdd{ms3066@physics.rutgers.edu}
%\emailAdd{runkai.tao@physics.rutgers.edu}
\abstract{In this paper, we relate the MacDonald index of a 4d $\CN=2$ \tsf{SCFT} with the Hilbert series of the arc space of the Zhu algebra of the corresponding Schur \tsf{VOA}. Using this, we conjecture a simple formula for the MacDonald index of $(A_1,D_{2n+1})$ Argyres-Douglas theory. We perform checks of the formula against the known Schur limits and \tsf{RG} flows. To match the Schur limit, we prove new $q$-series identities.} 
\maketitle
\section{Introduction}
%\emph{Generic paragraph about importance of indices, applications, earlier work ...}
We study some families of 4d $\mathcal{N}=2$ superconformal field theories known as Argyres-Douglas theories \cite{Argyres:1995jj, Argyres:1995xn, Witten:1997sc, Cecotti:2010fi, Xie:2012hs}. These theories are characterized by the lack of a description in terms of any Lagrangian preserving the $\mathcal{N}=2$ supersymmetry, and fractional scaling dimensions of Coulomb branch generators (a subset of well-understood local operators that are expected to be universal in any interacting 4d $\mathcal{N}=2$ theory). Despite the absence of a Lagrangian description, these theories are amenable to powerful methods of analysis that give us insight into the dynamics of the theory without recourse to any perturbative approximation. This is made possible by several special features of these theories. These include: 
\begin{enumerate}[(1)]
\item standard simplifying features of conformal symmetry like radial quantization, the state-operator correspondence and the finite radius of convergence of the operator product expansion \cite{Rattazzi:2008pe} 
\item applicability of the representation theory of the superconformal algebra \cite{Dolan:2002zh, Cordova:2016emh} in constraining and characterizing the local operator spectrum of the theory, \item the presence of a continuous moduli space of vacua in the quantum theory that is not lifted by any radiative corrections \cite{Argyres:1996eh}, and 
\item the existence of subsectors of local operators (BPS operators) that are protected against (or transform in a tractable manner under) large classes of deformations of the theory. Such deformations include certain supersymmetry-preserving renormalization group flows (henceforth, \tsf{RG} flows) associated with spontaneous breaking of the global symmetry of these theories, and these in particular will be of interest to us \cite{Maruyoshi:2016aim, Maruyoshi:2016tqk}. 
\end{enumerate}
In this context, we will make extensive use of a familiar counting function called the superconformal index \cite{Kinney:2005ej} that counts (albeit with ambiguities) the \tsf{BPS} operators of a 4d \tsf{SCFT} and has good behavior under supersymmetry-preserving \tsf{RG} flows.
\paragraph{}
It was demonstrated in \cite{Beem:2013sza} that a subset of the ring of \tsf{BPS} local operators (called the Schur ring) has close connections to a vertex operator algebra, the algebraic building block of a two-dimensional rational conformal field theory. Throughout this work, we will focus our analysis on this subring. The superconformal index on this subring also restricts to two special limits called the MacDonald index and the Schur index \cite{Gadde:2011uv, Buican:2015tda, Cordova:2015nma} (see Appendix \ref{app:conventions} for a brief review) that we will use to understand properties of this ring. 
\paragraph{}
The goal of the present paper is twofold:
\begin{enumerate}
    \item Provide a concrete algorithm to compute the MacDonald index of a 4d $\CN=2$ \tsf{SCFT} using the Schur \tsf{VOA}. To do this, we prove, in Section \ref{sec:Mac_Hil_ser_arc}, that the Hilbert series of the arc space of the Zhu algebra of a strongly-finitely generated \tsf{VOA} is equal to the refined index of the \tsf{VOA} defined by Song \cite{Song:2016yfd}. Then, Song's conjecture implies that the MacDonald index (with flavor fugacities set to 1) of a 4d $\CN=2$ \tsf{SCFT} is the Hilbert series of the arc space of the Zhu algebra (viewed as a variety) of the corresponding Schur \tsf{VOA}. We check this result for $(A_1,A_{2n})$ Argyres-Douglas theory for which the MacDonald index is known to take a simple form. One of the main technical obstacles in the evaluation of the Hilbert series of arc space of a variety is determining the dimension of varieties appearing in the arc space construction. We use the Gr\"{o}bner basis algorithm \cite{loustaunau1994introduction} for this computation, and prove a formula for the Hilbert series in terms of monomials in the Gr\"{o}bner basis, see \eqref{eq: hilbert_ser_Rinfty_gen} for the explicit expression.
    \item We propose a simple formula for the MacDonald index of $(A_1,D_{2n+1})$ Argyres-Douglas theory (see Section \ref{sec:AD_intro} for a brief account of Argyres-Douglas theories). The main observation to expect a simple formula is that the Higgs branch \tsf{RG} flow of these theories relate it to free hypermultiplet, see Section \ref{sec:RG_flow} for more details. We can then expect that the MacDonald index can be written as some modification of the MacDonald index of the free hypermultiplet\footnote{There is a subtle issue involving a $\Z_2$ projection of the operator content of free hypermultiplet that is required, see Footnote \ref{foot:proj_hyper}.}. Following this observation and using some numerical evidence, we propose a formula for MacDonald index of the $(A_1,D_3)$ theory in Section \ref{sec:A1D3_Mac_conj} and generalize our proposal to $(A_1,D_{2n+1})$ theory in Section \ref{sec:A1D2n+1_Mac_conj}.  We then perform several nontrivial checks of our formula. In particular, to match the Schur limit to known Schur indices of the theory, we proved a family of interesting $q$-series identities. Checking our formula against the Hilbert series of the arc space of the Zhu algebra of the Schur \tsf{VOA} of these theories takes us into very techincal discussion of Gr\"{o}bner bases and constraint counting. We also check that our formula is consistent with the \tsf{RG} flow to $(A_1,A_{2(n-1)})$ theory and a free hypermultiplet. More concretely, in Section \ref{sec:prod_A2n_fH_Mac} we show that the product of the MacDonald indices of $(A_1,A_{2(n-1)})$ theory and the free hypermultiplet matches with our formula up to a correction factor.  
    \end{enumerate} 
In the following, we summarize the results of the paper for readers who are only interested in the final results. Conjectures (or results which assume another conjecture) are indicated with a question mark in parenthesis (?).
\begin{enumerate}
    \item (?) The MacDonald index of a 4d $\CN=2$ \tsf{SCFT} is the Hilbert series for the arc space of the Zhu algebra of the corresponding Schur \tsf{VOA}. 
    \item (?) The MacDonald index for $(A_1,D_{2n+1})$ theory is given by (see Appendix \ref{app:special_func} for the definition of the special functions appearing in the formulas below) 
    \begin{equation}
    \CI^{(A_1,D_{2n+1})}_{\mr{M}} = 1 + \sum_{m=1}^\infty 
    t^m \sum_{l=1}^m
    \frac{f^n_{m,l}(q)}{(q;q)_l}\sum_{k=0}^{2l}\binom{2l}{k}_q~,
\end{equation}
where $f^n_{m,l}(q)$ is a class of polynomial functions in $q$ which satisfy the hypergeometrical relation 
\begin{equation}
    1 + \sum_{m=1}^\infty q^m \sum_{l=1}^m 
    \frac{f^n_{m,l}(q)}{(q;q)_l} (q^N; q)_l 
    (q^{1-N};q)_l = q^{nN(N-1)}~.
\end{equation}
A closed form expression for the functions $f^n_{m,l}$ is given by 
\begin{equation}
f^k_{m,l}(q)=\sum_{i_k,i_{k-1},\dots,i_3\geq 0}\frac{q^{Q(m-(k-1)i_k-(k-2)i_{k-1}-\dots-2i_3-l,i_3,\dots,i_k)}(q;q)_{l}}{(q;q)_{i_k}\dots (q;q)_{i_3}(q;q)_{m-(k-1)i_k-(k-2)i_{k-1}-\dots-2i_3-l}(q;q)_{2l+i_3+2i_4+\dots+(k-2)i_k-m}}~,    
\end{equation}
where 
\begin{equation}
    Q(i_2,i_3,\dots,i_k):=i_k^2+(i_k+i_{k-1})^2+\dots+(i_k+\dots+i_3+i_2)^2~.
\end{equation}
In particular
\begin{equation}
    \CI_{\mr{M}}^{(A_1,D_3)}=\sum_{m=0}^\infty 
    \frac{t^m}{(q;q)_m} \sum_{k=0}^{2m}\binom{2m}{k}_q~.
\end{equation}
\item To match the Schur limit of our formula with the known Schur indices, we proved the following identity:
\begin{equation}
\left(\frac{\left(q^{2n+1} ; q^{2n+1}\right)_{\infty}}{(q ; q)_{\infty}}\right)^3= 1 + \sum_{m=1}^\infty 
    q^m \sum_{l=1}^m
    \frac{f^n_{m,l}(q)}{(q;q)_l}\sum_{k=0}^{2l}\binom{2l}{k}_q~.    
\end{equation}
We also proved the identity
\begin{equation}
\left(\frac{\left(q^{2} ; q^{2}\right)_{\infty}}{(q ; q)_{\infty}}\right)^3= \frac{1}{(q;q)_\infty(q^2;q^2)_\infty}\sum_{m=0}^\infty(-1)^mq^{m^2+m}\left(\frac{1+q^{2m+1}}{1-q^{2m+1}}\right)~,    
\end{equation}
of potential mathematical interest in its own right. 
\end{enumerate}
% \subsection{Outline}
% \begin{enumerate}
%     \item Introduction 
%     \begin{itemize}
%         \item Motivation
%         \item Summary of Results
%         \item Techniques: Construction of these two families of theories as generalized Class S theories: $(J^b[k],Y)$ and within Type IIB string theory as: $(G,G^{'})$.
%         \item Techniques: Index $I_M(q,t)$ from Class S and 2d TQFT. Index from Maruyoshi-Song?
%         \item $\mathcal{V}_{\text{Schur}}$ and BLLPRvR, Strongly finitely generated and Quasi-Lisse \tsf{VOA}s.\\
%         Schur index from vacuum character, Schur index from KS monodromy operator.
%         \item Higgs branch RG Flows, Partial closure of punctures, qDS reduction. 
%         \item Free field realization of \tsf{VOA}, R-filtration and Beem-Rastelli vs Buican's perspectives on the role of free fields.
%         \item Arc space of associated variety
%     \end{itemize}
%     \item Background: Zhu algebra, Arc space of associated variety, Hilbert series, MacDonald index.
%     \item $(A_1,D_3)\to(A_1,A_1)$: Results, Evidence from Schur limit, Evidence from arc space Hilbert series.
%     \\
%     A subsection on Andrews' partition identities and Chu-Vandermonde sum formulae. 
%     \item $(A_1,D_{2n+3})\to(A_1,A_{2n})$: Results, Evidence from Schur limit, Evidence from arc space Hilbert series.
%     \item Conclusions: Open Questions and Possible generalizations. 
% \end{enumerate}
\section{Constructing Argyres-Douglas Theories}\label{sec:AD_intro}
%\paragraph{}
Argyres-Douglas theories were originally discovered as special singular points on the Coulomb branch of 4d $\mathcal{N}=2$ gauge theories where mutually non-local \tsf{BPS} particles became simultaneously massless, leading to an \tsf{IR} \tsf{SCFT} that could not have a Lagrangian description preserving $\mathcal{N}=2$ supersymmetry \cite{Argyres:1995jj,Argyres:1995xn}. A second feature of these theories was that one or more Coulomb branch operator dimensions of these \tsf{SCFT}s were fractional. Later, other constructions for these theories were found. It is often the case that different observables become more tractable within different constructions of these theories. Three such constructions include: 
\begin{enumerate}[(1)]
    \item From D-branes probing isolated hypersurface singularities in Type \tsf{IIB} string theory \cite{Cecotti:2010fi}. 
    \item As generalized Class $\CS$ theories\footnote{See \cite{Witten:1997sc,Gaiotto:2009hg,Gaiotto:2009we} for discussion on class $\CS$ theories.} obtained from twisted compactification of 6d $(2,0)$ theories on punctured Riemann surfaces with defects \cite{Xie:2012hs}.
    \item As endpoints of $\mathcal{N}=1$-preserving \tsf{RG} flows with accidental supersymmetry enhancement in the \tsf{IR} \cite{Maruyoshi:2016aim, Maruyoshi:2016tqk}.
\end{enumerate}    
    Having multiple constructions has allowed us to catalogue various properties of these theories, including global symmetry (0-form and generalized), anomaly coefficients, moduli spaces of vacua, \tsf{BPS} spectrum and chiral ring, closed sub-rings of local operators like the Schur ring and the associated Schur \tsf{VOA}, and partition functions on various 4-manifolds and various choices of couplings to background fields, allowing access to information about compactifications and dualities. In this section, we collect some details about each of these constructions for the theories we will be working with and catalogue the properties that will be of interest to us.  
\subsection{Construction within Type \tsf{IIB} string theory} 
4d $\mathcal{N}=2$ \tsf{SCFT}s can be engineered within string theory in multiple ways. One way is to consider Type \tsf{IIB} string theory on a noncompact, singular Calabi-Yau threefold (the allowed types of singularities are constrained by additional consistency conditions). For our purposes, we can restrict ourselves to isolated hypersurface singularities of certain kinds that have an A-D-E classification. More concretely, the isolated hypersurface singularities are given by $\C^2/\Gamma$ where $\Gamma\subset \mr{SU(2)}$ is a discrete subgroup. Such discrete subgroups have an \tsf{ADE} classification. For $\mf{g}\in\{A_n,D_n,E_{6,7,8}\}$, let $\Gamma_\mf{g}\subset\mr{SU(2)}$ be the discrete subgroup and let $f_{\mf{g}}(x,y)$ be a polynomial such that  $\C^2/\Gamma_{\mf{g}}$ is the zero set of $f_{\mf{g}}$. The corresponding singular Calabi-Yau threefold is given by the zero set of 
\begin{equation}
x^2+y^2+f_{\mf{g}}(z,w)~.    
\end{equation}
Explicit form of the polynomials is shown in Table \ref{tab:pol_sing_CY}.
\begin{table}[]
\centering
\begin{tabular}{|c|l|}
\hline
\(\mathfrak{g}\) & \(f_{\mathfrak{g}}(x,y)\) \\
\hline
\(A_n\)   & \(x^2 + y^{n+1} \) \\
\hline
\(D_n\)   & \(y\,(x^2 + y^{n-2}) \) \\
\hline
\(E_6\)   & \(x^3 + y^4 \) \\
\hline
\(E_7\)   & \(x\,(x^2 + y^3) \) \\
\hline
\(E_8\)   & \(x^3 + y^5 \) \\
\hline
\end{tabular}
\caption{Polynomials defining singular Calabi-Yau threefolds}
\label{tab:pol_sing_CY}
\end{table}
% \begin{equation}
%     \begin{split}
%         &A_n: \text{ }x^2+y^{n+1}=0~, \\
%         &D_n: \text{ }y(x^2+y^{n-2})=0~, \\
%         &E_6: \text{ }x^3+y^4=0~, \\
%         &E_7: \text{ }x(x^2+y^3)=0~, \\
%         &E_8: \text{ }x^3+y^5=0~, \\
%     \end{split}
% \end{equation}
\paragraph{$(A_1,A_{N-1})$ Theory.}
The $(A_1,A_{N-1})$ theories are constructed using the singular Calabi-Yau given by \cite{Cecotti:2010fi}
\begin{equation}
    x^2+y^2+z^2+w^{N}=0~.
\end{equation}
This allows us to define the Seiberg-Witten curve $\Sigma$ and differential $\lambda$ as,
\begin{equation}
    \Sigma:~~z^2+w^{N}=0, \quad \lambda=z\,dw~.
\end{equation}
The above form allows us to read off the scaling dimensions of $z$ and $w$ as follows:
\begin{equation}
    \begin{split}
        &[z]+[w]=[\lambda]=1~,\quad 2[z]=N[w]  \\
    &\implies [z]=\frac{N}{N+2}, \quad [w]=\frac{2}{N+2}~.
    \end{split}
\end{equation}
Consider the simplest case of the $(A_1,A_2)$ theory. The Coulomb branch spectrum can be read off from the following deformed curve:
\begin{equation}
    z^2+w^3=u+\cdots
\end{equation}
where $u$ is the unique relevant operator appearing on the \tsf{RHS} since any other coefficient $v$ in a monomial involving $z$ and $w$ would have $[v]<1$ and therefore would not correspond to an operator in a unitary 4d $\mathcal{N}=2$ \tsf{SCFT}. 
This gives the following scaling dimensions:
\begin{equation}
    [z]=\frac{3}{5}, \quad [w]=\frac{2}{5}, \quad [u]=\frac{6}{5}~.
\end{equation}
Generalizing the above to arbitrary $N$, we find the following spectra of Coulomb branch operators:
\begin{equation}
    \begin{split}
        &\Delta\in\left\{\frac{2N}{N+2},\frac{2N-2}{N+2},\frac{2N-4}{N+2},\cdots,\frac{N+3}{N+2}\right\},\quad  N\text{ odd}~, \\
        &\Delta\in\left\{\frac{2N}{N+2},\frac{2N-2}{N+2},\frac{2N-4}{N+2},\cdots,\frac{N+4}{N+2}\right\},\quad \text{ }N\text{ even}~.
    \end{split}
\end{equation}
% =\left\{\frac{2N-2k}{N+2}:k=0,1,\cdots,\frac{N-3}{2}\right\}=\left\{\frac{2N-2k}{N+2}:k=0,1,2,\cdots,\frac{N-4}{2}\right\}
We therefore see that the rank (complex dimension of its Coulomb branch) of the $(A_1,A_{N-1})$ \tsf{SCFT} is $r=\lfloor\frac{N-1}{2}\rfloor$. We will be interested in the cases with $N$ odd above. Denoting $N=2n+1$ with $n\geq 1$, we can write down the central charges for these theories,
\begin{equation}
    c=\frac{n(6n+5)}{6(2n+3)}, \quad a=\frac{n(24n+19)}{24(2n+3)}~.
\end{equation}
These central charges may be derived from the Shapere-Tachikawa formula \cite{Shapere:2008zf}: 
\begin{equation}
    2a-c=\frac{1}{4}\sum_{i=1}^{r}(2\Delta_i-1)~,
\end{equation}
and the modular properties of the Schur index\footnote{The Schur index of a 4d 
$\CN=2$ \tsf{SCFT} counts certain class of operators in the theory, see Appendix \ref{app:conventions} for more details. In Type \tsf{IIB} construction, the Schur index can be obtained by computing the trace of the Kontsevich-Soibelman wall-crossing invariant $\mathcal{I}_{\mr{S}}(q)=\mr{Tr}~\mathcal{O}(q)$ \cite{Cordova:2015nma,Cecotti:2010fi}.}: 
\begin{equation}
    \lim_{q\to 1}\text{log }\mathcal{I}_{\mr{S}}(q)\approx \frac{8\pi^2}{\beta}\cdot (c-a), \quad  \quad q=e^{-\beta}~.
\end{equation}
These theories have trivial flavor symmetry for all $n$. 
\paragraph{$(A_1,D_N)$ Theory.}
The $(A_1,D_N)$ Argyres-Douglas theories are obtained from the hypersurface \cite{Cecotti:2010fi} 
\begin{equation}
x^2+y^2+wz^2+w^{N-1}=0~.    
\end{equation}
This implies the scaling dimensions
\begin{equation}
    \begin{split}
    [z]=\frac{N-2}{N}, \quad [w]=\frac{2}{N}~.
    \end{split}
\end{equation}
For the case of $(A_1,D_3)$, this leads to the following deformed Seiberg-Witten curve:
\begin{equation}
    \Sigma:~~w^2+wz^2=u~,
\end{equation}
which gives us the following scaling dimensions:
\begin{equation}
    [w]=\frac{2}{3}, \quad [z]=\frac{1}{3}, \quad [u]=\frac{4}{3}~.
\end{equation}
For general $N$, we obtain the following spectrum of Coulomb branch dimensions \cite{Cecotti:2010fi}:
\begin{equation}
    \begin{split}
        &\Delta\in\left\{\frac{2N-2}{N},\frac{2N-4}{N},\frac{2N-6}{N},\cdots,\frac{N+1}{N}\right\}, \quad  N\text{ odd}~, \\
        &\Delta\in\left\{\frac{2N-2}{N},\frac{2N-4}{N},\frac{2N-6}{N},\cdots,\frac{N+2}{N}\right\},\quad  N\text{ even}~, \\
    \end{split}
\end{equation}
% =\{\frac{2N-2k}{N}\}\text{ }k=\{1,2,\cdots,\frac{N-1}{2}\},\text{ }=\{\frac{2N-2k}{N}\right\}k=\{1,2,\cdots,\frac{N-2}{2}\}, \text{ }
so that the rank of the $(A_1,D_N)$ \tsf{SCFT} is $\lfloor\frac{N-1}{2}\rfloor$.  We will be interested in the cases with $N$ odd. Denoting $N=2n+1$ with $n\geq 1$, we can also use the Type \tsf{IIB} construction to infer the central charges for these theories:
\begin{equation}
    c=\frac{n}{2}, \quad a=\frac{n(8n+3)}{8(2n+1)}~.
\end{equation}
These theories generically have $\mr{U(1)}$ flavor symmetry, which enhances to $\mr{SU(2)}$ for $n=1$. 
\subsection{Construction As Generalized Class $\CS$ Theories}
A second construction of these theories that will be useful for us is to start from the 6d $(2,0)$ theory denoted as $\mathcal{S}[\mf{g}]$ (where $\mf{g}=\mf{a,d,e}$ is a simply laced finite-dimensional Lie algebra) and compactify on a Riemann surface with regular and irregular punctures. The theories that form the building blocks of this family of theories involve compactifications on the two-punctured Riemann sphere with one irregular puncture and one regular puncture. The 6d theory has codimension 2 defects that live in the directions transverse to these punctures at the singular locus. In the case where $\mf{g}=\mf{a}_{N-1}$ (which is what we will be concerned with), regular punctures are labeled by the data of a Young diagram $Y$ of $\mf{a}_{N-1}$ consisting of $N$ boxes \cite{Gaiotto:2009we,Chacaltana:2010ks,Tachikawa:2013kta,LeFloch:2020uop}. This corresponds to an irreducible representation of $\mf{a}_{N-1}$, and describes a boundary condition at the puncture for the Higgs field $\Phi(z)$ of the Hitchin system living on the Riemann surface. A full puncture $Y=\mr{F}$ corresponds to the Young diagram labeled $[1^N]$ in our convention, i.e., with a single row containing $N$ boxes. A simple puncture $Y=\phi$ corresponds to the Young diagram denoted as $[N]$ in our convention, and it refers to the Young diagram with $N$ rows with one box each, or the trivial irreducible representation of $\mf{a}_{N-1}$. The irregular puncture is labeled by two integers $b$ and $p$ which refer to the pole structure of the singular Higgs field. Together, this allows a packaging of the data defining the construction into the form $D_p^b(\mf{g},Y)$, which labels the 4d $\mathcal{N}=2$ \tsf{SCFT}. Some well-studied families of Argyres-Douglas theories that will be relevant to our work, constructed as generalized Class $\CS$ theories, have the following parameters in this notation \cite{Beem:2023ofp}:
\begin{equation}
    \begin{split}
        (A_1,D_{2N+2})&=D_{2N+2}^2(\mf{sl}_2,[1^2])~, \\
        %=(A_{N+1}^{\text{ }N+2}[-1],(N,1^2)) ~, \\
        (A_1,A_{2N-1})&=D_{2N+2}^2(\mf{sl}_2,[2])~, \\
        %=(A_{N-1}^{\text{ }N}[1],(N-1,1))~, \\
        (A_1,D_{2N+3})&=D_{2N+3}^2(\mf{sl}_2,[1^2])~, \\
        %=(A_1^{\text{ }2}[2N+1],(2)) \\
        (A_1,A_{2N})&=D_{2N+3}^2(\mf{sl}_2,[2])~. \\
        %=(A_1^{\text{ }2}[2N+1],(1^2))~. \\
    \end{split}
\end{equation}
This construction is useful for deriving the MacDonald index (see Appendix \ref{app:conventions} for a brief account of MacDonald index) of these theories. For example, the MacDonald index for $(A_1,D_{2N-3})=D_{2N-3}^2(\mf{sl}_2,[1^2])$ theory is given by \cite{Buican:2015tda}:
\begin{equation}\label{eq:Mac_class_S-1}
    I_M(q,t,a)=\sum_{n=0}^{\infty}C_n\tilde{F}_n^{(N)}(q,t,a)~,
\end{equation}
where
\begin{equation}
    C_n=\bigg[\frac{(t;q)_{\infty}}{(q;q)_{\infty}} \bigg]^{1/2}\frac{P_n(q,t,\sqrt{t})}{(t^2;q)_{\infty}}~,
\end{equation}
and
\begin{equation}
    \tilde{F}_n^{(N)}(q,t,a)=\bigg[\frac{(t;q)_{\infty}}{(q;q)_{\infty}} \bigg]^{1/2}\frac{\tilde{P}_n^{(N)}(q,t,a)}{(t;q)_{\infty}}~,
\end{equation}
is the wavefunction associated to the irregular puncture. The special functions $P_n(q,t,a)$ and $\tilde{P}_n^{(N)}(q,t,a)$ respectively have the expressions
\begin{equation}
    \begin{split}
        &P_n(q,t,a)=N_n(q,t)\sum_{m=0}^{n}\frac{(t;q)_m(t;q)_{n-m}}{(q;q)_m(q;q)_{n-m}}a^{2m-n}~, \\
        &\tilde{P}_n^{(N)}(q,t,a)=N_n(q,t)\sum_{m=0}^n\frac{(t;q)_m(t;q)_{n-m}}{(q;q)_m(q;q)_{n-m}}a^{2m-n}q^{-N(\frac{n}{2}-m)^2}~,
    \end{split}
\end{equation}
and the normalization factor $N_n(q,t)$ is given as
\begin{equation}\label{eq:Mac_class_S-7}
    N_n(q,t)=\frac{(q;q)_n}{(t;q)_{\infty}}\sqrt{(1-q^nt)(q^{n+1};q)_{\infty}(q^nt^2;q)_{\infty}}
\end{equation}
For completeness, we also note that the wavefunction associated with a regular puncture is given as
\begin{equation}
    F_n(q,t,a)=\bigg[\frac{(t;q)_{\infty}}{(q;q)_{\infty}} \bigg]^{1/2}\frac{P_n(q,t,a)}{(ta^{\pm 2,0};q)_{\infty}}~.
\end{equation}

\subsection{Higgs Branch \tsf{RG} Flow And Quantum Hamiltonian Reduction}\label{sec:RG_flow}
The construction of the $(A_1,A_N)$ and the $(A_1,D_N)$ theories within the Class $\CS$ framework makes it transparent that the $(A_1,D_{2n+3})=D_{2n+3}^2(\mf{sl}_2,[1^2])$ theory flows to the $(A_1,A_{2n})=D_{2n+3}^2(\mf{sl}_2,[2])$ theory under a Higgs branch \tsf{RG} flow, implemented as a partial closure of the regular puncture from $Y=\mr{F}=[1^2]$ to $Y=\phi=[2]$. We briefly review the procedure of partial closure of the puncture applied for our case and refer the reader to \cite{Tachikawa:2015bga} for more details. Let us first recall the notion of Higgsing. Suppose a 4d $\CN=2$ \tsf{SCFT} with flavor symmetry group $G$. The flavor multiplet has a chiral operator $\mu^+$ in the adjoint of $G$. One can then give a nilpotent \tsf{VEV} to this operator which triggers a Higgs branch flow. By Jacobson-Morozov theorem, any nilpotent element of a semi-simple Lie
algebra $\mf{g}$ is given by the image $\rho(\sigma^+)$ of the raising operator $\sigma^+\in\mf{su}(2)$ of an embedding $\rho: \mf{su}(2) \longrightarrow \mf{g}$. The flavor symmetry algebra of the \tsf{IR} fixed point\footnote{Note that, in general, this procedure does not preserve the full $\mr{SU(2)_R}$-symmetry of the \tsf{UV} theory and is broken to $\mr{U(1)_R}$.} then is the commutant of $\rho(\sigma^+)$ in $\mf{g}:=\mr{Lie}(G)$. By the usual Higgs mechanism, this sponteneous breaking of flavor symmetry results in some Nambu-Goldstone hypermultiplets at the \tsf{IR} fixed point. 
\par In a class $\CS$ theory, flavor symmetries are associated to the defects at the punctures of the Riemann surface. In case of a class $\CS$ theory of type $A_{n-1}$, Higgsing of the flavor symmetry associated to a regular puncture is called partial closure of puncture. The $(A_1,D_{2n+3})$ Argyres-Douglas theory has $\mr{SU(2)}$ flavor symmetry. The chiral operator $\mu^+$ is in the adjoint of $\mr{SU(2)}$. The only choice of nilpotent \tsf{VEV} for $\mu^+$ is
\begin{equation}
    \langle\mu^+\rangle=\sigma^+~.
\end{equation}
The flavor symmetry then breaks from $\mr{SU(2)}$ to the commutant of $e^{i\sigma^+}$ which is $\mr{U(1)}$. At the level of defect at the puncture, the nilpotent \tsf{VEV} $\langle\mu^+\rangle=\sigma^+$ closes the puncture $[1^2]$ to $[2]$. The \tsf{IR} fixed point\footnote{One can show that this partial closure of puncture the \tsf{UV} $\mr{SU(2)_R}$-symmetry breaks to $\mr{U(1)_R}$ in the \tsf{IR} but sits as the Cartan of an $\mr{SU(2)'_R}$ which becomes the \tsf{R}-symmetry of the \tsf{IR}. The \tsf{IR} fixed point is thus a 4d $\CN=2$ \tsf{SCFT}. In the language of Tachikawa \cite{Tachikawa:2015bga}, this partial closure of puncture is \textit{good}.} is thus the $(A_1,A_{2n})$ Argyres-Douglas theory with flavor symmetry $\mr{U(1)}$. Following the counting of free Nambu-Goldstone hypermultiplets in \cite{Tachikawa:2015bga}, we see that the partial closure of puncture $[1^2]\to[2]$ results in a single free hypermultiplet. We thus conclude that the Higgs branch \tsf{RG} flow of the $(A_1,D_{2n+3})$ theory results in a single free hypermultiplet and the $(A_1,A_{2n})$ theory. We summurize this as 
\begin{equation}\label{eq:RG_flowD2n+1to2n}
(A_1,D_{2n+3})\xrightarrow[\tsf{RG flow}]{\langle\mu^+\rangle=\sigma^+} (A_1,A_{2n})\otimes 1\,\mr{HM}~.   
\end{equation}
At the level of the corresponding Schur \tsf{VOA}s, the analogous operation is the quantum Drinfeld-Sokolov (or quantum Hamiltonian) reduction (see Section 4 of \cite{Beem:2014rza}). This is a special case of the \tsf{BRST} procedure applied to \tsf{VOA}s, typically affine Lie \tsf{VOA}s $V_k(\mathfrak{g})$. It requires the input data of the affine \tsf{VOA} $V_k(\mathfrak{g})$ and a Lie algebra homomorphism $\rho:\mf{su}(2)\longrightarrow\mathfrak{g}$. The output is a second \tsf{VOA}, typically a $\mathcal{W}$-algebra obtained by ``gauging'' or the imposition of certain constraints (related to the nilpotent orbit $\rho$) on the parent \tsf{VOA} by the \tsf{BRST} procedure. Additional useful properties appear (e.g., modular properties of the characters of modules) when the parent affine \tsf{VOA} $V_k(\mathfrak{g})$ are at \textit{boundary admissible levels} in the sense of Kac and Wakimoto, i.e., there exists $q\geq h^\vee$ with $\mr{gcd}(q,h^\vee)=1$ such that
\begin{equation}
    \tilde{k}=-h^{\vee}+\frac{h^{\vee}}{q}~, 
\end{equation}
where $h^\vee$ is the dual Coxeter number of $\mf{g}$. For the \tsf{RG} flow \eqref{eq:RG_flowD2n+1to2n}, the quatum Drinfeld-Sokolov reduction is \cite{adamovic2019realizations}
\begin{equation}
    \mf{su}(2)_{-\frac{4n+4}{2n+3}}\to\mr{Vir}\left(-\frac{2n(6n+5)}{2n+3},0\right)~.
\end{equation}
Based on this, we expect that the indices in the theories related by \tsf{RG} flow \eqref{eq:RG_flowD2n+1to2n} will have some factorization property\footnote{See \cite{Deb:2025ypl} for an intriguing relation of Schur indices of theories lying on the same \tsf{RG} flow line.}. Indeed, we show in Section \ref{sec:RG_flow} that such a factorization of MacDonald index holds. 
\par
It has been suggested by the authors of \cite{Beem:2019tfp} that an inversion of the quantum Drinfeld-Sokolov reduction is implemented by their free field realizations\footnote{See also \cite{Beem:2024fom} for free-field realizations of rank 1 \tsf{SCFT}s.} for the corresponding Schur \tsf{VOA}s. This aligns with the spirit of this work, where we use analogous machinery to these free field realizations to come up with a new proposal for the MacDonald indices of these theories. Note that while the free field realizations of \cite{Beem:2019tfp} come with an R-filtration and therefore can be used to obtain the MacDonald index, our techniques do not seem to derive in any obvious way from this R-filtration, and the operator rearrangement scheme we work out here seems to have a different (though assuredly related) origin. \\
% \subsection{The MacDonald index from the Class S construction}
% The $(A_1,D_{2N-3})$ theory may be constructed as a generalized Class S theory by compactifying the 6d $(2,0)$ theory of Type $A_1$ on a Riemann sphere with one ``irregular" puncture. In the notation of (\textcolor{red}{Cite hep-th/2311.12123}), this is denoted as a $D_p^b(G,Y)$ theory with $G=A_1$, $b=2$, $p=2N-3$ and $Y=F$ (the full puncture corresponding to the $A_1$ Young diagram with one row and two columns or the irreducible representation of dimension 2). The MacDonald index of this theory is then given as ,

% \subsection{Goals}
% We have found a relationship between the MacDonald indices of the $(A_1,D_{2n+3})$ and the $(A_1,A_{2n})$ series of Argyres-Douglas theories. We want to understand the physical content of our results. \\
% \textbf{Basic Goal}: \emph{To add a story to support our computations that would justify submitting these results to a physics journal}.
% \begin{enumerate}
%     \item The RG flow triggered by flowing onto the Higgs branch of the $(A_1,D_{2n+1})$ theory spontaneously breaks the $\mathfrak{f}=\text{su}(2)$ flavor symmetry of the UV \tsf{SCFT}, and we discover an additional free hypermultiplet representing Goldstone modes in the IR besides the $(A_1,A_{2n})$ \tsf{SCFT}, as expected. However, this says nothing about the global form of the flavor symmetry. In particular, our expression for the UV index in terms of IR fields makes use of a $\mathds{Z}_2$ projection ($\mathds{Z}_2$ gauging?) of the free hypermultiplet, rather than the usual free hypermultiplet. How do we interpret this? \\
%     \emph{Describe the global form of the flavor symmetry in the UV and the IR \tsf{SCFT}s and use this to understand whether our index relations are the consequence of an RG flow or an operator mapping between the MacDonald rings of the two sets of theories.}
%     \item This particular RG flow is special in multiple ways, which may be why we have not managed to find other families of theories where such a relationship exists between the MacDonald indices. We note some of these special features below and ask: \\
%     \emph{At what stages in our computations do these features come into play, if at all?}
%     \begin{itemize}
%         \item The $(J^b[k],Y)$ notation for a 4d $\mathcal{N}=2$ \tsf{SCFT} of Class $\mathcal{S}$ amounts to the following data: We compactify the 6d $(2,0)$ theory of type $J\in A-D-E$ on a Riemann sphere with one regular puncture and one irregular puncture. A Hitchin system lives on this Riemann sphere. At the regular puncture, the singularity of the Higgs field of the Hitchin system is given by the Young diagram $Y$ of $J$, while at the irregular puncture, the singularity of the Higgs field is described by the integers $(b,k)$. In this notation, we have $(A_1,D_{2n+3})=(A_1^2[2n+1],(1^2))$ while $(A_1,A_{2n})=(A_1^2[2n+1],(2))$, so in the Class $\mathcal{S}$ description, this RG flow is implemented as a partial closure of the full $A_1$ puncture to the trivial $A_1$ puncture. For comparison, there exists a Higgs branch RG flow between $(A_1,D_{2n+2})=(A_{n+1}^{n+2}[-1],(n,1^2))$ and $(A_1,A_{2n-1})=(A_{n-1}^n[1],(n-1,1))$, but in the Class $\mathcal{S}$ description, this is not as simple as closing an $A_1$ regular puncture. \\
%         \textbf{Reference:} \emph{hep-th/1706.01607, Song, Xie, Yan. Also, Xie's original paper on generalized Argyres-Douglas theories}.
%         \item The Hall-Littlewood index of the $(A_1,D_{2n+3})$ theories are entirely generated by $\hat{\mathcal{B}}_1$ multiplets, while the $(A_1,A_{2n})$ theories have a trivial Hall-Littlewood chiral ring. On the other hand, in the $(A_1,D_{2n+2})$ theory, the Hall-Littlewood index is generated by four $\hat{\mathcal{B}}_1$ multiplets and four $\hat{\mathcal{B}}_{(n+1)/2}$ multiplets for $n>0$. The index itself is given by,
%         \begin{equation}
%             \frac{1+2t^{(n+1)/2}+t^{n+1}+2t^n+2t^{n-1}+\cdots+2t}{(1-t)^2(1-t^{(n+1)/2})^2}
%         \end{equation}
%         On the other hand, the Hall-Littlewood chiral ring of the $(A_1,A_{2n-1})$ theory is generated by one $\hat{\mathcal{B}}_1$ multiplet $M$ and two $\hat{\mathcal{B}}_{n/2}$ multiplets $B_1$ and $B_2$ with a single relation $B_1B_2+M^n=0$, and so the Higgs branch of the theory is,
%         \begin{equation}
%             \mathcal{M}_{\text{Higgs}}=\mathds{C}[M,B_1,B_2]/\langle B_1B_2+M^n\rangle\cong\mathds{C}^2/\mathds{Z}_{n}
%         \end{equation}
%         and the Hall-Littlewood index is,
%         \begin{equation}
%             I_{HL}(t)=\frac{1-t^n}{(1-t)(1-t^{n/2})^2}=\frac{1+t^{n/2}}{(1-t)(1-t^{n/2})}
%         \end{equation}
%         We see that not only are the Hall-Littlewood indices not the same, one also cannot make up the difference with free hypermultiplets because of the presence of generators with $R>1$. At the level of the Schur \tsf{VOA}s of these theories, this difference manifests in the presence of only conformal dimension $h\leq 2$ generators for both the $\hat{su}(2)_{-\frac{4n}{2n+1}}$ and the $\mathcal{M}(2,2n+3)$ Virasoro minimal models. On the other hand, for both the $\mathcal{B}_N$ Bershadsky-Polyakov algebras and the logarithmic $\mathcal{W}_M$ algebras (which are the Schur \tsf{VOA}s of the $(A_1,A_{2n-1})$ and the $(A_1,D_{2n+2})$ theories for specific values of $N$ and $M$), there exist $\mathcal{W}$-algebra generators with $h>2$.
%         \item The Higgs branch of the $(A_1,D_{2n+3})$ theories as a variety is $\mathds{C}[x,y,z]/\langle xz+y^2\rangle\cong\mathds{C}^2/\mathds{Z}_2$ and this is the principal nilpotent orbit of $\text{SU}(2)$. On the other hand, for the $(A_1,D_{2n+2})\to (A_1,A_{2n-1})$ RG flows, even in the one case where the Higgs branch is generated entirely by $\hat{\mathcal{B}}_1$ multiplets ($n=1$), the UV Higgs branch is not the principal nilpotent orbit of $\text{SU}(3)$ but a sub-maximal one. The difference is manifested in the dimension of the variety,
%         \begin{equation}
%             \begin{split}
%                 &(1^3)\text{ or principal nilpotent orbit}\to\frac{(1-t^2)(1-t^3)}{(1-t)^8}, \quad \mr{dim}_{\mathds{C}}\mathcal{V}=6 \\
%                 &(1,2)\text{ or minimal nilpotent orbit}\to\frac{1+4t+t^2}{(1-t)^4}, \quad \mr{dim}_{\mathds{C}}\mathcal{V}=4 \\
%             \end{split}
%         \end{equation}
%         \textbf{Reference:} \emph{hep-th/1601.0402, Hanany, Kalveks; Table 4.}
%     \end{itemize}
%     \item Explain the partition theory origin of the identity Andrews used to prove the $q$-series equivalence (Chu-Vandermonde sums). Maybe add a section on $q$-hypergeometric series and why they play a role in 4d $\mathcal{N}=2$ superconformal indices (Hitchin integrable systems are supposed to be the connecting link, expand on this). 
% \end{enumerate}
% % \section{4d $\mathcal{N}=2$ superconformal index}
% % We denote the supercharges of 4d $\mathcal{N}=2$ theory as 
% % \begin{equation}
% %     \begin{array}{llll}
% % \mathcal{Q}_\alpha^{\mathcal{I}}\quad & \widetilde{\mathcal{Q}}_{\mathcal{I} \dot{\alpha}}\quad & \mathcal{S}_{\mathcal{I}}^\alpha
% % \equiv (\mathcal{Q}^\dagger)^\alpha_\mathcal{I}
% % \quad& \widetilde{\mathcal{S}}^{\mathcal{I} \dot{\alpha}}
% % \equiv 
% % (\tQ^\dagger)^{\mathcal{I}\dot\alpha} 

% % ,
% % \end{array}
% % \end{equation}
% % where $\alpha =\pm$ is an 
% % $SU(2)_1$ index and 
% % $\dot\alpha = \dot{\pm}$ an
% % $SU(2)_2$ index and $\mathcal{I} = 1,2$ and $SU(2)_R$ index. 
% % In our convention, $\mathcal{Q}$ has $U(1)_r$ charge $-\frac{1}{2}$ and 
% % $\tQ$ has $U(1)_r$ charge $\frac{1}{2}$.

% % The 4d $\mathcal{N}=2$ superconformal index is defined to be 
% % \begin{equation}
% % \label{def:4dN=2scI}
% %     \mathcal{I}(p,q,t; a_i) = \mathrm{Tr}(-1)^F 
% %     t^{r+R} p^{j_2 - j_1 -r}
% %     q^{j_1+j_2 - r}
% %     \prod_{i}a_i^{f_i} 
% %     e^{-\beta \delta_{2\dot{-}}}
% % \end{equation}
% % where 
% % \begin{equation}
% %     2\delta_{2\dminus} = 
% %     \{ \tQ_{2\dminus}, (\tQ_{2\dminus})^\dagger\} =  
% %     E - 2j_2 - 2R + r.
% % \end{equation}
% % $(p,q,t)$ are superconformal fugacities, and $\{a_i\}$ are flavor fugacities. $E$ is the conformal dimension and $j_1$ and $j_2$ are the Cartan generators of the $SU(2)_1 \times SU(2)_2$.
% % $\{f_i\}$ are the flavor symmetry generators. 
% % The definition \eqref{def:4dN=2scI} is believed to be independent of $\beta$. That is the states that contribute to the index satisfy 
% % $\delta_{2\dminus} = 0$.
% % Since $\tQ_{2\dminus}$ has quantum numbers 
% % $E = R = - r = \frac{1}{2}$,
% % $(j_1,j_2) = (0,-\frac{1}{2})$, 
% % $r+R$, $j_2 -j_1 -r$ and 
% % $j_1+j_2 -r$ all commute with $\delta_{2\dminus}$.
% % \subsection{Supersymmetric limits}
% % \noindent {\bf Macdonald limit}\\
% % In the Macdonald limit, we take $p\to 0$.

% % \dots

% % \noindent {\bf Hall-little wood limit}\\
% % In the Hall-little wood limit, we take $p,q \to 0$.

% % \dots 

% % \noindent {\bf Schur limit}\\
% % In the Schur limit, we take $t\to q$.
% % The superconformal index \eqref{def:4dN=2scI} in this limit becomes 
% % \begin{equation}
% %     \mathcal{I}\left(p, q ; a_i\right)=\operatorname{Tr}(-1)^F  p^{j_2-j_1-r} q^{j_1+j_2 + R} \prod_i a_i^{f_i} e^{-\beta \delta_2-}
% % \end{equation}
% % \myRed{The Schur index will be independent of $p$.}

% % \noindent {\bf The Coulomb branch limit}
% % We take $p,q,t \to 0$ while 
% % $pq/t = u$ fixed.





% \section{Superconformal index for $(A_1,A_{2N-1})$ theory}
% \begin{equation}
%     \begin{aligned}
% \mathcal{I}_{\mathcal{N}=2}^{\left(A_1, A_{2 N-1}\right)}= & \prod_{i=1}^{N-1} \frac{\Gamma\left(\left(\frac{p q}{t}\right)^{\frac{2 N+1-i}{N+1}}\right)}{\Gamma\left(\left(\frac{p q}{t}\right)^{\frac{i+1}{N+1}}\right)} \\
% & \times \frac{\kappa^{N-1}}{N!} \Gamma\left(\left(\frac{p q}{t}\right)^{\frac{1}{N+1}}\right)^{N-1} \oint[d \boldsymbol{z}] \prod_{i \neq j} \frac{\Gamma\left(\frac{z_i}{z_j}\left(\frac{p q}{t}\right)^{\frac{1}{N+1}}\right)}{\Gamma\left(\frac{z_i}{z_j}\right)} \prod_{i=1}^N \Gamma\left(\left(z_i a\right)^{ \pm}\left(\frac{p q}{t}\right)^{\frac{1-N}{2 N+2}} t^{\frac{1}{2}}\right)
% \end{aligned}
% \end{equation}
%\section{Introduction}
\section{MacDonald Index And Zhu Algebra}\label{sec:Mac_Hil_ser_arc}
In \cite{Song:2016yfd}, Song conjectured an algorithmic procedure to compute the MacDonald index of a 4d $\CN=2$ theory using only the data of the corresponding \tsf{VOA}. \cite{Song:2016yfd} defined a certain refined character for a strongly finitely generated \tsf{VOA} (see \eqref{eq:ref_char_song}) and checked that it reproduces the MacDonald index for a large class of 4d $\CN=2$ \tsf{SCFT}.   In this section, we prove the following theorem.
\begin{thm}\label{thm:Hil_ser_mac}
Let $\CV$ be a strongly finitely generated $\mr{VOA}$. Then the refined character of $\CV$ with the fugacities for (non-affine) symmetry $\mf{g}$ set to 1 is equal to the Hilbert series for the arc space of the Zhu algebra of $\CV$. In particular, assuming Song's conjecture, 
the MacDonald index without flavor fugacities of a 4d $\CN=2~\mr{SCFT}$ is given by the Hilbert series of the arc space of the Zhu algebra of the corresponding chiral algebra.     
\end{thm}
%Our argument is based on the conjecture that the MacDonald index of a 4d $\CN=2$ \tsf{SCFT} is given by a certain refined vacuum character (defined below) of the corresponding \tsf{VOA}\cite{Song:2016yfd}. 
% \begin{enumerate}
%     \item The 
%     \item The Zhu algebra (defined below) of a $\mathcal{W}$-algebra $\CV$ is given by \cite{Xie:2019zlb} 
%     \begin{equation}
%         A(\CV)\cong\frac{\C[W_1,\dots,W_n]}{\langle f_1,\dots,f_n\rangle}~,
%     \end{equation}
% where $W_i$ are variables corresponding to generators of conformal dimension $d_i$ of the $\mathcal{W}$-algebra and the relations $f_i$ correspond to the singular vector relations among the generator.      
% \end{enumerate} 
\subsection{The Refined Character Of A \tsf{VOA}}
We begin by recalling the notion of the refined index of a \tsf{VOA}. Recall that a \tsf{VOA}$\CV$ is said to be \textit{strongly generated} by a subset $U\subseteq \CV$ if  
\begin{equation}
    \CV=\mathsf{Span}\left\{u^1_{(-n_1)}\cdots u^k_{(-n_k)}|0\rangle:u^i\in U,n_i\geq 1,k\geq 0\right\}/\{\mr{null~~ vectors}\}~,
\end{equation}
where $k=0$ corresponds to the vacuum $\ket{0}$ and our convention for the expansion of the vertex operators is 
\begin{equation}
    Y(u,z)=\sum_{n\in\Z}u_{(-n)}z^{-n-1}~.
\end{equation}
If $U$ is a finite-dimensional subspace, then $\CV$ is called \textit{strongly finitely generated}. Note that in this convention, we have 
\begin{equation}
    L_0\cdot u^1_{(-n_1)}\cdots u^k_{(-n_k)}|0\rangle=\left(\sum_{i=1}^k(d_i+n_i-1)\right)u^1_{(-n_1)}\cdots u^k_{(-n_k)}|0\rangle~,
\end{equation}
where $d_i$ is the conformal dimension of $u^i$.


The notion of the refined character, as we will see below, requires the \tsf{VOA} to be strongly generated. It is widely believed that the \tsf{VOA} arising from a 4d $\CN=2$ \tsf{SCFT} is strongly finitely generated \cite{Beem:2019tfp,Arakawa:2016hkg}. We will need this assumption to prove our statement.

A \tsf{VOA} $\CV$ strongly generated by $U$ admits the filtration
\begin{equation}
    \CV_0\subset\CV_1\subset\CV_2\subset\dots
\end{equation}
where 
\begin{equation}
    \CV_k=\mr{Span}\left\{u_{(-n_1)}^{1} u_{(-n_2)}^{2} \cdots u_{(-n_m)}^{m}|0\rangle: n_1 \geq n_2 \geq \ldots \geq n_m, m \leq k;u^i\in U\right\}/\{\mr{null~~ vectors}\}~.
\end{equation}
Notice that if $\CV$ has a one-dimensional vacuum space generated by the vacuum vector, then $\CV_0=\C|0\rangle$. 
Using this filtration, we can define the graded space 
\begin{equation}
    \mr{Gr}(\CV):=\CV^{(0)}\oplus\bigoplus_{n\geq 1}\CV^{(n)}~,
\end{equation}
where 
\begin{equation}
\CV^{(0)}:=\CV_0=\C|0\rangle~,\quad \CV^{(n)}:=\CV_n/\CV_{n-1}~.   
\end{equation}
We can further decompose $\mr{Gr}(\CV)$ into $L_0$-homogeneous subspaces:
\begin{equation}
\mr{Gr}(\CV)=\bigoplus_{n\geq 0}\bigoplus_{h\in\Z}\CV^{(n)}_h~,    
\end{equation}
where $\CV^{(n)}_h$ is the subspace of $\CV^{(n)}$ with $L_0$-eigenvalue $h$. 
Let $\mf{g}$ be the algebra spanned by the zeromodes of the generators of $\CV$. Then the homogeneous spaces $\CV^{(n)}_h$ are $\mf{g}$-modules. 

The refined character of $\CV$ is then defined by\footnote{Note the difference in the exponent of $q$ compared to \cite[Eq. (1.8)]{Song:2016yfd}. The difference arises because of our conventions and can be related to \cite{Song:2016yfd} using $T=t/q$. Our conventions are summarized in Appendix \ref{app:conventions}.}
\begin{equation}\label{eq:ref_char_song}
Z_{\mathcal{V}}^{\mathsf{ref}}(\boldsymbol{z} ; q, t)=\sum_{n=0}^{\infty} \sum_{h\in\Z} \mr{ch}\left(\CV_h^{(n)} ; \boldsymbol{z}\right) q^{h-n} t^n~,
\end{equation}
where  
\begin{equation}
\mr{ch}\left(\CV_h^{(n)} ; \boldsymbol{z}\right)=\mr{Tr}_{\CV_h^{(i)}}\left(\prod_a \boldsymbol{z}_a^{\mr{F}_a}\right)    
\end{equation}
denotes the character of the (non-affine) symmetry $\mathfrak{g}$ on $\CV_h^{(n)}$. Here $\mr{F}_a$ are the Cartan generators of $\mathfrak{g}$. 
\begin{conj}\cite[Song]{Song:2016yfd}
The MacDonald index of a 4d $\CN=2$ \tsf{SCFT} with associated \tsf{VOA} $\CV$ is the refined index of $\CV$:
\begin{equation}\label{eq:song_conj}
    \mathcal{I}_M(\boldsymbol{z};q,t)=Z_{\mathcal{V}}^{\mathsf{ref}}(\boldsymbol{z} ; q, t)~.
\end{equation}
\end{conj}
We remark that in defining the refined character, we assigned the bigrading 
\begin{equation}
    \mr{wt}\left(u_{(-n_1)}^{1} u_{(-n_2)}^{2} \cdots u_{(-n_m)}^{m}|0\rangle\right)=(h-m,m)~,
\end{equation}
where $h$ is the conformal dimension of the state. In other words, the second grading was assigned as $w(u^i)=1$ to each of the strong generators. This is strictly not correct for all strong generators \cite{Song:2016yfd}. For generators of an affine vertex operator algebra and the stress tensor, it is 1 but in general it is not inherently defined in the 2d chiral algebra. If we could track the 4d origin of the strong generators, then there is a consistent way of assigning the weight. Finding a consistent 2d assignment of the second weight is called the $\mr{R}$-filtration problem. Up to these limitations, Song's conjecture gives us a prescription to compute the MacDonald index from the 2d chiral algebra.  
\subsection{The Arc Space Of Zhu Algebra And Its Hilbert Series}
In this section, we recall the notion of Zhu algebra, Hilbert series of arc space of a variety and prove that MacDonald index of 4d $\CN=2$ \tsf{SCFT} is given by the Hilbert series of the arc of space of the Zhu algebra of the corresponding \tsf{VOA}.
\subsubsection{Zhu Algebra Of A (Finitely Strongly Generated) \tsf{VOA}}
Let us begin by recalling the notion of the Zhu algebra of a \tsf{VOA}. 

Let $\CV$ be a \tsf{VOA}. Define the subspace 
\begin{equation}
    C_2(\CV):=\{a_{(-2)}\cdot b:a,b\in\CV\}~.
\end{equation}
Then the Zhu algebra of $\CV$ is defined by 
\begin{equation}
    A(\CV):=\CV/C_2(\CV)~.
\end{equation}
It is well known that $A(\CV)$ is a commutative associative algebra with the multiplication induced by the normal ordered product of vertex operators \cite{Zhu1995ModularIO}. The \tsf{VOA} $\CV$ is called $C_2$-\textit{cofinite} (or \textit{lisse}) if $A(\CV)$ is finite dimensional. Note that $A(\CV)$ inherits a $\Z$-grading from that of $\CV$. 

It is clear that $\CV$ is finitely strongly generated if and only if $A(\CV)$ is finitely generated. In particular, if $\CV$ is generated by $U$ with basis 
\begin{equation}
    \mathcal{B}_U=\{u^1,\dots,u^N\}~,
\end{equation}
then 
\begin{equation}
    A(\CV)\cong\C[x^{(1)},\dots,x^{(N)}]/I~,
\end{equation}
for some ideal $I$. We want to determine the ideal $I$. Roughly speaking, the ideal corresponds to all the singular vector relations among the vertex operators of $u^1,\dots,u^N$. It was conjectured in  in \cite{Xie:2019zlb} that when $\CV$ is a $\mathcal{W}$-algebra then $I$ is generated by these null state relations. We now show that $I$ is generated by the singular vector relations assuming that the conformal vector and vacuum are not singular vectors. This proves a conjecture in \cite{Xie:2019zlb}. Let $\CN$ be the ideal of 
\begin{equation}
\ov{\CV}:= \mathsf{Span}\left\{u^{i_1}_{(-n_1)}\cdots u^{i_k}_{(-n_k)}|0\rangle:n_i\geq 1,k\geq 0\right\}~.   
\end{equation}
generated by the singular vectors, so that $\CV=\ov{\CV}/\CN$. 
By \cite[Proposition 1.4.2]{frenkel1992vertex}, we have 
\begin{equation}
A(\CV)=A(\ov{\CV})/A(\CN)~,    
\end{equation}
where $A(\CN)$ is the image of $\CN$ in $A(\ov{\CV})$ under the projection map. It is clear that the map
\begin{equation}
\begin{split}
    A(\ov{\CV})&\longrightarrow\C[x^{(1)},\dots,x^{(N)}] 
    \\
    [u^{i_1}_{(-1)}\cdots u^{i_k}_{(-1)}|0\rangle]&\longmapsto x^{(i_1)}\cdots x^{(i_k)},~~~~~k\geq 0~,
\end{split} 
\end{equation}
is an isomorphism of algebras. 
We claim that $A(\CN)$ is generated by singular vector relations. Given a singular vector, it corresponds to a polynomial 
\begin{equation}
    F(Y(u^1,z),\partial Y(u^1,z),\dots;Y(u^2,z),\partial Y(u^2,z),\dots; Y(u^N,z),\partial Y(u^N,z),\dots)~,
\end{equation}
where $\dots$ mean higher order derivatives of the vertex operators, 
and the product being normal ordered product of vertex operators. In the Zhu algebra, these polynomials are equivalent to  
\begin{equation}
    F'(Y(u^1,z),Y(u^2,z),\dots,Y(u^N,z))~,
\end{equation}
obtained from $F$ by setting all the monomials containing atleast one derivative to zero. Since $\CN$ is generated by singular vectors,  $A(\CN)$ is generated by the corresponding polynomials.
\subsubsection{Arc Spaces And Hilbert Series}
Let us now define the arc space of an algebraic variety
\begin{equation}\label{eq:variety_R}
    R=\C[x^{(1)},\dots,x^{(N)}]/\langle f_1,\dots,f_K\rangle~.
\end{equation}
We will follow \cite{bruschek2011arc} for our exposition. 
Suppose $R$ is graded with the grading defined by declaring 
\begin{equation}
    \mr{wt}(x^{(i)})=d_i\geq 1~.
\end{equation}
Without loss of generality, we can assume that the polynomials $f_1,\dots,f_K$ are homogeneous with respect to the above grading. We begin by writing the formal power series
\begin{equation}
    x^{(i)}(s)=\sum_{n=d_i-1}^\infty x^{(i)}_ns^n\in\left(\C[[s]]\right)^N~,\quad i=1,\dots,N
\end{equation}
% where 
% \begin{equation}
% \vec{x}_n=\left(x^{(1)}_n,\dots,x^{(N)}_n\right)~.    
% \end{equation}
Let us write the variables as a vector 
\begin{equation}
\vec{x}(s)=\left(x^{(1)}(s),\dots,x^{(N)}(s)\right)~.    
\end{equation}
For $n\in\mathds{Z}_{\geq 0}$, the $n$-th Jet space of $R$ is defined by
\begin{equation}
    R_n:=\C[x^{(1)}_{d_1-1},\dots,x^{(1)}_n,x^{(2)}_{d_2-1},\dots,x^{(2)}_n,\dots,x^{(N)}_{d_N-1},\dots,x^{(N)}_n]/I_n~,
\end{equation}
where 
\begin{equation}
    I_n:=\left\langle f_1(\vec{x}(s)),\dots,f_K(\vec{x}(s))\equiv 0\bmod s^{n+1}\right\rangle~.
\end{equation}
The $n$-th Jet scheme is defined as 
\begin{equation}
    X_n(R):=\mr{Spec}\,R_n~.
\end{equation}
The \textit{arc space} $X_\infty(R)$ of $R$ is defined to be 
\begin{equation}
X_\infty(R)=\lim_{\leftarrow}X_n(R)\cong\mr{Spec}\,R_\infty~,     
\end{equation}
where 
\begin{equation}
 R_\infty:=\C[x^{(1)}_{i_1},\dots,x^{(N)}_{i_N},i_k\geq d_{k}-1]/\left\langle f_1(\vec{x}(s)),\dots,f_K(\vec{x}(s))\right\rangle~.   
\end{equation}
$R_\infty$ admits a filtration
\begin{equation}
R_0\subset R_1\subset R_2\subset\dots    
\end{equation}
Assign bigrading to $R_\infty$ by declaring 
\begin{equation}\label{eq:wt_arc_space_var}
    \mr{wt}^{(2)}(x_j^{(i)}):=\left(\mr{wt}_1(x_j^{(i)}),\mr{wt}_2(x_j^{(i)})\right):=(j,1)~.
\end{equation}
The Hilbert series of $R_\infty$ is then defined by 
\begin{equation}\label{eq:Hilber_ser_def}
H_{R_\infty}(q,t)=\sum_{\ell,m\geq 0} \mathsf{dim}(R_\infty)_{\ell,m}q^\ell t^m~,    
\end{equation}
where $(R_\infty)_{\ell,m}\subset R_\infty$ is the homogeneous subspace with weight $(\ell,m)$.

When $R=A(\CV)$ is the Zhu algebra of a \tsf{VOA}, then the grading is simply the conformal dimension of the operators. 
\subsubsection{Proof Of Theorem \ref{thm:Hil_ser_mac}}
We now show that the Hilbert series of $(A(\CV))_\infty$ is the unflavored MacDonald index of the 4d $\CN=2$ \tsf{SCFT} with chiral algebra $\CV$ assuming that $\CV$ is strongly finitely generated. Let us begin by noting that 
\begin{equation}
    \mr{ch}(\CV^{(n)}_h;1)=\mr{dim}~\CV^{(n)}_h~.
\end{equation}
We now show that 
\begin{equation}\label{eq:hom_space_voa_arc_iso}
\CV^{(n)}_h\cong  \left((A(\CV))_\infty\right)_{h-n,n}~.   
\end{equation}
Let $u^1,\dots,u^N$ be a set of strong generators of $\CV$ with conformal dimension $d_1,\dots,d_N$ respectively, and $x^{(1)},\dots,x^{(N)}$ be the corresponding generators in the Zhu algebra. 

By definition, $\CV^{(n)}$ is given by 
\begin{equation}
    \CV^{(n)}_h=\mr{Span}\left\{u_{(-k_1)}^{i_1} u_{(-k_2)}^{i_2} \cdots u_{(-k_n)}^{i_n}|0\rangle: k_1 \geq k_2 \geq \ldots \geq k_n\geq 1;\sum_{i=1}^n(d_i+k_i-1)=h \right\}/\{\mr{null~~ vectors}\}~.
\end{equation}
Define the map
\begin{equation}\label{eq:voa_arc_space_zhu}
\begin{split}
    \mr{Gr}(\CV)&\longrightarrow (A(\CV))_\infty
    \\
    [u_{(-k_1)}^{i_1} u_{(-k_2)}^{i_2} \cdots u_{(-k_n)}^{i_n}|0\rangle]&\longmapsto [x^{(i_1)}_{k_1+d_{i_1}-2}\cdots x^{(i_n)}_{k_n+d_{i_n}-2}]~,
\end{split}    
\end{equation}
where $[\cdot]$ indicated the equivalence class.  Note that this restricts to a map 
\begin{equation}\label{eq:iso_ref_char_Hilbert_ser}
\CV^{(n)}_h\longrightarrow  \left((A(\CV))_\infty\right)_{h-n,n}~.    
\end{equation}
We first show that the map is well defined. Let us write
\begin{equation}
    A(\CV)=\C[x^{(1)},\dots,x^{(N)}]/\langle f_1,\dots,f_K\rangle~,
\end{equation}
where $f_1,\dots,f_K$ are homogeneous polynomials with respect to the induced grading on $A(\CV)$, corresponding to $K$ singular vectors in 
\begin{equation}
    \mathsf{Span}\left\{u^{i_1}_{(-n_1)}\cdots u^{i_k}_{(-n_k)}|0\rangle:n_i\geq 1,k\geq 1\right\}~.
\end{equation}
These polynomials correspond to polynomials 
\begin{equation}
    F_i(Y(u^1,z),\partial Y(u^1,z),\dots;Y(u^2,z),\partial Y(u^2,z),\dots; Y(u^N,z),\partial Y(u^N,z))~,\quad i=1,\dots,K~.
\end{equation}
with the product being normal ordered product of vertex operators obtained from singular vectors. In the Zhu algebra, these polynomials are equivalent to  
\begin{equation}\label{eq:Fi'_sing_rel}
    F'_i(Y(u^1,z),Y(u^2,z),\dots,Y(u^N,z))~,\quad i=1,\dots,K~,
\end{equation}
obtained from $F_i$ by setting all the monomials containing atleast one derivative to zero. In particular, note that $F_i'=f_i$. Now, noting that only the  modes $u^i_{(-n)}$ with $n\geq 1$ are relevant for $\CV$, we can make an identification 
\begin{equation}
    \sum_{n\geq 1}u^i_{(-n)}z^{n-1}\longleftrightarrow x^{(i)}(s)~,
\end{equation}
to conclude that the ideal generated by singular vectors on the \tsf{LHS} is equivalent to the ideal 
\begin{equation}
    \left\langle f_1(\vec{x}(s)),\dots,f_K(\vec{x}(s))\right\rangle~.
\end{equation}
Indeed, there are two ways of obtaining null vectors in $\CV^{(n)}_h$: 
\begin{enumerate}
    \item From a singular vector at level $h$ corresponding to a polynomial $F_i'$ of homogeneous degree $n$. The the coefficient of the lowest power of $s$ in the expansion of $f_i(\vec{x}(s))$ implements this constraint in $\left((A(\CV))_\infty\right)_{h-n,n}$.
    \item Taking a normal ordered product of other operators with $F_j'$ of degree lower than $n$  which corresponds to a singular vector at level less than $h$ and acting it on the vacuum. On the arc space side, this corresponds to forming polynomials of degree $n$ by multiplying monomials to coefficients of $f_j$ in the expansion in $s$ to form a degree $n$ polynomial.   
\end{enumerate}
These two are precisely the ways constraints appear in the arc space side. This establishes the isomorphism \eqref{eq:iso_ref_char_Hilbert_ser}. In particular, by definition of the Hilbert series \eqref{eq:Hilber_ser_def} and \eqref{eq:song_conj} we see that 
\begin{equation}
\mathcal{I}_M(1;q,t)= H_{A(\CV)_\infty}(q,t)~.   
\end{equation}
\subsection{Computation Of The Hilbert Series}\label{sec:comp_Hilb_ser}
In this section, we explain an algorithmic procedure for the computation of the Hilbert series of the arc space of a variety. The algorithm uses Gr\"{o}bner basis for a finitely generated ideal, relevant definitions and results have been reviewed in Appendix \ref{app:Grob}. 

Assume the notation of \eqref{eq:variety_R} -- \eqref{eq:Hilber_ser_def}. The Hilbert series of $R_n$ is defined as
\begin{equation}
H_{R_n}(q, t)=\sum_{\ell=0}^n \sum_{m \geq 0} \mr{dim}\left(R_n\right)_{\ell, m} q^{\ell} t^m~.
\end{equation}
Taking $n \rightarrow \infty$, we get the Hilbert series of  $R_{\infty}$:
\begin{equation}
H_{R_{\infty}}(q, t)=\sum_{\ell, m \geq 0} \mr{dim}\left(R_{\infty}\right)_{\ell, m} q^{\ell} t^m~.
\end{equation}
One way of computing the Hilbert series of an $n$-the jet space $R_n$ is by using Gr\"{o}bner basis of $I_n$.
Let us denote the Gr\"{o}bner basis of $I_n$ by $G_n$. Note that $G_n$ depends on the choice of ordering. Let us denote the set of leading monomials of the polynomials in $G_n$ by
\begin{equation}
\mr{LM}(G_n):=\{\mr{LM}(g)\mid g\in G_n\}~.
\end{equation}
We call the monomials in $\mr{LM}(G_n)$ as \textit{constraints}. 
Then, we can find a canonical basis of $R_n$ as explained in Appendix \ref{app:Grob}. Let us denote a monomial in $\C[x^{(1)}_{d_1-1},\dots,x^{(1)}_n,x^{(2)}_{d_2-1},\dots,x^{(2)}_n,\dots,x^{(N)}_{d_N-1},\dots,x^{(N)}_n]$ by
\begin{equation}
    X_1^{\vec{a}_1}X_2^{\vec{a}_2}\cdots X_N^{\vec{a}_N}~,\quad \vec{a}_k\in (\IZ_{\geq 0})^n
\end{equation}
where
\begin{equation}
    X_k^{\vec{a}_k}:=(x_{d_k-1}^{(k)})^{a_k^1}(x_{d_1}^{(k)})^{a_k^2}\cdots (x_n^{(k)})^{a_k^n}~.
\end{equation}
We call such monomials as \textit{words}. 
Then a canonical basis of $R_n/I_n$ is given by 
\begin{equation}\label{eq:basis_Rn/I_n}
    \CB_{R_n}:=\{\text{monomial}~~X_1^{\vec{a}_1}X_2^{\vec{a}_2}\cdots X_N^{\vec{a}_N}:g\nmid X_1^{\vec{a}_1}X_2^{\vec{a}_2}\cdots X_N^{\vec{a}_N}\text{ for all }g\in \mr{LM}(G_n)\}~,
\end{equation}
i.e., words modulo constraints. 
Let $(\CB_{R_n})_{\ell,m}\subseteq \CB_{R_n}$ be the subset of monomials of weight $(\ell,m)$. Then the Hilbert series for $R_n$ can be written as 
\begin{equation}
H_{R_n}(q, t)=\sum_{\ell=0}^n \sum_{m \geq 0} |(\CB_{R_n})_{\ell,m}|q^{\ell} t^m~.
\end{equation}
Taking $n \rightarrow \infty$, we get the Hilbert series of  $R_{\infty}$:
\begin{equation}
H_{R_{\infty}}(q, t)=\sum_{\ell, m \geq 0} |(\CB_{R_n})_{\ell,m}| q^{\ell} t^m~.
\end{equation}
In Appendix \ref{app:Hilb_ser_count}, we prove the following formula for the Hilbert series of $R_n$: 
\begin{equation}
\label{eq: hilbert_ser_Rn_gen}
H_{R_n}(q,t) = \mr{Ser}_{q=0,t=0}\frac{1+\sum_{S\subseteq \mr{LM}(G_n)}(-1)^{|S|}(q,t)^{\mr{wt}^{(2)}(\mr{LCM}(S))}}{\prod_{i=1}^N\prod_{j\geq 0}\left(1-q^{(d_i-1)+j}t\right)} \quad (\mr{mod} \;q^{n+1})~,
\end{equation}
where $\mr{Ser}_{q=0,t=0}(f(q,t))$ denotes the power series of $f(q,t)$ around $q=t=0$, the summation in the numerator is over all subsets of $\mr{LM}(G_n)$, $\mr{LCM}(S)$ is least common multiple of the monomials in $S$, and we have defined
\begin{equation}
(q,t)^{(\ell,m)} := q^{\ell} t^m ~.   
\end{equation}
The full Hilbert series is then given by taking $n\to \infty$:
\begin{equation}
\label{eq: hilbert_ser_Rinfty_gen}
H_{R_\infty}(q,t) = \mr{Ser}_{q=0,t=0}\frac{1+\sum_{S\subseteq \mr{LM}(G_\infty)}(-1)^{|S|}(q,t)^{\mr{wt}^{(2)}(\mr{LCM}(S))}}{\prod_{i=1}^N\prod_{j\geq 0}\left(1-q^{(d_i-1)+j}t\right)}~,
\end{equation}
where 
\begin{equation}
    \mr{LM}(G_\infty):=\lim _{\leftarrow } \mr{LM}(G_n)~.
\end{equation}
This formula will be primarily used in Section \ref{sec:Hilb_arc_MacD3} for computation of MacDonald index of $(A_1,D_3)$ theory. 
\subsubsection{Example: $(A_1,A_{2n})$ Argyres-Douglas Theory}
In this section, we work out the MacDonald index of $(A_1,A_{2n})$ theory using the Hilbert series of arc space of the Zhu algebra of the corresponding \tsf{VOA} as a demonstration of Theorem \ref{thm:Hil_ser_mac}. \par 
The MacDonald index for $(A_1,A_{2n})$ theory is given by \cite{Buican:2015tda}
\begin{equation}\label{eq:Mac_A_1A2n}
    \CI_{\mr{M}}^{(A_1,A_{2n})}(q,t)=\sum_{N_1\geq N_2\geq\dots\geq N_{n}\geq 0}t^{N_1+N_2+\dots+N_{n}}\frac{q^{N_1^2+\dots+N_{n}^2}}{(q;q)_{N_1-N_2}\dots (q;q)_{N_{n-1}-N_{n}}(q;q)_{N_{n}}}~.
\end{equation}
The associated \tsf{VOA} of the $(A_1,A_{2n})$ Argyres-Douglas theories is \cite{beem2014infinite,Beem:2017ooy}
\begin{equation}
    \CV_{(A_1,A_{2n})}=\mr{Vir}\left(c=-\frac{2 n(5+6 n)}{3+2 n},0\right)~,
\end{equation}
where $\mr{Vir}(c,0)$ denotes the Virasoro \tsf{VOA} with central charge $c$. The Zhu algebra of $\CV_{(A_1,A_{2n})}$ is given by \cite{Beem:2017ooy} 
\begin{equation}\label{eq:zhuA1A2}
A(\CV_{(A_1,A_{2n})})\cong \C[J]/\langle J^{n+1}\rangle~.    
\end{equation}
Here $J$ is the variable corresponding to the stress tensor in the \tsf{VOA}.
The MacDonald index for $(A_1,A_2)$ theory was reproduced from the Hilbert series for the arc space of $A(\CV_{(A_1,A_{2})})$ in \cite{Bhargava:2023hsc}, the physical reasoning was an exact map of operators of the theory to the free vector multiplet. Here we reproduce some of the calculations from \cite{Bhargava:2023hsc} but using the algorithm given in Section \ref{sec:comp_Hilb_ser} and show that it reproduces the MacDonald index in line with our general Theorem \ref{thm:Hil_ser_mac}. For $(A_1,A_4)$ theory, we compute the Hilbert series to $O(t^6)$ and show that it matches with the MacDonald index. For $n>1$, we have done numerical checks and confirmed that it reproduces the MacDonald index.  
\\\\
$\underline{\boldsymbol{n=1}.}$ The Hilbert series for the arc space of \eqref{eq:zhuA1A2} for $n=1$ was computed rigorously in  \cite{bruschek2011arc,bai2020quadratic}. Here we present a computation using our formula \eqref{eq: hilbert_ser_Rinfty_gen}.  We start with the formal series 
\begin{equation}
    J(x)=\sum_{m=1}^\infty J_mx^m~.
\end{equation}
We choose the ordering of variables to be
\begin{equation}
J_1>J_2>\dots
\end{equation}
With the degree lexicographic ordering of monomials, Mathematica calculations give is the leading terms of the Gr\"{o}bner basis to be
\begin{equation}
\mr{LM}\left(G_s^{\left(A_1, A_2\right)}\right)=\bigcup_{m \geq 2} \bigcup_{0 \leq \ell \leq s} I_{\ell, m}^{\left(A_1, A_2\right)}~,
\end{equation}
where $I_{\ell, m}^{\left(A_1, A_2\right)}$ is given by
\begin{equation}
I_{\ell, m}^{\left(A_1, A_2\right)}:=\{J_{m-1} \}\cdot F_{m-1}(\ell-m+1, m-1)~,
\end{equation}
where for sets $A,B$ we have introduced the notation $A\cdot B:=\{ab:a\in A,b\in B\}$, $F_s(\ell, m)$ is the set of all monomials generated by $\left\{J_i: i \geq s\right\}$ of weight $(\ell, m)$.
The set $F_m(\ell, m)$ is thus the set of monomials of weight $(\ell, m)$. We now want to count the number of basis elements of $(A(\CV_{(A_1,A_2)})_n/I_n)_{\ell,m}$ obtained from Gr\"{o}bner basis using the formula \eqref{eq:basis_Rn/I_n}. The set $F(\ell,m)$ of all monomials of weight $(\ell,m)$ are given by 
\begin{equation}
    F(\ell,m):=\left\{J_{k_1}^{a_1}\cdots J_{k_s}^{a_s}:s\geq 0,k_1\leq k_2\leq\dots\leq k_s,\sum_{i=1}^sa_ik_i=\ell,\sum_{i=1}^sa_i=m\right\}~.
\end{equation}
From $F(\ell,m)$, we need to remove the monomials
\begin{equation}
\{J_k\}\cdot F_k(\ell-k,m-1)~,\quad k=1,\dots,m-1~,    
\end{equation}
because these monomials are part of the leading monomials of the Gr\"{o}bner basis. This means that that 
\begin{equation}
|(A(\CV_{(A_1,A_2)})_n/I_n)_{\ell,m}|=|F_m(\ell,m)|~.    
\end{equation}
The Hilbert series is then given by
\begin{equation}
H_{(A(\CV_{(A_1,A_2)})_\infty}(q,t)=\sum_{\ell \geq 0, m \geq 0}\left|F_m(\ell, m)\right| q^\ell t^m~.
\end{equation}
We now claim that 
\begin{equation}
    \sum_{\ell \geq 0, m \geq 0}\left|F_m(\ell, m)\right| q^\ell t^m=\sum_{m \geq 0} \frac{q^{m^2}}{(q)_m} t^m~,
\end{equation}
where 
\begin{equation}
    (q)_m:=(q;q)_m=\prod_{i=1}^m(1-q^i)~.
\end{equation}
To prove this, let us first define 
\begin{equation}
\label{def:Ism}
\mathcal{I}_m^s(q):=\sum_{\ell \geq 0}\left|F_s(\ell, m)\right| q^\ell~.
\end{equation}
Then we claim that
\begin{equation}\label{eq:Ism}
\mathcal{I}_m^s(q)=\frac{q^{sm}}{(q)_m}~.
\end{equation}
The generating function for $|F(\ell,m)|$ is given by 
\begin{equation}
    \sum_{\ell,m\geq 0}|F(\ell,m)|q^\ell t^m=\mr{Ser}_{t=0, q=0} \frac{1}{\prod_{i=1}^{\infty}\left(1-q^it\right)}~.
\end{equation}
It is also easy to see that 
\begin{equation}\label{eq:free_gen_ser_A1A4}
\mr{Ser}_{t=0, q=0} \frac{1}{\prod_{i=s}^{\infty}\left(1-q^it\right)}=\sum_{\ell,m\geq 0}|F_s(\ell,m)|q^\ell t^m=\sum_{m\geq 0}\CI_m^s(q)t^m~.    
\end{equation}
Using the formula 
\begin{equation}
    \frac{1}{(z;q)_\infty}=\sum_{m\geq 0}\frac{z^m}{(q;q)_m}~,
\end{equation}
for $z=tq$, we deduce that 
\begin{equation}\label{eq:I1m}
    \CI^1_{m}(q)=\frac{q^m}{(q)_m}~.
\end{equation}
It is easy to see that 
\begin{equation}\label{eq:Is0}
    \CI^s_0=1~.
\end{equation}
Next note that 
\begin{equation}
 \left(1-q^s\right) \sum_{m \geq 0} \mathcal{I}_m^s(q) t^m=\sum_{m \geq 0} \mathcal{I}_m^{s+1}(q) t^m ~.  
\end{equation}
This gives us the recursion relation
\begin{equation}
\mathcal{I}_m^{s+1}(q)=\mathcal{I}_m^s(q)-q^s \mathcal{I}_{m-1}^s(q)~,\quad m>0~.
\end{equation}
The solution to this recursion relation subject to initial conditions \eqref{eq:I1m} and \eqref{eq:Is0} is given by \eqref{eq:Ism}. Indeed
\begin{equation}
\mathcal{I}_m^{s+1}(q)=q^{s m}\left(\frac{1}{(q)_m}-\frac{1}{(q)_{m-1}}\right)=\frac{q^{(s+1)_m}}{(q)_m}~.
\end{equation}
This proves that 
\begin{equation}\label{eq:Hil_A1A2}
\sum_{\ell \geq 0, m \geq 0}\left|F_m(\ell, m)\right| q^\ell t^m=\sum_{m\geq 0}\CI_m^m(q)t^m=\sum_{m \geq 0} \frac{q^{m^2}}{(q)_m} t^m~.    
\end{equation}
Thus the Hilbert series \eqref{eq:Hil_A1A2} matches with the MacDonald index \eqref{eq:Mac_A_1A2n} of $(A_1,A_2)$ theory. 
\\\\
\underline{$\boldsymbol{n=2}$.} For $(A_1,A_4)$ theory, using numerical checks we obtain the following expression for the Hilbert series of the arc space of $A(\CV_{(A_1,A_4)})$:
\begin{equation}
H_{(A(\CV_{(A_1,A_4)})_\infty}(q,t)=\sum_{m \geq 0}\CI^{(A_1,A_4)}_m(q) t^m~,   
\end{equation}
where 
\begin{equation}
    \mathcal{I}_m^{\left(A_1, A_4\right)}= \begin{cases}\mathcal{I}_m^{m / 2}+\sum_{i=1}^{(m-2) / 2} I_m^{2 i, m / 2-i}, & m \text { is even }, \\ \mathcal{I}_m^{(m+1) / 2}+\sum_{i=1}^{(m-1) / 2} I_m^{2 i-1,(m+1) / 2-i}, & m \text { is odd } ,\end{cases}
\end{equation}
and $I_m^{c,s}(q)$ can be obtained from the recursion relation:
\begin{equation}
I_m^{c, s}= \begin{cases}q^s \mathcal{I}_{m-1}^{s+1}, & c=1, \\ q^s \mathcal{I}_{m-1}^{s+2}, & c=2,3, \\ q^s \mathcal{I}_{m-1}^{s+c+1}+q^s \sum_{i=1}^{c / 2-1} I_{m-1}^{2 i-1, s+c / 2+1-i}, & c>3 \text { and even }, \\ q^s \mathcal{I}_{m-1}^{s+(c+1) / 2}+q^s \sum_{i=1}^{(c-3) / 2} I_{m-1}^{2 i, s+(c+1) / 2-i}, & c>3 \text { and odd~. }\end{cases}    
\end{equation}
See Appendix \ref{app:Hilb_A1A4} for details of the calculation. 
This formula matches the MacDonald index \eqref{eq:Mac_A_1A2n} for the theory $(A_1,A_4)$ up to $O(t^{16})$. More concrelty, we can change variables to $N_1+N_2=m,N_2=k$ and write the MacDonald index as
\begin{equation}
    \CI_{\mr{M}}^{(A_1,A_4)}(q,t)=\sum_{m\geq 0}t^m\left(\sum_{k=0}^{\lfloor\frac{m}{2}\rfloor}\frac{q^{(m-k)^2+k^2}}{(q;q)_{m-2k}(q;q)_k}\right)~.
\end{equation}
Then, our formula for MacDonald index using the Hilbert series predicts the following identity:
\begin{equation}
\mathcal{I}_m^{\left(A_1, A_4\right)}(q)=\sum_{k=0}^{\lfloor\frac{m}{2}\rfloor}\frac{q^{(m-k)^2+k^2}}{(q;q)_{m-2k}(q;q)_k}~.    
\end{equation}
This was checked numerically for large values of $m$. 
We conjecture that this is true for all $m$. 
% \paragraph{$(A_1,D_{2n+1})$ Argyres-Douglas Theory.} 
% \myRed{Seems repetitive. Can be dropped.}\\
% The MacDonald index for $(A_1,D_{2n+1})$ obtained using the class 
% $\CS$ description has been recorded in \eqref{eq:Mac_class_S-1} -- \eqref{eq:Mac_class_S-7}. The associated \tsf{VOA} is given by  
% \begin{equation}
%     \CV_{(A_1,D_{2n+1})}=\mathfrak{su}(2)_{-\frac{4n}{2n+1}}~.
% \end{equation}
% The Zhu algebra is known for $n=1,2$ \cite{}
% \begin{align}
%     A(\CV_{(A_1,D_{3})})\cong \frac{\C[j^1,j^2,j^3]}{\langle j^a((j^1)^2+(j^2)^2+(j^3)^3),((j^1)^2+(j^2)^2+(j^3)^2)^2,a=1,2,3\rangle}~,
%     \\
%     A(\CV_{(A_1,D_{5})})\cong \frac{\C[j^1,j^2,j^3]}{\langle j^a((j^1)^2+(j^2)^2+(j^3)^3),((j^1)^2+(j^2)^2+(j^3)^2)^3,a=1,2,3\rangle}~,
% \end{align}
% where the generators $j^a$ correspond to the Kac-Moody currents of affine $\mf{su}(2)$. In Section \ref{sec:A1D3_Mac_conj} and Section \ref{sec:A1D2n+1_Mac_conj}, compute the Hilbert series of the arc space of this Zhu algebra to support our proposal for the MacDonald index of $(A_1,D_3)$ and $(A_1,D_5)$ theory. 
\section{Macdonald Index For $(A_1,D_{3})$ Theory}\label{sec:A1D3_Mac_conj}
In this section, we conjecture a very simple expression for the MacDonald index of $(A_1,D_3)$ and provide powerful checks of our formula. Our starting point is the observation, supported by Higgs branch \tsf{RG} flow explained in Section \ref{sec:RG_flow}, that the MacDonald index for the $(A_1,D_3)$ theory must be related to the MacDonald index for the (projected) free hypermultiplet (see Footnote \ref{foot:proj_hyper} for a remark about the term ``projected''.). 
\subsection{Relation To Free Hypermultiplet Index}\label{sec:D3_free_hyper_rel}
We begin by reviewing the MacDonald index for (projected) free hypermultiplet.
\subsubsection{Free Hypermultiplet}
The 4d $\mathcal{N}=2$ hypermultiplet consists of two 4d $\mathcal{N}=1$ chiral (and anti-chiral) multiplets $Q$ $(q,q^\dagger,\psi,\bar\psi)$ and 
$\widetilde{Q}$ $(\tilde{q},\tilde{q}^\dagger, \chi,\bar\chi)$.
$(q,\tilde{q}^\dagger)$ and 
$(- \tilde{q},q^\dagger)$
transform as doublets of $\mr{SU(2)_R}$. We choose the $\mr{U(1)}_r$ charge of $q$ and 
$\widetilde{q}$ to be zero. 
There is a global $\mr{U(1)_B}$ baryonic symmetry that acts differently on each half-hypermultiplet, 
\begin{equation}
 Q \to e^{i\varphi } Q,\quad \widetilde{Q}\to e^{-i\varphi} \widetilde Q~.   
\end{equation}
The letters contributing the index of the free 
$\mathcal{N}=2$ (half) hypermultiplet  are shown in Table \ref{tab:half_hyper_charge}.
\begin{table}[H]
    \centering
    \begin{tabular}{|c|c|c|c|c|c|c|}
    \hline
        \text{Letters} & 
        $E$ &
        $j_1$& $j_2$ & $R$ & $r$ &
        $\mathcal{I}(p,q,t)$\\
        \hline
         $q$& $1$ & $0$ & $0$ & $\frac{1}{2}$ &  $0$ & $t^{1/2}$\\
         $\bar\psi_{\dot +}$& 
         $\frac{3}{2}$
         & $0$ & $\frac{1}{2}$ & $0$ &$-\frac{1}{2}$  &
         $-pqt^{-1/2} $
         \\
         $\partial_{\pm \dplus}$
         & $1$ & $\pm \frac{1}{2}$ &$\frac{1}{2}$  & $0$ & $0$ & $p,q$\\
         \hline
    \end{tabular}
    \caption{Charge assignments of various operators in free (half) hypermultiplet.}
    \label{tab:half_hyper_charge}
\end{table}
\noindent We have the following single particle indices for free hypermultiplet
\begin{equation}
    \mathcal{I}^{\mr{HM}}_{s . p .}=t^{\frac{1}{2}} \frac{1-\frac{p q}{t}}{(1-p)(1-q)}\left(a+a^{-1}\right) \chi_{\Lambda}(\boldsymbol{x})~,
\end{equation}
where $a$ is the $U(1)_B$ fugacity and 
$\chi_\Lambda(\boldsymbol{x})$ is the character of the representation of some global symmetry.
In this paper, we will ignore flavor fugacities and consequently take $\boldsymbol{x} = 1,a=1$. We will take $a\to 1$ at the end. 
We get the superconformal index of a free hypermultiplet 
\begin{equation}
\begin{aligned}
    \mathcal{I}_{\mr{SC}}^{\mr{HM}}(p,q,t;a ) &= 
    \mathsf{PE}\left[
    t^{\frac{1}{2}} \frac{1-\frac{p q}{t}}{(1-p)(1-q)}\left(a+a^{-1}\right)
    \right]\\
    & = 
    \Gamma(t^{\frac{1}{2}} a; p,q)
    \Gamma(t^{\frac{1}{2}} a^{-1}; p,q)~,
\end{aligned}
\end{equation}
where the plethystic exponential is given by 
\begin{equation}
    \mr{PE}[f(x, y, \cdots)]:=\exp \left[\sum_{\ell=1}^{\infty} \frac{1}{\ell} f\left(x^{\ell}, y^{\ell}, \cdots\right)\right]~,
\end{equation}
and the elliptic gamma function is defined by 
\begin{equation}
    \Gamma(z ; p, q):=\mathrm{PE}\left[\frac{z-\frac{p q}{z}}{(1-p)(1-q)}\right]=\prod_{i, j=0}^{\infty} \frac{1-p^{i+1} q^{j+1} z^{-1}}{1-p^i q^j z}~.
\end{equation}
The Macdonald index of a free hypermultiplet is obtained by taking $p\to 0$ in the superconformal index: 
\begin{equation}
    \mathcal{I}_{\mr{SC}}^{\mr{HM}}( p\to 0, q,t; a) = 
    \prod_{n=0}^\infty 
    \frac{1}{\left(1-q^n t^{\frac{1}{2}} a\right)}
    \frac{1}{\left(1-q^n t^{\frac{1}{2}} a^{-1}\right)}.
\end{equation}
Taking $a \to 1$, we can draw a coefficient diagram for Macdonald index of a free hypermultiplet as in Table \ref{tab:coeff_mac_FH}.
\begin{table}[h]
    \centering
    \begin{tabular}{|c|ccccccccccc|}
    \hline
   Variables & $t^0$ & $t^{1/2}$ 
    & $t^{1}$ & $t^{3/2}$
    & $t^{2}$ & $t^{5/2}$
    & $t^{3}$ & $t^{7/2}$
    & $t^{4}$ & $t^{9/2}$
    & $t^{5}$ \\
    \hline
 $q^{0}$ &1 & 2 & 3 & 4 & 5 & 6 & 7 & 8 & 9 & 10 & 11 \\
 $q^{1}$ &0 & 2 & 4 & 6 & 8 & 10 & 12 & 14 & 16 & 18 & 20 \\
 $q^{2}$ &0 & 2 & 7 & 12 & 17 & 22 & 27 & 32 & 37 & 42 & 47 \\
 $q^{3}$ &0 & 2 & 8 & 18 & 28 & 38 & 48 & 58 & 68 & 78 & 88 \\
 $q^{4}$ &0 & 2 & 11 & 26 & 46 & 66 & 86 & 106 & 126 & 146 & 166 \\
 $q^{5}$ &0 & 2 & 12 & 34 & 64 & 100 & 136 & 172 & 208 & 244 & 280 \\
 $q^{6}$ &0 & 2 & 15 & 46 & 94 & 152 & 217 & 282 & 347 & 412 & 477 \\
 $q^{7}$ &0 & 2 & 16 & 56 & 124 & 214 & 316 & 426 & 536 & 646 & 756 \\
 $q^{8}$ &0 & 2 & 19 & 70 & 167 & 302 & 464 & 640 & 825 & 1010 & 1195 \\
 $q^{9}$ &0 & 2 & 20 & 84 & 212 & 406 & 648 & 922 & 1212 & 1512 & 1812 \\
 $q^{10}$ &0 & 2 & 23 & 100 & 271 & 542 & 899 & 1314 & 1766 & 2236 & 2717 \\\hline
\end{tabular}
\caption{Table of first few coefficients of the Macdonald index for a free hypermultiplet.}
\label{tab:coeff_mac_FH}
\end{table}
The unflavoured index can also be written as an infinite sum
\begin{equation}
    \CI_{\mr{M}}^{\mr{HM}}(q,t)=\sum_{2 n=0}^{\infty} t^n \sum_{k=0}^{2 n} \frac{1}{(q ; q)_k(q ; q)_{2 n-k}}~.
\end{equation}
The projected free hypermultiplet is obtained from the free hypermultiplet by projecting out the operators with half-integral $\mr{U(1)}_r$-charge\footnote{\label{foot:proj_hyper}More precisely, the free hypermultiplet has a global $\mathds{Z}_2\subset \mr{U(1)}_r$ symmetry. The projected free hypermultiplet that we will use throughout this work is obtained by projecting onto the local operators invariant under this symmetry. However, we note that we only work with the Schur ring of these theories, so for our purposes, this is not a projection (or $\mathds{Z}_2$-gauging) at the level of the full theory, but merely a convenient way to set up the operator map between the Schur rings of the \tsf{UV} and \tsf{IR} theories. We do not set up enough structure to treat the projected free hypermultiplet as a gauged version of the free hypermultiplet.}.
Consequently, we need to remove the columns with $t^{n+\frac{1}{2}},~n\in\Z_{\geq 0}$ from Table \ref{tab:coeff_mac_FH} to obtain the coefficient list of projected free hypermultiplet, we reproduce it here for later reference in Table \ref{tab:coeff_mac_proj_FH}.
\begin{table}[h]
    \centering
    \begin{tabular}{|c|cccccc|}
    \hline
   Variables & $t^0$ & $t^{1}$ & $t^{2}$ & $t^{3}$ & $t^{4}$ & $t^{5}$ \\
    \hline
    $q^{0}$  & 1 & 3 & 5 & 7 & 9 & 11 \\
    $q^{1}$  & 0 & 4 & 8 & 12 & 16 & 20 \\
    $q^{2}$  & 0 & 7 & 17 & 27 & 37 & 47 \\
    $q^{3}$  & 0 & 8 & 28 & 48 & 68 & 88 \\
    $q^{4}$  & 0 & 11 & 46 & 86 & 126 & 166 \\
    $q^{5}$  & 0 & 12 & 64 & 136 & 208 & 280 \\
    $q^{6}$  & 0 & 15 & 94 & 217 & 347 & 477 \\
    $q^{7}$  & 0 & 16 & 124 & 316 & 536 & 756 \\
    $q^{8}$  & 0 & 19 & 167 & 464 & 825 & 1195 \\
    $q^{9}$  & 0 & 20 & 212 & 648 & 1212 & 1812 \\
    $q^{10}$ & 0 & 23 & 271 & 899 & 1766 & 2717 \\
    \hline
    \end{tabular}
    \caption{Table of selected coefficients of the Macdonald index for a projected free hypermultiplet.}
    \label{tab:coeff_mac_proj_FH}
\end{table}

The resulting MacDonald index can be written as
\begin{equation}
    \CI_{\mr{M}}^{\mr{HM}/\mathds{Z}_2}(q,t)=\sum_{\ell=0}^{\infty}\frac{t^\ell}{(q;q)_{2\ell}}\sum_{k=0}^{2\ell}{2\ell\choose k}_q~,
\end{equation}
where 
\begin{equation}
\binom{2m}{k}_q=\frac{(q;q)_{2m}}{(q;q)_{2m-k}(q;q)_k}~,    
\end{equation}
is the $q$-binomial coefficient.
%\subsubsection{MacDonald Index Of $(A_1,D_3)$ Theory}

%and the index is obtained as a partition function of a 2d \tsf{TQFT} (valued in 4d $\mathcal{N}=2$ theories of Class $\CS$; \textcolor{red}{Cite GMN and Arakawa's paper on Moore-Tachikawa varieties}) 
% The formula for $(A_1,D_{2N-3})$ is recorded in \eqref{eq:Mac_class_S-1} -- \eqref{eq:Mac_class_S-7}, which for the case of $(A_1,D_{3})$ theory is given by
% \begin{equation}
%     \CI_{\mr{M}}^{(A_1,D_3)}(q, t, x)=\sum_{\lambda=0}^{\infty} C_\lambda \tilde{f}_\lambda^{(n)}(q, t, a)~,
% \end{equation}
% where $a$ is the fugacity for the $\mf{u}(1)$ Cartan subalgebra of the $\mf{su}(2)$ flavor symmetry and 
% \begin{align}
% C_\lambda&=\left[\frac{(t ; q)}{(q ; q)}\right]^{1 / 2} \frac{P_\lambda(q, t, \sqrt{t})}{\left(t^2 ; q\right)}~, \\
% \tilde{f}_\lambda^{(n)}(q, t, a)&=\left(\frac{(t ; q)}{(q ; q)}\right)^{1 / 2} \operatorname{PE}\left[\frac{t}{1-q}\right] \tilde{P}_\lambda(q, t, a)~, \\
% P_\lambda(q, t, a)&=N_\lambda(q, t) \sum_{m=0}^\lambda \frac{(t ; q)_m(t ; q)_{\lambda-m}}{(q ; q)_m(q ; q)_{\lambda-m}} a^{2 m-\lambda}~, \\
% \tilde{P}_\lambda(q, t, x)&=N_\lambda(q, t) t^{n \lambda / 2} q^{n \lambda^2 / 4} \sum_{m=0}^\lambda \frac{(t ; q)_m(t ; q)_{\lambda-m}}{(q ; q)_m(q ; q)_{\lambda-m}^{2 m-\lambda} q^{-n\left(\frac{\lambda}{2}-m\right)^2}}~, \\
% N_\lambda(q, t)&=\frac{(q ; q)_\lambda}{(t ; q)_{\infty}} \sqrt{\left(1-t q^\lambda\right) \prod_{k=\lambda}^{\infty}\left(1-q^{k+1}\right)\left(1-t^2 q^k\right)}~.
% \end{align}
% We can expand it in powers of $t$ to write 
% \begin{equation}
% \begin{aligned}
% & \CI_{\mr{M}}(q, t, a)=1+t\left(a^2+1+a^{-2}\right)+q t+q t\left(a^2+1+a^{-2}\right)+q^2 t+q^2 t\left(a^2+1+a^{-2}\right) \\
% & +q^3 t+q^3 t\left(a^2+1+a^{-2}\right)+t^2\left(a^4+a^2+1+a^{-2}+a^{-4}\right)+q t^2\left(a^4+2 a^2+2+2 a^{-2}+a^{-4}\right) \\
% & +q^2 t^2\left(2 a^4+4 a^2+5+4 a^{-2}+2 a^{-4}\right)+t^3\left(a^6+a^4+a^2+1+a^{-2}+a^{-4}+a^{-6}\right)+\mathcal{O}\left(t^4\right)~.
% \end{aligned}
% \end{equation}
% We can organize the sum into characters of $\mf{su}(2)$ to write:
% \begin{equation}
% \begin{aligned}
% & \CI_{\mr{M}}(q, t, a)=1+\left(t \chi_1(a)+t^2 \chi_2(a)+t^3 \chi_3(a)\right)+\left(q t+q^2 t+q^3 t\right)\left(1+\chi_1(a)\right) \\
% & +q t^2\left(\chi_1(a)+\chi_2(a)\right)+q^2 t^2\left(1+2 \chi_1(a)+2 \chi_3(a)\right)+\cdots
% \end{aligned}
% \end{equation}
% Taking $a\to 1$, 

\subsubsection{A Closed Formula For The MacDonald Index}
Let us start with the closed form expression for the MacDonald index for $(A_1,D_3)$ theory obtained from Class $\CS$ construction. The formula for $(A_1,D_{2N-3})$ is recorded in \eqref{eq:Mac_class_S-1} -- \eqref{eq:Mac_class_S-7}. Simplifying it for $N=3$, we obtain Table \ref{tab:coeff_mac_A1D3} of coefficients of the MacDonald index for $(A_1,D_3)$ theory. 
\begin{table}[]
    \centering
    \begin{tabular}{|c|ccccccccc|} 
    \hline
Variables & $t^0$ & $t^1$ & $t^2$ & $t^3$ & $t^4$ & $t^5$ & $t^6$ & $t^7$ & $t^8$ \\\hline
$q^0$ & 1 & 3 & 5 & 7 & 9 & 11 & 13 & 15 & 17 \\
$q^1$ & 0 & 4 & 8 & 12 & 16 & 20 & 24 & 28 & 32 \\
$q^2$ & 0 & 4 & 17 & 27 & 37 & 47 & 57 & 67 & 77 \\
$q^3$ & 0 & 4 & 23 & 48 & 68 & 88 & 108 & 128 & 148 \\
$q^4$ & 0 & 4 & 33 & 79 & 126 & 166 & 206 & 246 & 286 \\
$q^5$ & 0 & 4 & 39 & 117 & 199 & 280 & 352 & 424 & 496 \\
$q^6$ & 0 & 4 & 49 & 171 & 322 & 466 & 607 & 737 & 867
\\
\hline
\end{tabular}
    \caption{Table of selected coefficients of the Macdonald index for $(A_1,D_3)$ theory.}
    \label{tab:coeff_mac_A1D3}
\end{table}

A crucial observation is that the coefficients in Table \ref{tab:coeff_mac_A1D3} and Table \ref{tab:coeff_mac_proj_FH} above the diagonal are in perfect agreement. This suggests that each column in the MacDonald index for $(A_1,D_3)$ can be written as a product of the corresponding column in $\CI_{\mr{M}}^{\mr{HM}/\mathds{Z}_2}(q,t)$ and some other factor. Our goal in this section is to figure out the factor by checking the first few columns. 

The contribution of the second column to the MacDonald index is 
\begin{equation}
t\left(3+4 q+4 q^2+4 q^3+\cdots\right)~,    
\end{equation}
which can be written as 
\begin{equation}
\begin{split}
    &t\left(1-q^2\right)\left(3+4 q+7 q^2+8 q^3+11 q^4+12 q^5+\cdots\right)\\&=t\left(1-q^2\right)\left(\frac{1}{(q ; q)_0(q ; q)_2}+\frac{1}{(q ; q)_1(q ; q)_1}+\frac{1}{(q ; q)_2(q ; q)_0}\right)~.
\end{split}
\end{equation}
This has the desired form as it has the contribution of the second column of the projected free hypermultiplet. At order $t^2$, we can the write the contribution as 
\begin{equation}
\begin{split}
    &t^2\left(5+8 q+17 q^2+23 q^3+33 q^4+39 q^5+49 q^6+\cdots\right)\\=~&t^2\left(1-q^3\right)\left(1-q^4\right)\left(\frac{1}{(q ; q)_0(q ; q)_4}+\frac{1}{(q ; q)_1(q ; q)_3}+\frac{1}{(q ; q)_2(q ; q)_2}+\frac{1}{(q ; q)_3(q ; q)_1}+\frac{1}{(q ; q)_4(q ; q)_0}\right)~,
\end{split}
\end{equation}
which again has the desired form. We further notice that 
\begin{equation}
\begin{aligned}
1&=(q ; q)_0 \\
\left(1-q^2\right)&=\left(q^2 ; q\right)_1 \\
\left(1-q^3\right)\left(1-q^4\right)&=\left(q^3 ; q\right)_2~.
\end{aligned}
\end{equation}
This pattern motivates our formula for the MacDonald index:
\begin{equation}
    \CI_{\mr{M}}^{(A_1,D_3)}(q, t)=\sum_{n=0}^{\infty} t^n\left(\sum_{k=0}^{2 n} \frac{1}{(q ; q)_k(q ; q)_{2 n-k}}\right)\left(q^{n+1} ; q\right)_n~.
\end{equation}
Using the formula
\begin{equation}
(a ; q)_{n+k}=(a ; q)_n\left(a q^n ; q\right)_k~,    
\end{equation}
we obtain our final proposal for the MacDonald index:
\begin{conj}
The MacDonald index for the $(A_1,D_3)$ Argyres-Douglas theory is given by 
\begin{equation}\label{eq:Mac_ind_A1D3}
\CI_{\mr{M}}^{(A_1,D_3)}=\sum_{m=0}^\infty 
    \frac{t^m}{(q;q)_m} \sum_{k=0}^{2m}\binom{2m}{k}_q~.    
\end{equation}
\end{conj}
This formula reproduces Table \ref{tab:coeff_mac_A1D3}. 
\subsection{Checks Of Our Formula}
In this section, we show that the formula \eqref{eq:Mac_ind_A1D3} satisfies several expected properties. 
\subsubsection{Matching The Schur Limit}
The first check is the matching of the Schur index obtained by taking the Schur limit. The Schur index of $(A_1,D_{2n+1})$ is given by (for details, we refer the reader to \cite{Cordova:2015nma, Song:2016yfd}, and to \cite{Buican:2015ina, Buican:2015hsa} for the specific case of the $(A_1,D_3)$ theory)
\begin{equation}
    \CI_{\mathsf{S}}^{(A_1,D_{2n+1})}(q,a)=\mr{PE}\left[\frac{q-q^{2 n+1}}{(1-q)\left(1-q^{2 n+1}\right)} \chi_{\mr{adj}(a)}\right]~,
\end{equation}
where 
\begin{equation}
    \chi_{\mathsf{adj}}(a)=a^2+\frac{1}{a^2}+1~, 
\end{equation}
is the character of the adjoint representation of the flavor symmetry $\mf{su}(2)$. We claim that 
\begin{equation}\label{eq:schur_D2n+1_pochh_form}
\CI_{\mathsf{S}}^{(A_1,D_{2n+1})}(q,1)=\left(\frac{\left(q^{2n+1} ; q^{2n+1}\right)_{\infty}}{(q ; q)_{\infty}}\right)^3~.    
\end{equation}
We have 
\begin{equation}
 \CI_{\mathsf{S}}^{(A_1,D_{2n+1})}(q,1)=\mr{PE}\left[3\frac{q-q^{2 n+1}}{(1-q)\left(1-q^{2 n+1}\right)} \right]   
\end{equation}
Notice that
\begin{equation}
\begin{split}
\frac{3(q-q^{2 n+1})}{(1-q)(1-q^{2 n+1})}&=\left(\frac{3 q}{1-q}-\frac{3 q^{2 n+1}}{1-q^{2 n+1}}\right) ~.
\end{split}
\end{equation}
Using the formula 
\begin{equation}
    \mr{PE}\left[\frac{kq}{1-q}\right]=\frac{1}{(q;q)^k_\infty} ~,\quad \PE[f-g]=\frac{\PE[f]}{\PE[g]}~,
\end{equation}
we get
\begin{equation}
    \CI_{\mathsf{S}}^{(A_1,D_{2n+1})}(q,1)=\mr{PE}\left[3\frac{q-q^{2 n+1}}{(1-q)\left(1-q^{2 n+1}\right)} \right]  =\left(\frac{\left(q^{2n+1} ; q^{2n+1}\right)_{\infty}}{(q ; q)_{\infty}}\right)^3~.
\end{equation}
Taking $t\to q$ in \eqref{eq:Mac_ind_A1D3}, we get 
\begin{equation}
    \CI_{\mr{S}}^{(A_1,D_3)}(q)=\sum_{m=0}^\infty 
    \frac{q^m}{(q;q)_m} \sum_{k=0}^{2m}\binom{2m}{k}_q~.    
\end{equation}
Thus, proving that our formula \eqref{eq:Mac_ind_A1D3} satisfies the correct Schur limit amounts to proving the following product-sum identity:
\begin{equation}\label{eq:new_idn1}
\sum_{m=0}^\infty 
    \frac{q^m}{(q;q)_m}\sum_{k=0}^{2m}\binom{2m}{k}_q=\left(\frac{\left(q^3 ; q^3\right)_{\infty}}{(q ; q)_{\infty}}\right)^3~.    
\end{equation}
We now prove a more general identity which implies \eqref{eq:new_idn1}.
\begin{thm}\label{thm:main_id}
We have the identity

\begin{equation}\label{eq:main_id}
\frac{\left(z ; q^3\right)_{\infty}\left(q^3 / z ; q^3\right)_{\infty}\left(q^3 ; q^3\right)_{\infty}}{(z ; q)_{\infty}(q / z ; q)_{\infty}(q ; q)_{\infty}}=\sum_{n=0}^{\infty} \frac{q^n}{(q ; q)_n} \sum_{k=-n}^n{
2 n \choose
n+k}_q z^k~.
\end{equation}    
\end{thm}
Before proving this theorem, we note that plugging $z=1$ in \eqref{eq:main_id} gives \eqref{eq:new_idn1}. On the \tsf{LHS} of \eqref{eq:main_id}, we need to use 
\begin{equation}
    \lim_{z\to 1}\frac{\left(z ; q^3\right)_{\infty}}{\left(z ; q\right)_{\infty}}=\frac{\left(q^3 ; q^3\right)_{\infty}}{\left(q ; q\right)_{\infty}}~.
\end{equation}
\begin{proof}
We start with the following form of the  Jacobi's triple product identity
%
\begin{equation}
\sum_{n \in \mathds{Z}} q^{\frac{n(n+1)}{2}} z^n=(q ; q)_{\infty}\left(-1/z ; q\right)_{\infty}(-z q ; q)_{\infty}
\end{equation}
which can be obtained from \cite[Theorem 2.8, Page 21]{andrews1998theory} up on using the definition of $q$-Pochhammer symbol. Substituting $z \mapsto-z / q$, we obtain 
\begin{align}
\sum_{n \in \mathds{Z}}(-1)^n q^{n\left(\frac{n+1}{2}-1\right)} z^n&=(q ; q)_{\infty},(q / z ; q)_{\infty}(z ; q)_{\infty}~,
\\\implies
\sum_{n \in \mathds{Z}}(-1)^n q^{\binom{n}{2}} z^n&=(z ; q)_{\infty}(q / z ; q)_{\infty}(q ; q)_{\infty} .
\end{align}
So we may equivalently prove
\begin{equation}\label{eq:intermediate_main_id}
\begin{aligned}
& \sum_{N=-\infty}^{\infty}(-1)^N q^{3\binom{N}{2}} z^N=\sum_{m=-\infty}^{\infty}(-1)^m q^{\binom{m}{2}} z^m \sum_{n=0}^{\infty} \frac{q^n}{(q ; q)_n} \sum_{k=-n}^n{
2 n \choose
n+k}_q z^k .
\end{aligned}
\end{equation}
To prove this, we will show that the coefficient of $z^N$ of both sides is equal for $N \geq 0$. Then for $N < 0$, we will use the $z \mapsto z^{-1}$ symmetry on both sides of \eqref{eq:intermediate_main_id}.
The coefficient of $z^N$ on LHS of \eqref{eq:intermediate_main_id} is
\begin{equation}\label{eq:LHS_zN_coeff}
(-1)^N q^{3{N\choose 2}}~.
\end{equation}
The coefficient of $z^N$ on RHS of \eqref{eq:intermediate_main_id} is
\begin{equation}\label{eq:RHS_zN_coeff}
\begin{aligned}
& \sum_{n=0}^{\infty} \frac{q^n}{(q ; q)_n} \sum_{k=-n}^n{
2 n \choose
n+k}_q(-1)^{N-k} q^{\binom{N-k}{2}} \\
& \quad=(-1)^N q^{\binom{N}{2}} \sum_{n=0}^{\infty} \frac{q^n}{(q ; q)_n} \sum_{k=-n}^n{
2 n \choose
n+k}_q(-1)^k q^{\binom{k}{2}+k(1-N)} \\
& \quad=(-1)^N q^{\binom{N}{2}} \sum_{n=0}^{\infty} \frac{q^n}{(q ; q)_n}\left(q^{1-N} ; q\right)_n\left(q^N ; q\right)_n~,
\end{aligned}
\end{equation}
where we used
\begin{equation}
\begin{aligned}
\binom{N}{2}+\binom{k}{2}+k(1-N) & =\frac{N(N-1)}{2}+\frac{k(k-1)}{2}+k(1-N) \\
& =\frac{N(N-1)}{2}+k \frac{(k-1+2-2 N)}{2} \\
& =\frac{N^2-N+k^2+k-2 N k}{2} \\
& =\frac{(N-k)(N-k-1)}{2} \\
& =\binom{N-k}{2}
\end{aligned}
\end{equation}
and \cite[Page 49, Example 1]{andrews1998theory}
\begin{equation}\label{eq:identity_qx}
\sum_{k=-n}^n q^{k^2} x^k{
2 n \choose
n+k}_{q^2}=\left(-x^{-1} q ; q^2\right)_n\left(-x q ; q^2\right)_n~,
\end{equation}
with $q \rightarrow q^{1 / 2}, x=-q^{1-N-\frac{1}{2}}$ so that LHS of \eqref{eq:identity_qx} becomes
\begin{equation}
\sum_{k=-n}^n q^{\frac{k^2}{2}+k\left(1-N-\frac{1}{2}\right)}(-1)^k{
2 n \choose
n+k}_q=\sum_{k=-n}^n q^{\binom{k}{2}+k(1-N)}(-1)^k{
2 n \choose
n+k}_q~,
\end{equation}
and the RHS of \eqref{eq:identity_qx} becomes
\begin{equation}
\left(q^{N-1+\frac{1}{2}+\frac{1}{2}} ; q\right)_n\left(q^{1-N-\frac{1}{2}+\frac{1}{2}} ; q\right)_n=\left(q^N ; q\right)_n\left(q^{1-N} ; q\right)_n~.
\end{equation}
Now for $N=0$, all terms in the sum in the last line of \eqref{eq:RHS_zN_coeff}, except $n=0$, vanish since $(1 ; q)_n=0$. The $n=0$ term is 1 which is same as \eqref{eq:LHS_zN_coeff} for $N=0$.
For $N>0,1-N \leq 0$ and hence we can use the identity \cite[Appendix II, Eq. (II.6)]{gasper2004basic}
\begin{equation}\label{eq:phi_identity}
{ }_2 \phi_1\left(a, q^{-n} ; c ; q, q\right)=\frac{(c / a ; q)_n}{(c ; q)_n} a^n~,\quad n\geq 0~,
\end{equation}
where \cite[Chapter 1, Eq.(1.2.22)]{gasper2004basic}
\begin{equation}
{ }_2 \phi_1\left(a_1, a_2 ; c ; q, z\right)=\sum_{n=0}^{\infty} \frac{\left(a_1 ; q\right)_n\left(a_2 ; q\right)_n}{(q ; q)_n(c ; q)_n} z^n~.
\end{equation}
Using \eqref{eq:phi_identity} for $a=q^N, c=0,z=q$, we get,
\begin{equation}\label{eq:S_1_sum}
{ }_2 \phi_1\left(q^N, q^{1-N} ; 0 ; q , q\right)=\sum_{n=0}^{\infty} \frac{\left(q^N ; q\right)_n\left(q^{1-N} ; q\right)_n}{(q ; q)_n} q^n=q^{N(N-1)}=q^{N^2-N}~.
\end{equation}
Thus the last line of \eqref{eq:RHS_zN_coeff} evaluates to
\begin{align}
(-1)^N q^{\frac{N(N-1)}{2}} q^{N^2-N}=(-1)^N q^{3\binom{N}{2}}~,
\end{align}
which is same as \eqref{eq:LHS_zN_coeff}.
\end{proof}
\paragraph{An interesting identity.} While proving \eqref{eq:new_idn1}, we noticed that we can prove another interesting $q$-series identity which might be of independent interest.
\begin{thm}
The following identity holds:
\begin{equation}
\left(\frac{\left(q^{2} ; q^{2}\right)_{\infty}}{(q ; q)_{\infty}}\right)^3= \frac{1}{(q;q)_\infty(q^2;q^2)_{\infty}}\sum_{n=0}^\infty(-1)^nq^{n^2+n}\left(\frac{1+q^{2n+1}}{1-q^{2n+1}}\right)~.    
\end{equation}
\begin{proof}
Let us start by noting that 
\begin{equation}
    \frac{(z,q^2/z,q^2;q^2)_\infty}{(z,q/z,q;q)_\infty}=\frac{1}{(zq,q/z,q;q^2)_\infty}~,
\end{equation}
where 
\begin{equation}
    (a_1,a_2,\dots,a_n;q)_\infty:=(a_1;q)_\infty(a_2;q)_\infty\cdots (a_n;q)_\infty~.
\end{equation}
We used the fact that 
\begin{equation}
    \frac{(z;q^2)_\infty}{(z;q)_\infty}=\frac{1}{(zq;q^2)_\infty},\quad \frac{(q^2/z;q^2)_\infty}{(q/z;q)_\infty}=\frac{1}{(q/z;q^2)_\infty}~,\quad \frac{(q^2;q^2)_\infty}{(q;q)_\infty}=\frac{1}{(q;q^2)_\infty}~,
\end{equation}
which follows from the definition. Next, we have 
\begin{equation}
\frac{1}{(zq,q/z,q;q^2)_\infty}=\frac{1}{(q;q^2)_\infty}\sum_{m,n=0}^\infty \frac{z^{m-n}q^{m+n}}{(q^2;q^2)_m(q^2;q^2)_n}~,    
\end{equation}
where we used 
\begin{equation}
    \frac{1}{(z;q)_\infty}=\sum_{m=0}^\infty\frac{z^m}{(q;q)_\infty}~.
\end{equation}
Changing variable to $m-n=k$, we get 
\begin{equation}
\begin{split}
\frac{1}{(zq,q/z,q;q^2)_\infty}&=\frac{1}{(q;q^2)_\infty}\sum_{k\in\Z}z^k\sum_{n=0}^\infty \frac{q^{2n+|k|}}{(q^2;q^2)_{n+|k|}(q^2;q^2)_n}
\\&=\frac{1}{(q;q^2)_\infty}\sum_{k\in\Z}z^k\frac{q^{|k|}}{(q^2;q^2)_{|k|}}{ }_2 \phi_1\left[\begin{array}{c}0,0\\ q^{2|k|+2}\end{array} ; q^2, q^2\right]~,
\end{split}
\end{equation}
where we used the definition in \eqref{eq:phirs_def}. Using Heine's transform \eqref{eq:Heine_exp}, we get 
\begin{equation}
\begin{split}
\frac{1}{(zq,q/z,q;q^2)_\infty}&=\frac{1}{(q;q^2)_\infty}\sum_{k\in\Z}z^k\frac{q^{|k|}}{(q^2;q^2)^2_{\infty}}\sum_{n=0}^\infty(-1)^nq^{n^2+n+2|k|n}
\\&=\frac{1}{(q;q)_\infty(q^2;q^2)_{\infty}}\sum_{k\in\Z}z^k\sum_{n=0}^\infty(-1)^nq^{n^2+n+|k|(2n+1)}
\\&=\frac{1}{(q;q)_\infty(q^2;q^2)_{\infty}}\sum_{n=0}^\infty(-1)^nq^{n^2+n}\sum_{k\in\Z}z^kq^{|k|(2n+1)}
\\&=\frac{1}{(q;q)_\infty(q^2;q^2)_{\infty}}\sum_{n=0}^\infty(-1)^nq^{n^2+n}\left(\sum_{k=0}^\infty\left((zq^{(2n+1)})^k+(z^{-1}q^{(2n+1)})^k\right)-1\right)
\\&=\frac{1}{(q;q)_\infty(q^2;q^2)_{\infty}}\sum_{n=0}^\infty(-1)^nq^{n^2+n}\left(\frac{1}{1-zq^{2n+1}}+\frac{1}{1-z^{-1}q^{2n+1}}-1\right)
\\&=\frac{1}{(q;q)_\infty(q^2;q^2)_{\infty}}\sum_{n=0}^\infty(-1)^nq^{n^2+n}\frac{1-q^{4n+2}}{(1-zq^{2n+1})(1-z^{-1}q^{2n+1})}~.
\end{split}    
\end{equation}
We then have 
\begin{equation}
\begin{split}
    \left(\frac{\left(q^{2} ; q^{2}\right)_{\infty}}{(q ; q)_{\infty}}\right)^3&=\lim_{z\to 1}\frac{(z,q^2/z,q^2;q^2)_\infty}{(z,q/z,q;q)_\infty}=\lim_{z\to 1}\frac{1}{(zq,q/z,q;q^2)_\infty}
    \\&=\frac{1}{(q;q)_\infty(q^2;q^2)_{\infty}}\sum_{n=0}^\infty(-1)^nq^{n^2+n}\frac{1-q^{4n+2}}{(1-q^{2n+1})^2}
    \\&=\frac{1}{(q;q)_\infty(q^2;q^2)_{\infty}}\sum_{n=0}^\infty(-1)^nq^{n^2+n}\left(\frac{1+q^{2n+1}}{1-q^{2n+1}}\right)~.
\end{split}    
\end{equation}
\end{proof}
\end{thm}
\subsubsection{MacDonald Index From Hilbert Series Of Arc Space}\label{sec:Hilb_arc_MacD3}
In this section, we check our formula for the MacDonald index by computing the first three terms in the $t$-expansion of the Hilbert series of the arc space of the Zhu algebra of the chiral algebra associated to the $(A_1,D_3)$ Argyres-Douglas theory. This supports our formula following the discussion in Section \ref{sec:Mac_Hil_ser_arc}. 

The associated chiral algebra of $(A_1,D_{2n+1})$ theory is \cite{beem2014infinite,Beem:2017ooy} 
\begin{equation}
    \CV_{(A_1,D_{2n+1})}=\mf{su}(2)_{-\frac{4n}{2n+1}}~.
\end{equation}
To get the Zhu algebra of $\mf{su}(2)_{-\frac{4n}{2n+1}}$, we note that this \tsf{VOA}ahs three strong generators, namely, the $\mf{su}(2)$ currents $j^0,j^\pm$. The stress tensor is given by the Sugawara construction 
\begin{equation}
    T\sim j^+j^++j^-j^-+j^0j^0~.
\end{equation}
For $n=1,2$, there are four null vectors \cite{Beem:2017ooy} given by 
\begin{align}
j^aT&=0,\quad n=1~,
\\
j^aT^2&=0,\quad n=2~,
\end{align}
and 
\begin{align}
T^2&=0~,\quad n=1
\\
T^3&=0~,\quad n=2~.
\end{align}
We will change the basis for $\mf{su}(2)$ and take the generators to be 
\begin{equation}\label{eq:basis_change_su(2)}
    j^x:=j^++ij^-,\quad j^z:=j^+-ij^-,\quad j^y:=ij^0~.
\end{equation}
Mapping 
\begin{equation}\label{eq:var_su(2)_zhu}
    j^x\mapsto X,\quad j^y\mapsto Y,\quad j^z\mapsto Z,\quad T\mapsto J~,
\end{equation}
we see that the Zhu algebra of $\mf{su}(2)_{-\frac{4}{3}}$ is given by 
\begin{equation}
    R:=A\left(\CV_{(A_1,D_{3})}\right)\cong\frac{\C[X,Y,Z,J]}{\langle f_1,f_2,f_3,f_4,f_5\rangle}~,
\end{equation}
where 
\begin{equation}
    f_1:=XZ-Y^2-J,\quad f_2:=JX,\quad f_3:=JY,\quad f_4:=JZ,\quad f_5:=J^2~.
\end{equation}
The grading on the Zhu algebra is given by 
\begin{equation}
    \mr{wt}(X,Y,Z)=1,\quad \mr{wt}(J)=2~.
\end{equation}
To compute the Hilbert series of the arc space of $R$, we start with power series 
\begin{equation}
    X(x)=\sum_{n=0}^\infty X_nx^n~,\quad Y(x)=\sum_{n=0}^\infty Y_nx^n~,\quad Z(x)=\sum_{n=0}^\infty Z_nx^n~,\quad J(x)=\sum_{n=1}^\infty J_nx^n~.
\end{equation}
Note that the power series for $J$ starts from $n=1$ owing to the fact that the $J$ represents the stress tensor in the physical theory. 
Define 
\begin{equation}
    R_n:=\C[X_0,\dots,X_n,Y_0,\dots,Y_n,Z_0,\dots,Z_n,J_1,\dots,J_n]/I_n~,
\end{equation}
where 
\begin{equation}
    I_n:=\left\langle f_1(x),\dots,f_5(x)\equiv 0 \bmod x^{n+1}\right\rangle~,
\end{equation}
and
\begin{equation}
    f_i(x):=f_i(X(x),Y(x),Z(x),J(x))~,\quad i=1,\dots,5~.
\end{equation}
% We have
% \begin{equation}
%     R_\infty:=\lim_{\leftarrow }R_n~.
% \end{equation}
% The bigrading on $R_\infty$ is given by
% \begin{equation}
%     \mathsf{wt}(X_i,Y_i,Z_i,J_i)=(i,1)~,\quad \mathsf{wt}(1)=(0,0)~.
% \end{equation}
% The Hilbert series is then defined as 
% \begin{equation}
%     H_R(q,t)=\sum_{n,m\geq 0} \mathsf{dim}(R_\infty)_{n,m}q^nt^m~.
% \end{equation}
For $n=0,1$, we obtain
\begin{equation}
\begin{aligned}
R_0 & =\mathds{C}\left[X_0, Y_0, Z_0\right] /\left\langle X_0 Z_0-Y_0^2\right\rangle~, \\
R_1 & =\mathds{C}\left[X_0, Y_0, Z_0, J_1, X_1, Y_1, Z_1\right] / I_1~,
\end{aligned}
\end{equation}
where 
\begin{equation}
    I_1 = \langle X_0 Z_0 - Y_0^2, X_0 Z_1-2 Y_0 Y_1+Z_0 X_1-J_1, X_0 J_1,Y_0 J_1,Z_0J_1 \rangle~.
\end{equation}
The arc space of $R$ is defined by
\begin{equation}
    R_{\infty}:=\lim _{ \leftarrow } R_n~.
\end{equation}
We assign bigrading to the $n$-jet space $R_n$ by defining
\begin{equation}
\begin{split}
\mathsf{wt}^{(2)}\left(X_i\right)&=\left(\mathsf{wt}_1\left(X_i\right),\mathsf{wt}_2\left(X_i\right)\right)=(i, 1)~,\quad i\geq 0~, 
\\
\mathsf{wt}^{(2)}\left(Y_i\right)&=\left(\mathsf{wt}_1\left(Y_i\right),\mathsf{wt}_2\left(Y_i\right)\right)=(i, 1)~,\quad i\geq 0~, 
\\
\mathsf{wt}^{(2)}\left(Z_i\right)&=\left(\mathsf{wt}_1\left(Z_i\right),\mathsf{wt}_2\left(Z_i\right)\right)=(i, 1)~,\quad i\geq 0~, 
\\
\mathsf{wt}^{(2)}\left(J_i\right)&=\left(\mathsf{wt}_1\left(J_i\right),\mathsf{wt}_2\left(J_i\right)\right)=(i, 1)~,\quad i\geq 1~, 
\\
\mathsf{wt}^{(2)}\left(1\right)&=\left(\mathsf{wt}_1\left(1\right),\mathsf{wt}_2\left(1\right)\right)=(0, 0)~.
\end{split}
\end{equation}
We denote $(R_n)_{\ell,m}$, the homogeneous subspace of weight $(\ell,m)$.
The Hilbert series of $R_n$ is defined as
\begin{equation}
H_{R_n}(q, t)=\sum_{\ell=0}^n \sum_{m \geq 0} \mr{dim}\left(R_n\right)_{\ell, m} q^{\ell} t^m~.
\end{equation}
Taking $n \rightarrow \infty$, we get the Hilbert series of  $R_{\infty}$
\begin{equation}
H_{R_{\infty}}(q, t)=\sum_{\ell, m \geq 0} \mr{dim}\left(R_{\infty}\right)_{\ell, m} q^{\ell} t^m~.
\end{equation}
% One way of computing the Hilbert series of an $n$-jet scheme $R_n$ is by using Gr\"{o}bner basis of $I_n$.
% Let us denote the Gr\"{o}bner basis of $I_n$ by $G_n$. Note that $G_n$ depends on the choice of ordering. Let us denote the set of leading monomials of the polynomials in $G_n$ by
% \begin{equation}
% \mr{LM}(G_n):=\{\mr{LM}(g)\mid g\in G_n\}~.
% \end{equation}
% \myRed{Need to explain the formula.}
% \\
From the general formula \eqref{eq: hilbert_ser_Rn_gen} (proved in Appendix \ref{app:Hilb_ser_count}), the Hilbert series of $R_n$ for the case at hand is given by
\begin{equation}
\label{eq: hilbert_ser_Rn}
H_{R_n}(q,t) = \mr{Ser}_{q=0,t=0}\frac{1+\sum_{S\subseteq \mr{LM}(G_n)}(-1)^{|S|}(q,t)^{\mr{wt}^{(2)}(\mr{LCM}(S))}}{(1-t)^3 \prod_{i=1}^{n}\left(1-q^i t\right)^4} \quad (\mr{mod} \;q^{n+1})~,
\end{equation}
where $\mr{Ser}_{q=0,t=0}(f(q,t))$ denotes the power series of $f(q,t)$ around $q=t=0$, the summation in the numerator is over all subsets of monomials in $\mr{LM}(G_n)$, $|S|$ denotes the number of monomials in $S$, $\mr{LCM}(S)$ is least common multiple of the monomials in $S$, and we have defined
\begin{equation}
(q,t)^{(\ell,m)} := q^{\ell} t^m ~.   
\end{equation}
The Hilbert series of $R_\infty$ can then be given by taking $n\to \infty$
\begin{equation}
\label{eq: hilbert_ser_Rinf}
H_{R_\infty}(q,t) = \mr{Ser}_{q=0,t=0}\frac{1+\sum_{S\subseteq \mr{LM}(G_\infty)}(-1)^{|S|}(q,t)^{\mr{wt}^{(2)}(\mr{LCM}(S))}}{(1-t)^3 \prod_{i=1}^{\infty}\left(1-q^i t\right)^4}~.
\end{equation}
To apply \eqref{eq: hilbert_ser_Rn}, we need a general form for Gr\"{o}bner basis.
Through Mathematica computations, we have the following conjecture for $\mr{LM}(G_n)$:
\begin{conj}
Choose the ordering of variables to be
\begin{equation}\label{eq:var_ord}
J_1>J_2>\cdots>X_0>X_1>\cdots>Y_0>Y_1>\cdots>Z_0>Z_1>\ldots~.
\end{equation}
Then with the degree lexicographic (deglex) ordering of monomials (see Appendix \ref{app:Grob}), the leading terms of the Gr\"{o}bner basis of  is given by 
\begin{equation}
\label{eq:Gr\"{o}bner_basis_Gn}
\mr{LM}\left(G_n\right)=\bigcup_{\substack{0 \leq \ell \leq n \\ 2 \leq m \leq \ell+2}} I_{\ell, m}~,
\end{equation}
where $I_{\ell,m}$ is a homogeneous set of weight $(\ell,m)$ given by
\begin{equation}
I_{\ell,m} := \left(\{X_{m-2}\}\cdot F_z(\ell-m+2, m-1)\right)\bigcup \left(\{J_{m-1}\}\cdot F(\ell-m+1, m-1)\right)~,
\end{equation}
where 
\begin{enumerate}
\item for sets $A,B$, we have introduced the notation
\begin{equation}
    A\cdot B:=\{ab:a\in A,b\in B\}~,
\end{equation}
\item for $\ell,m\in\IZ_{\geq 0}$, we have defined
\begin{equation}
F_z(\ell,m) := \{Z_{i_1}\dots Z_{i_m}: i_1+\dots +i_m = \ell\}~,
\end{equation}
\item for $\ell,m\in\IZ_{\geq 0}$, $F(\ell,m)$ is the set of all degree $m$ monomials that can be formed using the variables $X_0,\dots,X_\ell,Y_0\dots,Y_\ell,Z_0,\dots,Z_\ell,J_1,\dots,J_\ell$. More explicitly, 
\begin{equation}
\begin{aligned}
F(0,0) &= \{1\}~, \\
F(\ell,0) &= \emptyset, \quad \text{for}\quad \ell\neq 0~,\\
\end{aligned}
\end{equation}
and for $m>0$,
\begin{equation}
F(\ell,m):=\left\{X_0^{n_{X_0}}Y_0^{n_{Y_0}}Z_0^{n_{Z_0}} \prod_{i=1}^\ell X_i^{n_{X_i}}Y_i^{n_{Y_i}}Z_i^{n_{Z_i}}J_i^{n_{J_i}}: \sum_{i=0}^\ell (n_{X_i}+n_{Y_i}+n_{Z_i})+\sum_{i=1}^\ell n_{J_i}=m\right\}~.
\end{equation} 
\end{enumerate}
\end{conj}
\par
We now work order by order in $t$. At order $t^m$, the only terms in the numerator that contribute are the constraints $S\subset\mr{LM}(G_\infty)$ which have weight $(\ell,m)$.
\paragraph{Order $t^0$.} At order $t^0$, we only have one generator 1. Thus, the series starts with 1. 
\paragraph{Order $t^1$.} From \eqref{eq:Gr\"{o}bner_basis_Gn}, we see there is no constraint $S\subset\mr{LM}(G_\infty)$  weight $(\ell,1)$. Then from \eqref{eq: hilbert_ser_Rinf}, the contribution for the first order in $t$ can be given by
\begin{equation}\label{eq:t1_comp_D3}
    \begin{aligned}
H_{R_\infty}(q,t) &= \mr{Ser}_{q=0,t=0}\frac{1}{(1-t)^3 \prod_{i=1}^{\infty}\left(1-q^i t\right)^4}\quad (\mr{mod} \;t^2)\\
& = 1 + t\bigg(3+4\sum_{i=1}^\infty q^i\bigg)\\
& = 1+t\left(\frac{3+q}{1-q}\right) \\
& = 1+\frac{t}{(q ; q)_1} \sum_{k=0}^2\binom{2}{k}_q~,
\end{aligned}
\end{equation}
where the last equality follows from a straightforward Mathematica calculation. 
\paragraph{Order $t^2$.} The constraints $S\subset\mr{LM}(G_\infty)$ of weight $(\ell,m=2)$ are given by
\begin{equation}
I_{\ell,m=2} = \{X_{0}\}\times F_z(\ell, 1)\bigcup \{J_{1}\}\times F(\ell-1, 1)~.
\end{equation}
More explicitly, these are 
\begin{equation}
\begin{aligned}
I_{\ell=0, m = 2} &=\{X_0Z_0\}~,\\
I_{\ell=1, m= 2}  &= \{X_0Z_1, J_1 X_0,J_1Y_0,J_1Z_0\}~,\\
I_{\ell, m=2}  &= \{X_0Z_\ell, J_1 X_\ell,J_1Y_\ell,J_1Z_\ell, J_1J_\ell\}, \quad \ell>1~.
\end{aligned}
\end{equation}
The Hilbert series to the second order in $t$ can thus be obtained from \eqref{eq: hilbert_ser_Rn}:
\begin{equation}
\label{eq:Hil_ser_t2}
\begin{aligned}
H_{R_\infty}(q,t) &= \mr{Ser}_{q=0,t=0}\frac{1-t^2(1+4q + 5\sum_{i=2}^\infty q^i)}{(1-t)^3 \prod_{i=1}^{\infty}\left(1-q^i t\right)^4}\quad (\mr{mod} \;t^3).
\end{aligned}
\end{equation}
Focusing on the $t^2$ term,
\begin{equation}
H_{R_{\infty}}(q, t)|_{t^2}=\sum_{\ell \geq 0} \mr{dim}\left(R_{\infty}\right)_{\ell, 2} q^\ell~,
\end{equation}
since the denominator is the generating function for all monomials, the coefficients of $t^2$ term in \eqref{eq:Hil_ser_t2} is given by
\begin{equation}\label{eq:dimRl2}
\mr{dim}\left(R_{\infty}\right)_{\ell, 2} = |F(\ell,2)|-|F(\ell,0)|-4|F(\ell-1,0)|-5\sum_{i=2}^\infty|F(\ell-i,0)|~.
\end{equation}
Note that 
\begin{equation}
\begin{aligned}
|F(\ell,0)| = \delta_{\ell,0}~.
\end{aligned}
\end{equation}
Our goal is to compute $|F(\ell,2)|$. 
The first few set of free monomials is given  by
\begin{equation}
\begin{aligned}
F(0,2) &= \{X_0^2, Y_0^2,Z_0^2,X_0Y_0,X_0Z_0,Y_0Z_0\}~,\\
F(1,2) &= \{X_0, Y_0, Z_0\}\cdot \{X_1,Y_1,Z_1,J_1\}~,\\
F(2,2) & = \{X_0,Y_0,Z_0\}\cdot \{X_2,Y_2,Z_2,J_2\}\cup\{X_1^2, Y_1^2,Z_1^2,J_1^2,X_1Y_1,\dots\}~,\\
F(3,2) & = \{X_0,Y_0,Z_0\}\cdot \{X_3,Y_3,Z_3,J_3\}\cup\{X_1,Y_1,Z_1,J_1\}\cdot \{X_2,Y_2,Z_2,J_2\}~,\\
&\dots
\end{aligned}
\end{equation}
% where for sets $A,B,C,D$, we introduced the notation
% \begin{align}
%     A\cdot B\cup C&:=\{ab+c:a\in A,b\in B,c\in C\}~,
%     \\
%     A\cdot B\cup C\cdot D&:=\{ab+cd:a\in A,b\in B,c\in C,d\in D\}~.
% \end{align}
We anticipate the formula
\begin{equation}
\label{eq:num_Fl2}
|F(\ell,2)|=\begin{cases}
8\ell + 6, & \text{if $\ell$ is even}\\
8\ell + 4, & \text{if $\ell$ is odd}~.
\end{cases}
\end{equation}
We combine two cases 
\begin{equation}\label{eq:free_words_t2}
|F(\ell,2)| = 8\ell + 5 + (-1)^\ell~.
\end{equation}
Let us prove this. Note that
a monomial of total weight $(\ell, 2)$ must satisfy:
\begin{equation}
\sum (\text{indices} \times \text{exponents}) = \ell \quad \text{and} \quad \sum (\text{exponents}) = 2.
\end{equation}
We analyze two cases:
\\
\textbf{Case 1: Squared Variables.}
A monomial of the form $V_i^2$ contributes weight $(2i, 2)$. For this to equal $(\ell, 2)$, we require $2i = \ell$, i.e., $i = \ell/2$. 
\begin{itemize}
    \item If $\ell$ is even, $i = \ell/2$ is an integer. The variables available are $X_{\ell/2}, Y_{\ell/2}, Z_{\ell/2}$ (3 options) and $J_{\ell/2}$ (if $\ell/2 \geq 1$, since $J_0$ doesn't exist). 
    \begin{equation}
    \text{Count} = 
    \begin{cases} 
    4 & \text{if } \ell \geq 2 \text{ even}, \\
    3 & \text{if } \ell = 0.
    \end{cases}
    \end{equation}
    \item If $\ell$ is odd, no such $i$ exists. Count $= 0$.
\end{itemize}
\textbf{Case 2: Product of Two Distinct Variables.}
A monomial $V_i W_j$ with $i \neq j$ contributes weight $(i + j, 2)$. For $i + j = \ell$, we count pairs $(V_i, W_j)$:
\begin{itemize}
    \item For $i = 0$ or $j=0$: Variables $X_0, Y_0, Z_0$ (3 choices). Partner $W_j$ must satisfy $j = \ell$. For $j \geq 1$, there are 4 variables ($X_j, Y_j, Z_j, J_j$).
    \item For $i \geq 1$: Both $V_i$ and $W_j$ have 4 choices each ($X, Y, Z, J$).
\end{itemize}
Total count for this case:
\begin{equation}
\text{Count} = 
\begin{cases} 
\frac{1}{2}(3 \times 4) + \frac{n}{2}(4 \times 4) & \text{if } \ell \text{ even}, \\
(3 \times 4) + (4 \times 4) \times \left(\frac{\ell - 1}{2}\right) & \text{if } \ell \text{ odd}.
\end{cases}
\end{equation}
Simplifying, we get
\begin{equation}
\text{Count} = 
\begin{cases} 
8\ell + 6 & \text{if } \ell \text{ even}, \\
8\ell + 4 & \text{if } \ell \text{ odd}~.
\end{cases}
\end{equation}
This completes the proof. 
Using \eqref{eq:num_Fl2} and \eqref{eq:dimRl2}, we have 
\begin{equation}
H_{R_\infty}(q,t)|_{t^2} = t^2\left(5+8 q+\sum_{n \geq 2}\left(8 n+(-1)^n\right) q^n\right)~.
\end{equation}
Noting that
\begin{equation}
\sum_{n \geq 1} n q^n=\frac{q}{(1-q)^2}, \quad \sum_{n \geq 1}(-1)^n q^n=-\frac{q}{1+q}~.
\end{equation}
We have
\begin{equation}\label{eq:I_2-1_calc}
\begin{aligned}
5+8 q+\sum_{n \geq 2}\left(8 n+(-1)^n\right) q^n & =5+8 q+\left(\frac{q}{(1-q)^2}-q\right)+\left(-\frac{q}{1+q}+q\right) \\
& =4+q+\frac{8 q}{(1-q)^2}+\frac{1}{1+q}~.
\end{aligned}
\end{equation}
Using straightforward Mathematica calculation, we find that 
\begin{equation}\label{eq:I_2-2_calc}
H_{R_\infty}(q,t)|_{t^2} = \frac{1}{(q ; q)_2} \sum_{k=0}^4\binom{4}{k}_q.
\end{equation}
Because of the techincality of the discussion, we present the evaluation of the Hilbert series at order $t^3$ in Appendix \ref{app:Hilb_ser_t3}. 
% At order $t$, we have generators $X_0,\dots,Y_0,\dots,Z_0,\dots,J_1,\dots$. At $q^0$, we have three generators $X_0,Y_0,Z_0$ without generators. At order $q^n$ with $n>0$, we have four generators $X_n,Y_n,Z_n,J_n$ without any relations. Thus we have
% \begin{equation}
% \begin{split}  H_R(q,t)&=1+t\left(3+4\sum_{n=1}^\infty q^n\right)+O(t^2)
% \\&=1+t\left(3+\frac{4}{1-q}-4\right)+O(t^2)
% \\&=1+t\left(\frac{3+q}{1-q}\right)+O(t^2)~.
% \end{split}  
% \end{equation}
%  From Mathematica computations, we noticed that at order $t^m$, if we choose the following ordering of variables 
% \begin{equation}
%     J_1>J_2>\dots>X_0>X_1>\dots>Y_0>Y_1>\dots>Z_0>Z_1>\dots~,
% \end{equation}
% then leading term of the Gr\"{o}bner basis for $\langle f_1(x),\dots,f_5(x)\rangle$ is given by 
% \begin{equation}
%     I_m:=\bigcup_{n\geq 0}I_{m,n}~,
% \end{equation}
% where 
% \begin{equation}
%     I_{m,n}=\begin{cases}
%         \emptyset; & m-n>2 \\
%         \{X_nZ_0^{m-1}\}; &m-n=2 
%     \end{cases}~,
% \end{equation}
% and for $m-n<2$, we have 
% \begin{equation}\label{eq:Im-n<2}
%     I_{m,n}=\left\{X_{m-2}Z_{i_1}\dots Z_{i_{m-1}},X_{\vec{w}}^{(i_1,\dots,i_{m-1})}:\begin{array}{c}
%      i_1+\dots+i_{m-1}=n-m+1      \\
%       \vec{w}=(w_1,\dots,w_{m-1})\in(\mathds{F}_2^4)^{m-1}~, \mathsf{wt}(w_i)=1   
%     \end{array} \right\}~,
% \end{equation}
% where for $\vec{w}=(w_1^1,w_1^2,w_1^3,w_1^4;\dots,w_{m-1}^1w_{m-1}^2w_{m-1}^3w_{m-1}^4)\in (\mathds{F}_2^4)^{m-1}$ we introduced the notation
% \begin{equation}
% X_{\vec{w}}^{(i_1,\dots,i_{m-1})}:=J_{m-1}\left( X_{i_1}^{w_1^1}Y_{i_1}^{w_1^2}Z_{i_1}^{w_1^3}J_{i_1}^{w_1^4}\cdots X_{i_{m-1}}^{w_{m-1}^1}Y_{i_{m-1}}^{w_{m-1}^2}Z_{i_{m-1}}^{w_{m-1}^3}J_{i_{m-1}}^{w_{m-1}^4}\right)~.  
% \end{equation}
% At order $t^2$, the Gr\"{o}bner basis for $\langle f_1(x),\dots,f_5(x)\rangle$ is given by 
% \begin{equation}
%     I_2=\{X_0Z_0,X_0Z_n,J_1X_n,J_1Y_n,J_1Z_n,J_1J_n:n\geq 1\}~.
% \end{equation}
% At $q^0$, there are three generators $X_0,Y_0,Z_0$. Thus there are 6 independent words, namely
% \begin{equation}
%     X_0^2,Y_0^2,Z_0^2,X_0Y_0,X_0Z_0,Y_0Z_0~.
% \end{equation}
% There is one constraint $Z_0X_0$. Thus $\mathsf{dim}(R_\infty)_{0,2}=6-1=5$. At order $q$, there are $7$ generators and hence 12 words, namely,
% \begin{equation}
% \begin{split}
%     X_0X_1,X_0Y_1,X_0Z_1,X_0J_1
%     \\
%     Y_0X_1,Y_0Y_1,Y_0Z_1,Y_0J_1
%     \\
%     Z_0X_1,Z_0Y_1,Z_0Z_1,Z_0J_1~.
% \end{split}    
% \end{equation}
% There are 4 constraints coming from \eqref{eq:Im-n<2}:
% \begin{equation}
% X_0Z_1,J_1X_0,J_1Y_0,J_1Z_0~.    
% \end{equation}
% So $\mathsf{dim}(R_\infty)_{1,2}=12-4=8$. Thus we have 
% \begin{equation}
% H_R(q,t)=1+t\left(\frac{3+q}{1-q}\right)+t^2(5+8q+O(q^2))+O(t^3)~.    
% \end{equation}

% We now claim that for $n\geq 0$, the number of free words with bigrade $(n,2)$ is $8n+5+(-1)^n$. 
% The two cases can be combined and written as 
% \begin{equation}
% \boxed{8n + 5 + (-1)^n}
% \end{equation}
% At order $q^2$, there are 5 constraints obtained from \eqref{eq:Im-n<2}:
% \begin{equation}
%     X_0Z_2, J_1X_1,J_1Y_1,J_1Z_1,J_1^2~.
% \end{equation}
% Thus, we have 
% \begin{equation}
% \mathsf{dim}(R_\infty)_{n,2}  =8n+(-1)^n,\quad n\geq 2~.  
% \end{equation}
% We can then write
% \begin{equation}
% H_R(q,t)=1+t\left(3+4\sum_{n=1}^\infty q^n\right)+t^2(5+8q+\sum_{n\geq 2}(8n+(-1)^n)q^n)+O(t^3)~.    
% \end{equation}
% Noting that 
% \begin{equation}
%     \sum_{n\geq 1}nq^n=\frac{q}{(1-q)^2},\quad \sum_{n\geq 1}(-1)^nq^n=-\frac{q}{1+q}~,
% \end{equation}
% we have 
% \begin{equation}
% \begin{split}
%     5+8q+\sum_{n\geq 2}(8n+(-1)^n)q^n&=5+8q+\left(\frac{q}{(1-q)^2}-q\right)+\left(-\frac{q}{1+q}+q\right)
%     \\
%     &=4 + q + \frac{8q}{(1 - q)^2} + \frac{1}{1 + q}~.
% \end{split}    
% \end{equation}
% Lastly, we note that 
% \begin{equation}
% \frac{1}{(q;q)_1}\sum_{k=0}^2{2\choose k}_q=\frac{3+q}{1-q}~.    
% \end{equation}
% and 
% \begin{equation}
% \begin{split}
%     \frac{1}{(q;q)_2}\sum_{k=0}^4{4\choose k}_q&=4+\frac{8}{(1-q)^2}-\frac{1}{1-q}+q+\frac{1}{1+q}\\&=4 + q + \frac{8q}{(1 - q)^2} + \frac{1}{1 + q}~.
% \end{split}
% \end{equation}
% Thus, we have 
% \begin{equation}
% H_R(q,t)=1+\frac{t}{(q;q)_1}\sum_{k=0}^2{2\choose k}_q+\frac{t^2}{(q;q)_2}\sum_{k=0}^4{4\choose k}_q+O(t^3)~.    
% \end{equation}
% This matches the first three terms of our formula \eqref{eq:Mac_ind_A1D3}. 


\section{MacDonald Index For $(A_1,D_{2n+1})$ Theory}\label{sec:A1D2n+1_Mac_conj}
In this section, we conjecture a generalization of formula \eqref{eq:Mac_ind_A1D3} for the MacDonald index for the general $(A_1,D_{2n+1})$ theory. We motivate our formula using some observations from the Hilbert series of arc space for the Zhu algebra of chiral algebra of $(A_1,D_{5})$ theory and its relation to the MacDonald index of projected free hypermultiplet. We then check the consistency of our conjecture with the Schur limit. This requires us to prove a generalization of the product-sum formula \eqref{eq:new_idn1}, see \eqref{eq:new_id_gen_n} below. 
Let us first state the conjecture.
\begin{conj}
The MacDonald index for the $(A_1, D_{2n+1})$ Argyres-Douglas theory is given by
\begin{equation}\label{eq:Mac_ind_gen_n}
    \CI^{(A_1,D_{2n+1})}_{\mr{M}} = 1 + \sum_{m=1}^\infty 
    t^m \sum_{l=1}^m
    \frac{f^n_{m,l}(q)}{(q;q)_l}\sum_{k=0}^{2l}\binom{2l}{k}_q~,
\end{equation}
where $f^n_{m,l}(q)$ is a class of polynomial functions in $q$ which satisfy the hypergeometrical relation 
\begin{equation}\label{eq:hyper_fmnl_rel}
    1 + \sum_{m=1}^\infty q^m \sum_{l=1}^m 
    \frac{f^n_{m,l}(q)}{(q;q)_l} (q^N; q)_l 
    (q^{1-N};q)_l = q^{nN(N-1)}~.
\end{equation}
\end{conj}
As we explain in Section \ref{sec:schur_gen_n}, the condition \eqref{eq:hyper_fmnl_rel} on $f^n_{m,l}$ makes sure that our proposed formula satisfy the Schur limit. 
In Appendix \ref{app:new_id}, we prove that the functions
\begin{equation}\label{eq:fnml_closed_form_sec5}
f^k_{m,l}(q)=\sum_{i_k,i_{k-1},\dots,i_3\geq 0}\frac{q^{Q(m-(k-1)i_k-(k-2)i_{k-1}-\dots-l,i_3,\dots,i_k)}(q;q)_{l}}{(q;q)_{i_k}\dots (q;q)_{i_3}(q;q)_{m-(k-1)i_k-(k-2)i_{k-1}-\dots-l}(q;q)_{2l+i_3+2i_4+\dots+(k-2)i_k-m}}~,    
\end{equation}
where 
\begin{equation}\label{eq:Q_def_sec5}
    Q(i_2,i_3,\dots,i_k):=i_k^2+(i_k+i_{k-1})^2+\dots+(i_k+\dots+i_3+i_2)^2~,
\end{equation}
satisfy \eqref{eq:hyper_fmnl_rel}. It was checked on Mathematica that with $f^k_{m,l}$ given above, the proposed MacDonald index \eqref{eq:Mac_A1D2n+1_conj} reproduces the known ccoefficients in Tables \ref{tab:coeff_mac_A1D5} -- \ref{tab:coeff_mac_A1D9}.
%\myRed{Remark about RG flow and reference to App. F.}
\subsection{Motivating The Formula Using Hilbert Series}
As discussed in Section \ref{sec:RG_flow} and illustrated in Section \ref{sec:D3_free_hyper_rel}, the MacDonald index for $(A_1,D_{2n+1})$ theory must be related to that of the (projected) free hypermultiplet. %The relation is more evident from the , the numbers above the diagonal are in perfect match with the numbers above diagonal in Table \ref{tab:coeff_mac_FH}.
\begin{table}[]
    \centering
    \begin{tabular}{|l|lllllllll|}
\hline Variables & $t^0$ & $t^1$ & $t^2$ & $t^3$ & $t^4$ & $t^5$ & $t^6$ & $t^7$ & $t^8$ \\
\hline
$q^0$ & 1 & 3 & 5 & 7 & 9 & 11 & 13 & 15 & 17 \\
$q^1$ & 0 & 4 & 11 & 17 & 23 & 29 & 35 & 41 & 47 \\
$q^2$ & 0 & 4 & 21 & 40 & 56 & 72 & 88 & 104 & 120 \\
$q^3$ & 0 & 4 & 27 & 73 & 114 & 150 & 186 & 222 & 258 \\
$q^4$ & 0 & 4 & 37 & 119 & 218 & 303 & 381 & 459 & 537 \\
$q^5$ & 0 & 4 & 43 & 173 & 361 & 546 & 709 & 863 & 1017 \\
$q^6$ & 0 & 4 & 53 & 243 & 583 & 942 & 1279 & 1586 & 1882 \\
\hline
\end{tabular}
    \caption{Table of selected coefficients of the Macdonald index for $(A_1,D_5)$ theory.}
    \label{tab:coeff_mac_A1D5}
\end{table}
\begin{table}[]
    \centering
    \begin{tabular}{|l|lllllllll|}
\hline Variables & $t^0$ & $t^1$ & $t^2$ & $t^3$ & $t^4$ & $t^5$ & $t^6$ & $t^7$ & $t^8$ \\
\hline $q^0$ & 1 & 3 & 5 & 7 & 9 & 11 & 13 & 15 & 17 \\
$q^1$ & 0 & 4 & 11 & 17 & 23 & 29 & 35 & 41 & 47 \\
$q^2$ & 0 & 4 & 21 & 43 & 61 & 79 & 97 & 115 & 133 \\
$q^3$ & 0 & 4 & 27 & 77 & 127 & 169 & 211 & 253 & 295 \\
$q^4$ & 0 & 4 & 37 & 123 & 243 & 349 & 443 & 537 & 631 \\
$q^5$ & 0 & 4 & 43 & 177 & 401 & 638 & 846 & 1038 & 1230 \\
$q^6$ & 0 & 4 & 53 & 247 & 639 & 1109 & 1552 & 1952 & 2332 \\
\hline
\end{tabular}
    \caption{Table of selected coefficients of the Macdonald index for $(A_1,D_7)$ theory.}
    \label{tab:coeff_mac_A1D7}
\end{table}
\begin{table}[]
    \centering
    \begin{tabular}{|l|lllllllll|}
\hline Variables & $t^0$ & $t^1$ & $t^2$ & $t^3$ & $t^4$ & $t^5$ & $t^6$ & $t^7$ & $t^8$ \\
\hline $q^0$ & 1 & 3 & 5 & 7 & 9 & 11 & 13 & 15 & 17 \\
$q^1$ & 0 & 4 & 11 & 17 & 23 & 29 & 35 & 41 & 47 \\
$q^2$ & 0 & 4 & 21 & 43 & 61 & 79 & 97 & 115 & 133 \\
$q^3$ & 0 & 4 & 27 & 77 & 130 & 174 & 218 & 262 & 306 \\
$q^4$ & 0 & 4 & 37 & 123 & 247 & 362 & 462 & 562 & 662 \\
$q^5$ & 0 & 4 & 43 & 177 & 405 & 663 & 892 & 1100 & 1308 \\
$q^6$ & 0 & 4 & 53 & 247 & 643 & 1149 & 1644 & 2089 & 2507 \\
\hline
\end{tabular}
    \caption{Table of selected coefficients of the Macdonald index for $(A_1,D_9)$ theory.}
    \label{tab:coeff_mac_A1D9}
\end{table}
Let us write down the Macdonald index for $(A_1, D_3)$ as 
\begin{equation}
\mathcal{I}_{\mathrm{M}}^{\left(A_1, D_3\right)}=\sum_{m=0}^{\infty} \frac{t^m}{(q ; q)_m} \sum_{k=0}^{2 m}\binom{2 m}{k}_q =: \sum_{m=0}^\infty t^m \mathcal{I}_m(q)~,
\end{equation}
where 
\begin{equation}
\mathcal{I}_m(q):= \frac{1}{(q ; q)_m} \sum_{k=0}^{2 m}\binom{2 m}{k}_q~.   
\end{equation}
We anticipate that the relation of the MacDonald index of $(A_1,D_{2n+1})$ theory to that of the (projected) free hypermultiplet is through the functions $\CI_m$. So, we want to express the Macdonald index of $(A_1, D_{2n+1})$ in terms of $\mathcal{I}_m$.

We start from $(A_1, D_5)$ theory, where the Zhu algebra is known for its corresponding chiral algebra $\mathcal{V}_{(A_1, D_5)}=\mathfrak{s u}(2)_{-\frac{8}{5}}$ \cite{Beem:2017ooy}\footnote{Again, we are using the basis of $\mf{su}(2)$ given in \eqref{eq:basis_change_su(2)} and \eqref{eq:var_su(2)_zhu}.}:
\begin{equation}
R^{(A_1, D_5)}:=A\left(\mathcal{V}_{\left(A_1, D_5\right)}\right) \cong \frac{\mathbb{C}[X, Y, Z, J]}{\left\langle f_1, f_2, f_3, f_4, f_5\right\rangle}~,
\end{equation}
where 
\begin{equation}\label{eq:const_D4}
f_1:= XZ - Y^2 - J, \quad  f_2:= J^2 X, \quad f_3:= J^2Y, \quad f_4:= J^2 Z, \quad f_5:= J^3~.
\end{equation}
By Theorem \ref{thm:Hil_ser_mac}, the Macdonald index for $(A_1,D_5)$ theory is equal to the Hilbert series of the arc space of the associated Zhu algebra:
\begin{equation}
\mathcal{I}^{(A_1, D_5)}_{\mr{M}}(q,t) = H_{R^{(A_1, D_5)}_{\infty}}(q, t)~.
\end{equation}
Let us compute the Hilbert series of $R^{(A_1, D_5)}_\infty$ till $O(t^3)$. We will again use the deglex monomial ordering with the ordering on variables given by \eqref{eq:var_ord} along with the general formula \eqref{eq: hilbert_ser_Rinfty_gen}.

\paragraph{Order $t^0$.} At order $t^0$, we only have one generator 1. Thus, the series starts with 1. 
\paragraph{Order $t^1$.} From \eqref{eq:const_D4}, we see that there are no elements of the Gr\"{o}bner basis of $I_n$ for any $n$, which has weight $(\ell,1)$. Thus we have 
\begin{equation}
\begin{aligned}
H_{R^{(A_1, D_5)}_{\infty}}(q, t)&=\mr{Ser}_{q=0, t=0} \frac{1}{(1-t)^3 \prod_{i=1}^{\infty}\left(1-q^i t\right)^4} \quad\left(\mr{mod} ~t^2\right)\\
& = 1+ t \mathcal{I}_1~,
\end{aligned}
\end{equation}
where we used the computations of \eqref{eq:t1_comp_D3}.
\paragraph{Order $t^2$.} From Mathematica computations, we find that the leading monomials of a Gr\"{o}bner basis $G_\infty$ with weight  $(\ell,m=2)$ are given by 
\begin{equation}\label{eq:Grob_D5t2}
\{X_0 Z_\ell:\ell\geq 0\}~.    
\end{equation}
Then using \eqref{eq: hilbert_ser_Rn_gen}, we get 
\begin{equation}
\begin{aligned}
\left.H_{R^{(A_1, D_5)}_{\infty}}(q, t)\right|_{t^2}&=\left.\mr{Ser}_{q=0, t=0} \frac{1-t^2 \sum_{\ell=0}^{\infty}q^\ell}{(1-t)^3 \prod_{i=1}^{\infty}\left(1-q^i t\right)^4}\right|_{t^2}~.
\end{aligned}
\end{equation}
Note that 
\begin{equation}
\sum_{\ell=0}^{\infty}q^\ell=1+q+\sum_{\ell=2}^\infty q^\ell=1+4q+5\sum_{\ell=2}^{\infty}q^\ell-q\left(3+4\sum_{\ell=1}^{\infty}q^\ell\right)=1+4q+5\sum_{\ell=2}^{\infty}q^\ell-q\CI_1~.    
\end{equation}
Thus using \eqref{eq:free_words_t2}, we get 
\begin{equation}
\begin{split}
\left.H_{R^{(A_1, D_5)}_{\infty}}(q, t)\right|_{t^2}&=\left. \mr{Ser}_{q=0, t=0} \frac{1-t^2 (1+ 4 q + 5\sum_{\ell=2}^{\infty}q^\ell)+t^2q\mathcal{I}_1}{(1-t)^3 \prod_{i=1}^{\infty}\left(1-q^i t\right)^4}\right|_{t^2}
\\
&=t^2\left(\sum_{\ell=0}^\infty (8\ell+5+(-1)^\ell)q^\ell-1- 4 q - 5\sum_{\ell=2}^{\infty}q^\ell+q\CI_1\right)
\\
&=t^2\left(5+8q+\sum_{\ell=2}^\infty (8\ell+(-1)^\ell)q^\ell+q\CI_1\right)
\\
& = t^2\left(\mathcal{I}_2+q\mathcal{I}_1\right)~,
\end{split}    
\end{equation}
where we used the calculation from \eqref{eq:I_2-1_calc} -- \eqref{eq:I_2-2_calc}.
\paragraph{Order $t^3$.} From Mathematica computations, we find that the leading monomials of a Gr\"{o}bner basis $G_\infty$ with weight  $(\ell,m=3)$ are given by 
\begin{align}
&\{J_1^2X_0,J_1^2 Y_0,J^2_1 Z_0\},\quad \ell=2~,
\\
&\{J_1^2 X_{\ell-2},J_1^2 Y_{\ell-2},J^2_1 Z_{\ell-2}, J_1^2 J_{\ell-2}\},\quad \ell>2~.
\end{align}
Combined with \eqref{eq:Grob_D5t2} and using \eqref{eq: hilbert_ser_Rinfty_gen}, we have
\begin{equation}
\begin{aligned}
H_{R^{(A_1, D_5)}_{\infty}}(q, t)|_{t^3} & = \mr{Ser}_{q=0, t=0} \frac{1-t^2 \sum_{\ell=0}^{\infty}q^\ell -t^3(3 q^2 +4\sum_{\ell=3}^\infty q^\ell)}{(1-t)^3 \prod_{i=1}^{\infty}\left(1-q^i t\right)^4}\bigg|_{t^3}~.
\end{aligned}
\end{equation}
From \eqref{eq:Hil_ser_modt4} and \eqref{eq:I-3_ser} we see that 
\begin{equation}
    \CI_3=\left.\mr{Ser}_{q=0,t=0}\frac{1-t^2(1+4q + 5\sum_{\ell=2}^\infty q^\ell)+t^3(4q + 8 q^2 + 11 q^3 + 12 \sum_{\ell\geq 4}q^\ell)}{(1-t)^3 \prod_{i=1}^{\infty}\left(1-q^i t\right)^4}\right|_{t^3}~.
\end{equation}
Noting that 
\begin{equation}
\begin{split}
    1-t^2 \sum_{\ell=0}^{\infty}q^\ell -t^3\left(3 q^2 +4\sum_{\ell=3}^\infty q^\ell\right)&=t^2 q \mathcal{I}_1 - t^3(3q \mathcal{I}_1 +q^3 \mathcal{I}_1-5 q -q^2)+1-t^2\left(1+4q + 5\sum_{\ell=2}^\infty q^\ell\right)\\&+t^3\left(4q + 8 q^2 + 11 q^3 + 12 \sum_{\ell\geq 4}q^\ell\right)~,
\end{split}    
\end{equation}
we get 
\begin{equation}
\begin{split}
 H_{R^{(A_1, D_5)}_{\infty}}(q, t)|_{t^3} & = \mathcal{I}_3 +\mr{Ser}_{q=0, t=0} \frac{t^2 q \mathcal{I}_1 - t^3(3q \mathcal{I}_1 +q^3 \mathcal{I}_1-5 q -q^2)}{(1-t)^3 \prod_{i=1}^{\infty}\left(1-q^i t\right)^4}\bigg|_{t^3}
 \\
 &=\mathcal{I}_3 + q \mathcal{I}_1\mr{Ser}_{q=0, t=0} \frac{1}{(1-t)^3 \prod_{i=1}^{\infty}\left(1-q^i t\right)^4}\bigg|_{t} -\mr{Ser}_{q=0, t=0} \frac{t^3(3q \mathcal{I}_1 +q^3 \mathcal{I}_1-5 q -q^2)}{(1-t)^3 \prod_{i=1}^{\infty}\left(1-q^i t\right)^4}\bigg|_{t^3}
 \\
& = \mathcal{I}_3 + q \mathcal{I}_1^2 -3q \mathcal{I}_1 -q^3 \mathcal{I}_1+5 q +q^2~.   \end{split}    
\end{equation}
We can express $\mathcal{I}_1, \mathcal{I}_1^2$ by using $\mathcal{I}_2$ by
\begin{equation}
\mathcal{I}_1 = \frac{(1-q^2) \mathcal{I}_2 - 5+ q}{3 q + q^3}~, \quad 
\mathcal{I}_1^2 = \frac{1}{q}\left((1+q)\mathcal{I}_2 -5 - 4q - q^2\right)~. 
\end{equation}
Then
\begin{equation}
H_{R^{(A_1, D_5)}_{\infty}}(q, t)|_{t^3} = \mathcal{I}_3 + (q + q^2)\mathcal{I}_2~.
\end{equation}
The calculation above suggests the following form for the MacDonald index:
\begin{equation}\label{eq:Mac_A1D5_conj}
    \CI_{\mr{M}}^{(A_1,D_5)}(q,t)=1+\sum_{m=1}^\infty 
    q^m \sum_{l=1}^m
    \frac{f^2_{m,l}(q)}{(q;q)_l}\sum_{k=0}^{2l}\binom{2l}{k}_q~,
\end{equation}
where $f^2_{m,l}(q)$ is a polynomial which we found for some special values of $m,l$ above:
\begin{equation}
\begin{split}
    f^2_{1,1}&=1,
    \\
    f^2_{2,1}&=q, \quad f^2_{2,2}=1~,
    \\
    f^2_{3,1}&=0, \quad f^2_{3,2}=q+q^2, \quad f^2_{3,3}=1~. 
\end{split}    
\end{equation}
Using Mathematica computations, and matching the results to the MacDonald index \eqref{eq:Mac_class_S-1} -- \eqref{eq:Mac_class_S-7} obtained from class $\CS$ description, we were able to compute the functions $f^2_{m,l}(q)$ for high values of $m,l$. The result is recorded in Appendix \ref{app:data_fmnl} and are consistent with the closed formula \eqref{eq:fnml_closed_form}.
\par
%In the next section, we will find a closed formula for $f_{m,l}^n(q)$ by requiring that \eqref{eq:Mac_A1D5_conj} satisfies the Schur limit. 

Numerical evidence suggests the general form of MacDonald index for $(A_1,D_{2n+1})$ theory to be
\begin{equation}\label{eq:Mac_A1D2n+1_conj}
\CI_{\mr{M}}^{(A_1,D_{2n+1})}(q,t)=1+\sum_{m=1}^\infty 
    t^m \sum_{l=1}^m
    \frac{f^n_{m,l}(q)}{(q;q)_l}\sum_{k=0}^{2l}\binom{2l}{k}_q~,    
\end{equation}
for some polynomials $f^n_{m,l}(q)$. Again, matching with the MacDonald index \eqref{eq:Mac_class_S-1} -- \eqref{eq:Mac_class_S-7}, we were able to compute the polynomials $f_{m,l}^n(q)$, the results are recorded in Appendix \ref{app:data_fmnl} for $n=2,3$. 
\par 
In addition, we checked on Mathematica that the Hilbert series of the arc space of the variety 
\begin{equation}\label{eq:variety_gen_n}
    R:=\frac{\C[X,Y,Z,J]}{\langle g_1,g_2,g_3,g_4,g_5\rangle}~,
\end{equation}
where 
\begin{equation}\label{eq:const_gen_n}
    g_1:=XZ-Y^2-J\quad g_2:=J^{n+1}X,\quad g_3:=J^{n+1}Y,\quad g_4:=J^{n+1}Z,\quad g_5:=J^{n+1}~,
\end{equation}
reproduces the MacDonald index \eqref{eq:Mac_ind_gen_n} to high orders in $t$. This suggests the following conjecture:
\begin{conj}
The Zhu algebra of $\mf{su}(2)_{-\frac{4n}{2n+1}}$ is given by 
\begin{equation}
    A\left(\mathfrak{su}(2)_{-\frac{4n}{2n+1}}\right)\cong R~,
\end{equation}
with $R$ given by \eqref{eq:variety_gen_n}.
\end{conj}

\subsection{Matching The Schur Limit}\label{sec:schur_gen_n}
%The discussion above again suggests that the MacDonald index for $(A_1,D_{2n+1})$ theory can be obtained from the MacDonald index for the projected free hypermultiplet by appropriate modifications. Our generalization is guided by the matching of the Schur limit. 
Matching the Schur limit of the formula \eqref{eq:Mac_A1D2n+1_conj} amounts to proving the following product-sum identity:
\begin{equation}\label{eq:new_id_gen_n}
   \left(\frac{\left(q^{2n+1} ; q^{2n+1}\right)_{\infty}}{(q ; q)_{\infty}}\right)^3= 1 + \sum_{m=1}^\infty 
    q^m \sum_{l=1}^m
    \frac{f^n_{m,l}(q)}{(q;q)_l}\sum_{k=0}^{2l}\binom{2l}{k}_q~.
\end{equation}
Closely examining the proof of Theorem \ref{thm:main_id}, we notice that \eqref{eq:new_id_gen_n} holds if
$f^n_{m,l}$ satisfies the hypergeometric relation:
\begin{equation}\label{eq:fmnl_def}
    1 + \sum_{m=1}^\infty q^m \sum_{l=1}^m 
    \frac{f^n_{m,l}(q)}{(q;q)_l} (q^N; q)_l 
    (q^{1-N};q)_l = q^{nN(N-1)}~.
\end{equation}
Indeed, we can prove the more general identity 
\begin{equation}\label{eq:main_id_gen}
\frac{\left(z ; q^{2n+1}\right)_{\infty}\left(q^{2n+1} / z ; q^{2n+1}\right)_{\infty}\left(q^{2n+1} ; q^{2n+1}\right)_{\infty}}{(z ; q)_{\infty}(q / z ; q)_{\infty}(q ; q)_{\infty}}=\sum_{m=0}^{\infty}q^m\sum_{l=1}^m \frac{f^n_{m,l}(q)}{(q ; q)_l} \sum_{k=-l}^l{
2 l \choose
l+k}_q z^k~.    
\end{equation}
The idea is to again compare the powers of $z^N$ on both sides of \eqref{eq:main_id_gen}:
\begin{equation}\label{eq:genid_intermediate}
\sum_{N=-\infty}^{\infty}(-1)^N q^{(2n+1)\binom{N}{2}} z^N=\sum_{r=-\infty}^{\infty}(-1)^r q^{\binom{r}{2}} z^r \sum_{m=0}^{\infty}q^m \sum_{l=1}^m \frac{f^n_{m,l}(q)}{(q ; q)_l} \sum_{k=-l}^l{
2 l \choose
l+k}_q z^k~.     
\end{equation}
Again, we will only show that coefficients of $z^N$ match on both sides for $N\geq 0$ and use the $z\to z^{-1}$ symmetry on both sides for $N<0$. The coefficient of $z^N$ in the \tsf{LHS} of \eqref{eq:genid_intermediate} is 
\begin{equation}
    (-1)^Nq^{(2n+1){N\choose 2}}
\end{equation}
The coefficient of $z^N$ on the \tsf{RHS} of \eqref{eq:genid_intermediate} is
\begin{equation}\label{eq:RHS_zN_coeff}
\begin{aligned}
& \sum_{m=0}^{\infty}q^m\sum_{l=1}^m \frac{f^n_{m,l}(q)}{(q ; q)_l} \sum_{k=-l}^l{
2 l \choose
l+k}_q(-1)^{N-k} q^{\binom{N-k}{2}} \\
& \quad=(-1)^N q^{\binom{N}{2}} \sum_{m=0}^{\infty} \sum_{l=1}^m \frac{f^n_{m,l}(q)}{(q ; q)_l} \sum_{k=-l}^l{
2 l \choose
l+k}_q(-1)^k q^{\binom{k}{2}+k(1-N)} \\
& \quad=(-1)^N q^{\binom{N}{2}} \sum_{m=0}^\infty q^m \sum_{l=1}^m 
    \frac{f^n_{m,l}(q)}{(q;q)_l} (q^N; q)_l 
    (q^{1-N};q)_l ~.
\end{aligned}
\end{equation}
Now using \eqref{eq:fmnl_def} we get 
\begin{equation}
\sum_{m=0}^{\infty}q^m\sum_{l=1}^m \frac{f^n_{m,l}(q)}{(q ; q)_l} \sum_{k=-l}^l{
2 l \choose
l+k}_q(-1)^{N-k} q^{\binom{N-k}{2}} =(-1)^N q^{\binom{N}{2}}q^{nN(N-1)}=(-1)^Nq^{(2n+1){N\choose 2}}~.    
\end{equation}
This proves the identity \eqref{eq:main_id_gen}. Taking $z\to 1$ in \eqref{eq:main_id_gen} gives \eqref{eq:new_id_gen_n}. 
\subsection{Consistency With \tsf{RG} Flow}\label{sec:prod_A2n_fH_Mac}
From the discussion in Section \ref{sec:RG_flow}, since $(A_{1},D_{2n+1})$ flows to free hypermultiplet and $(A_1,A_{2(n-1)})$ theory, we expect that the MacDonald index satisfy
\begin{equation}
    \CI^{(A_1,D_{2n+1})}_{\mr{M}}\sim \CI^{(A_1,A_{2(n-1)})}_{\mr{M}}\cdot \CI^{\mr{HM}/\Z_2}_{\mr{M}}~.
\end{equation}
As an example, we compute the coefficients of the product of MacDonald indices of $(A_1,A_2)$ theory and the projected free hyper in Table \ref{tab:coeff_mac_A1A2_fH}. Comparing it with the coefficients of the MacDonald index for $(A_1,D_5)$ theory in Table \ref{tab:coeff_mac_A1D5}, we find that the coefficients above the diagonal match. 
\par
We now show that the product of the MacDonald indices of $(A_1,A_{2(n-1)})$ theory and the projected free hypermultiplet matches the MacDonald index of $(A_1,D_{2n+1})$ up to a correction factor. We begin by showing this for $n=2$ where the analysis is cleaner. 
\begin{table}[]
    \centering
    \begin{tabular}{|l|lllllllll|}
\hline Variables & $t^0$ & $t^1$ & $t^2$ & $t^3$ & $t^4$ & $t^5$ & $t^6$ & $t^7$ & $t^8$ \\
\hline
$q^0$ & 1 & 3 & 5 & 7 & 9 & 11 & 13 & 15 & 17 \\
$q^1$ & 0 & 5 & 11 & 17 & 23 & 29 & 35 & 41 & 47 \\
$q^2$ & 0 & 8 & 24 & 40 & 56 & 72 & 88 & 104 & 120 \\
$q^3$ & 0 & 9 & 42 & 78 & 114 & 150 & 186 & 222 & 258 \\
$q^4$ & 0 & 12 & 69 & 147 & 225 & 303 & 381 & 459 & 537 \\
$q^5$ & 0 & 13 & 98 & 247 & 401 & 555 & 709 & 863 & 1017 \\
$q^6$ & 0 & 16 & 141 & 402 & 698 & 994 & 1290 & 1586 & 1882 \\
\hline
\end{tabular}
    \caption{Table of selected coefficients of the Macdonald index for $(A_1,A_2)\otimes\mr{HM}/\mathds{Z}_2$ theory. The coefficients above the diagonal matches with the index for $(A_1,D_5)$ in Table \ref{tab:coeff_mac_A1D5}.}
    \label{tab:coeff_mac_A1A2_fH}
\end{table}
\subsubsection{$(A_1,D_5)\to(A_1,A_2)\otimes\mr{HM}/\mathds{Z}_2$}
The MacDonald index of the $(A_1,A_2)$ theory is given by \cite{Bhargava:2023hsc}\footnote{Note again that because of our convention, we do no have a factor of $q^n$, compared to \cite{Bhargava:2023hsc}.}:
\begin{equation}
    \CI_{\mr{M}}^{(A_1,A_2)}(q,t)=\sum_{n=0}^{\infty}t^n\frac{q^{n^2}}{(q;q)_n}~.
\end{equation}
The MacDonald index of the $\mathds{Z}_2$-projected free hypermultiplet is,
\begin{equation}
    \CI_{\mr{M}}^{\mr{HM}/\mathds{Z}_2}(q,t)=\sum_{l=0}^{\infty}\frac{t^l}{(q;q)_{2l}}\sum_{k=0}^{2l}{2l\choose k}_q
\end{equation}
Taking a product of these two indices, we obtain,
\begin{equation}
    \sum_{n,l=0}^{\infty}t^{n+l}\frac{q^{n^2}}{(q;q)_{2l}(q;q)_n}\sum_{k=0}^{2l}{2l\choose k}_q
\end{equation}
Substituting $m=n+l$,
\begin{equation}
    \sum_{m=0}^{\infty}t^m\sum_{l=0}^m\frac{q^{(m-l)^2}}{(q;q)_{m-l}}\frac{\textcolor{red}{1}}{\textcolor{red}{(q;q)_{2l}}}\sum_{k=0}^{2l}{2l\choose k}_q
\end{equation}
Now, we turn to the $(A_1,D_5)$ index. Using the fact that 
\begin{equation}
    f^2_{m,l}(q)=q^{(m-l)^2}{l\choose m-l}_q=q^{(m-l)^2}\frac{(q;q)_{l}}{(q;q)_{m-l}(q;q)_{2l-m}}~,
\end{equation}
we have 
\begin{equation}
\begin{split}
    \CI_{\mr{M}}^{(A_1,D5)}(q,t)&=\sum_{m=0}^{\infty}t^m\sum_{l=0}^{m}\frac{q^{(m-l)^2}}{(q;q)_l}\frac{(q;q)_{l}}{(q;q)_{m-l}(q;q)_{2l-m}}\sum_{k=0}^{2l}{2l\choose k}_q
    \\
    &=\sum_{m=0}^{\infty}t^m\sum_{l=0}^{m}\frac{q^{(m-l)^2}}{(q;q)_{m-l}}\frac{1}{\textcolor{red}{(q;q)_{2l-m}}}\sum_{k=0}^{2l}{2l\choose k}_q
\end{split}
\end{equation}
The correction factor is given by 
\begin{equation}
    \textcolor{red}{\frac{(q;q)_{2l}}{(q;q)_{2l-m}}}~.
\end{equation}
\subsubsection{The General $n$ Case}
We now show this for general $n$. The MacDonald index of the $(A_1,A_{2(n-1)})$ theory is given by \cite{Buican:2015tda}
\begin{equation}
    \CI_{\mr{M}}^{(A_1,A_{2(n-1)})}(q,t)=\sum_{N_1\geq N_2\geq\dots\geq N_{n-1}\geq 0}t^{N_1+N_2+\dots+N_{n-1}}\frac{q^{N_1^2+\dots+N_{n-1}^2}}{(q;q)_{N_1-N_2}\dots (q;q)_{N_{n-2}-N_{n-1}}(q;q)_{N_{n-1}}}~.
\end{equation}
The MacDonald index of the $\mathds{Z}_2$-projected free hypermultiplet is,
\begin{equation}
    \CI_{\mr{M}}^{\mr{HM}_2/\mathds{Z}_2}(q,t)=\sum_{l=0}^{\infty}\frac{t^l}{(q;q)_{2l}}\sum_{k=0}^{2l}{2l\choose k}_q~.
\end{equation}
Taking a product of these two indices, we obtain,
\begin{equation}\label{eq:prod_A1D2n+1HMZ2}
    \sum_{\substack{N_1\geq N_2\geq\dots\geq N_{n-1}\geq 0\\l\geq 0}}t^{N_1+\dots+N_{n-1}+l}\frac{q^{N_1^2+\dots+N_{n-1}^2}}{(q;q)_{N_1-N_2}\dots (q;q)_{N_{n-2}-N_{n-1}}(q;q)_{N_{n-1}}(q;q)_{2l}}\sum_{k=0}^{2l}{2l\choose k}_q~.
\end{equation} 
We claim that the correction factor is 
\begin{equation}
    \textcolor{red}{q^{(n-1)^2N_{n-1}^2-(n-1)N_{n-1}^2-2N_{n-1}(N_1+\dots+N_{n-2})}\frac{(q;q)_{N_{n-1}}(q;q)_{2l}}{(q;q)_{(n-2)N_{n-1}-N_1}(q;q)_{l-(n-1)N_{n-1}}}}~.
\end{equation}
To prove this, let us insert the correction factor in the expression \eqref{eq:prod_A1D2n+1HMZ2}. We get 
\begin{equation}
\begin{split}
\sum_{\substack{N_1\geq N_2\geq\dots\geq N_{n-1}\geq 0\\l\geq 0}}t^{N_1+\dots+N_{n-1}+l}\frac{q^{N_1^2+\dots+N_{n-1}^2+(n-1)^2N_{n-1}^2-(n-1)N_{n-1}-2N_{n-1}(N_1+\dots+N_{n-2})}}{(q;q)_{N_1-N_2}\dots (q;q)_{N_{n-2}-N_{n-1}}(q;q)_{(n-2)N_{n-1}-N_1}(q;q)_{l-(n-1)N_{n-1}}}\\\times\sum_{k=0}^{2l}{2l\choose k}_q~. 
\end{split}
\end{equation}
Note that 
\begin{equation}
\begin{split}
    N_1^2+\dots+N_{n-1}^2+(n-1)^2N_{n-1}^2-(n-1)N_{n-1}^2-2N_{n-1}(N_1+\dots+N_{n-2})\\=(N_1-N_{n-1})^2+\dots+(N_{n-2}-N_{n-1})^2+(n-1)^2N_{n-1}^2~.
\end{split}
\end{equation}
Put 
\begin{equation}
    m:=N_1+\dots+N_{n-1}+l,\quad i_{j+2}:=N_{j}-N_{j+1}~,\quad j=1,\dots,n-2~.
\end{equation}
Then we have 
\begin{equation}
    N_1+\dots N_{n-2}-(n-2)N_{n-1}=i_3+2i_4+\dots+(n-2)i_{n}~.
\end{equation}
Using this, we have 
\begin{equation}
    (n-1)N_{n-1}=m-i_3-2i_4-\dots-(n-2)i_{n}~.
\end{equation}
Using the fact that 
\begin{equation}
    N_1-N_{n-1}=i_3+\dots+i_n~,
\end{equation}
we get 
\begin{equation}
    (n-2)N_{n-1}-N_1=m-2i_3-3i_4-\dots-(n-1)i_n~.
\end{equation}
Using all these relations, the sum \eqref{eq:prod_A1D2n+1HMZ2} becomes 
\begin{equation}
\begin{split}
\sum_{m=0}^\infty t^m\sum_{l=0}^m\frac{1}{(q;q)_l}\sum_{\substack{i_3, i_4,\dots,i_n\geq 0}}&\frac{q^{Q(m-(n-1)i_n-(n-2)i_{n-1}-\dots-l,i_3,\dots,i_n)}(q;q)_{l}}{(q;q)_{i_n}\dots (q;q)_{i_3}(q;q)_{m-(n-1)i_n-(n-2)i_{n-1}-\dots-l}(q;q)_{2l+i_3+2i_4+\dots+(n-2)i_n-m}}\\&\times\sum_{k=0}^{2l}{2l\choose k}_q
\\&=\sum_{m=0}^\infty t^m\sum_{l=0}^m\frac{f^n_{m,l}(q)}{(q;q)_l}\sum_{k=0}^{2l}{2l\choose k}_q
\\&=\mathcal{I}^{(A_1,D_{2n+1})}_{\mr{M}}(q,t)~,
\end{split}
\end{equation}
where 
\begin{equation}
    Q(i_2,i_3,\dots,i_k):=i_k^2+(i_k+i_{k-1})^2+\dots+(i_k+\dots+i_3+i_2)^2~,
\end{equation}
as defined in \eqref{eq:Q_def}.
\\\\
\textbf{Acknowledgments.} We thank Matthew Buican for collaboration at the initial stages of the project. A.B. would like to thank Philip Argyres for helpful discussions. R.K.S. would like to thank Ken Ono for some useful correspondence and Anirudh Deb and Spencer Stubbs for discussions. R.T. would like to thank Heeyeon Kim for pointing out some interesting directions for future work. R.K.S. would like to thank ICTS, Bangalore for hospitality where this work was completed. The work of A.B. is supported by the US Department of Energy under grant DE-SC1019775. The work of C.B. was partially supported by funds from the Queen Mary University of London. The work of R.K.S. and R.T. is supported by the US
Department of Energy under grant DE-SC0010008. 
%This motivates the following conjecture.   
% In Appendix \ref{app:new_id}, we prove the identity 
% \begin{equation}\label{eq:new_id_gen_n}
%    \left(\frac{\left(q^{2n+1} ; q^{2n+1}\right)_{\infty}}{(q ; q)_{\infty}}\right)^3= 1 + \sum_{m=1}^\infty 
%     q^m \sum_{l=1}^m
%     \frac{f^n_{m,l}(q)}{(q;q)_l}\sum_{k=0}^{2l}\binom{2l}{k}_q~.
% \end{equation}
%Compared with the Schur index \eqref{eq:schur_D2n+1_pochh_form}, we see that our formula \eqref{eq:Mac_ind_gen_n} for the MacDonald index satisfies the correct Schur limit. 
\begin{comment}
\subsection{$(A_1, D_5)$}
We observe the following recursion relation for $f^{n=2}_{m,l}(q)$
\begin{equation}
    \begin{aligned}
& f^2_{m, m}=1, \quad m > 0 \\
& f^2_{2 m, m}=q^{m^2}, \quad m \geq 1 \\
& f^2_{m, l}=0, \quad m>2 l\\
&f^2_{m,m-k} = \sum_{l=2k-2}^{m-2}
q^{l+1} f_{l,l-k+1},  \quad m\geq 2k, \quad 1\leq k < m/2 
\end{aligned}
\end{equation}
Equivalently, this is 
\begin{equation}
    f^{n=2}_{m,l} = \begin{cases}
        0, \quad m \leq 0\; \text{or}\; l \leq 0\;, \text{or}\; m > 2l,\\
        1,\quad m=l>0, \\
        q^{l^2}, \quad m = 2l\\
        q^{m-1} f^{n=3}_{m-2,l-1} + f^{n=3}_{m-1,l-1} , \quad m/2< l < m,
    \end{cases}
\end{equation} 
%
\subsection{$(A_1, D_7)$}
The following recursion relation for 
$f^{n=3}_{m,l}(q)$ is conjectured.
\begin{equation}
    f^{n=3}_{m,l} = \begin{cases}
        0, \quad m \leq 0\; \text{or}\; l \leq 0\;, \text{or}\; m > 3l,\\
        1,\quad m=l> 0, \\
        q^{2 l^2}, \quad m = 3l\\
        q^l f^{n=2}_{m-1,l} + f^{n=2}_{m-1,l-1}-q^{2m-l-4} f^{n=2}_{m-4,l-1}, \quad m/3< l < m,
    \end{cases}
\end{equation}
\end{comment}

\begin{appendix}
\section{Some Special Functions}\label{app:special_func}
We make use of several special functions in defining and working with the superconformal index and its limits. Here we compile a list of functions used in this work. 
\paragraph{Plethystic Exponential.} This function is defined as:
\begin{equation}
    \mr{PE}[f(t)]=\exp\bigg(\sum_{n=1}^{\infty}\frac{f(t^n)}{n}\bigg)~.
\end{equation}
For example
\begin{equation}
    \mr{PE}[t]=\exp\bigg(\sum_{n=1}^{\infty}\frac{t^n}{n}\bigg)=\exp\bigg(-\log(1-t)\bigg)=\frac{1}{1-t}~.
\end{equation}
Following are two other properties of this function that will be useful for us:
\begin{equation}
    \begin{split}
        &\mr{PE}[f + g]=\mr{PE}[f]\mr{PE}[g]~, \\
        &\mr{PE}[-f]=1/\mr{PE}[f]~. \\
    \end{split}
\end{equation}
\paragraph{Elliptic Gamma Function.} This function is defined as:
\begin{equation}
    \Gamma_e(t):=\Gamma_e(t;p,q)=\prod_{m,n=0}^{\infty}\frac{1-p^{m+1}q^{n+1}t^{-1}}{1-p^mq^nt}=\mr{PE}\bigg[\frac{t-
    \frac{pq}{t}}{(1-p)(1-q)}\bigg]~.
\end{equation}
\paragraph{$q$-Pochhammer Symbol.} These are $q$-analogues of the Pochhammer symbols. This function is defined as:
\begin{equation}
    (z;q)_n=\prod_{k=0}^{n-1}(1-zq^k)~.
\end{equation}
We can similarly define the infinite $q$-Pochhammer symbol:
\begin{equation}
    (z;q)_{\infty}=\prod_{k=0}^{\infty}(1-zq^k)~.
\end{equation}
We also define 
\begin{equation}
    (a_1,a_2,\dots,a_n;q)_\infty:=(a_1;q)_\infty(a_2;q)_\infty\cdots (a_n;q)_\infty~.
\end{equation}
Some results about $q$-Pochhamer symbols that we use are recorded below, see \cite{gasper2004basic,andrews1998theory} for proofs.
\begin{align}
    (a ; q)_n &=\frac{(a ; q)_{\infty}}{\left(a q^n ; q\right)_{\infty}}~.
    \\
    (a ; q)_{n+k}&=(a ; q)_n\left(a q^n ; q\right)_k~.
    \\
    \left(a q^k ; q\right)_{n-k}&=\frac{(a ; q)_n}{(a ; q)_k}~.
    \\
    \frac{1}{(z;q)_\infty}&=\sum_{m\geq 0}\frac{z^m}{(q;q)_m}~.
    \\
    \sum_{n \in \mathds{Z}} q^{\frac{n(n+1)}{2}} z^n&=(q ; q)_{\infty}\left(-1/z ; q\right)_{\infty}(-z q ; q)_{\infty}~.
\end{align}

\paragraph{Functions involving $q$-analogues.}
We first define the $q$-analogues of integers:
\begin{equation}
    [n]_q:=\frac{1-q^n}{1-q}=1+q+q^2+\cdots +q^{n-1}, \quad \lim_{q\to 1}[n]_q=n~.
\end{equation}
We can then use this to define the $q$-analogues of factorials:
\begin{equation}
    [n]_q!:=\prod_{k=1}^{n}\frac{(1-q^k)}{(1-q)}=\frac{(q;q)_n}{(1-q)^n}
\end{equation}
We can use this to define $q$-binomial coefficients as follows:
\begin{equation}
    {n\choose k}_q:=\frac{[n]_q!}{[k]_q![n-k]_q!}=\frac{(q;q)_n}{(q;q)_k(q;q)_{n-k}}~.
\end{equation}
\paragraph{Basic Hypergeometric Series.}
The original Gaussian hypergeometric series $_2F_1(a,b;c;z)$ is given as:
\begin{equation}
    _2F_1(a,b;c;z)=\sum_{n=0}^{\infty}\frac{z^n}{n!}\frac{\Gamma(a)\Gamma(b)}{\Gamma(c)}~,
\end{equation}
and is well-studied in the context of the hypergeometric differential equation and Sturm-Liouville problems. 
This can be generalized to involve $q$-analogues of the integers as defined above: for $r>s+1,q\neq 0$,
\begin{equation}\label{eq:phirs_def}
\begin{aligned}  { }_r \phi_s\left(a_1, a_2, \ldots, a_r ; b_1, \ldots, b_s ; q, z\right) &\equiv{ }_r \phi_s\left[\begin{array}{c}a_1, a_2, \ldots, a_r \\ b_1, \ldots, b_s\end{array} ; q, z\right] \\ & :=\sum_{n=0}^{\infty} \frac{\left(a_1 ; q\right)_n\left(a_2 ; q\right)_n \cdots\left(a_r ; q\right)_n}{(q ; q)_n\left(b_1 ; q\right)_n \cdots\left(b_s ; q\right)_n}\left[(-1)^n q^{\binom{n}{2}}\right]^{1+s-r} z^n ~, 
\end{aligned}
\end{equation}
We will use the following relations \cite{gasper2004basic}:
\begin{enumerate}
\item The $q$-Vandermonde ($q$-Chu-Vandermonde) sums:
\begin{equation}
{ }_2 \phi_1\left(a, q^{-n} ; c ; q, q\right)=\frac{(c / a ; q)_n}{(c ; q)_n} a^n~.
\end{equation}
\item Reversing the order of summation,
\begin{equation}
{ }_2 \phi_1\left(a, q^{-n} ; c ; q, c q^n / a\right)=\frac{(c / a ; q)_n}{(c ; q)_n}~ .
\end{equation}
\item Heine's transform:
\begin{equation}
    { }_2 \phi_1(a, b ; c ; q, z)=\frac{(b, a z ; q)_{\infty}}{(c, z ; q)_{\infty}}{ }_2 \phi_1(c / b, z ; a z ; q, b)~.
\end{equation}
\item Expanding the \tsf{RHS} of Heine's transform gives:
\begin{equation}\label{eq:Heine_exp}
    { }_2 \phi_1(a, b ; c ; q, z)=\frac{(a z ; q)_{\infty}}{(z ; q)_{\infty}} \sum_{n=0}^{\infty} \frac{(a, c / b ; q)_n}{(q, c, a z ; q)_n}(-b z)^n q^{\binom{n}{2}}~.
\end{equation}
\end{enumerate}
% A useful special case is the following,
% \begin{equation}
%     _{k+1}\phi_k(a_1,\cdots,a_j;b_1,\cdots,b_k;q,z)=\sum_{n=0}^{\infty}t^n\frac{(a_1;q)_n\cdots (a_k;q)_n(a_{k+1};q)_n}{(b_1;q)_n\cdots (b_k;q)_n(q;q)_n}
% \end{equation}
\section{Conventions For Superconformal Index And Various Limits}\label{app:conventions}
There exist multiple conventions for the superconformal index in the literature, so in this appendix, we lay out our choice of conventions. See Table 1 of \cite{Gadde:2011uv} for other choices. \\
The superconformal index evaluates the partition function of a 4d $\mathcal{N}=2$ \tsf{SCFT} on $S^3\times S^1_{\beta}$ with certain fugacities for background fields turned on. In the Hamiltonian perspective, this involves four commuting supercharges and takes the form,
\begin{equation}
    \CI_{\mr{SCI}}(\rho,\sigma,\tau)=\text{Tr}_{\mathcal{H}(S^3)}[(-1)^F\sigma^{\frac{1}{2}\tilde{\delta}^1_{\dot{-}}}\rho^{\frac{1}{2}\tilde{\delta}^1_{\dot{+}}}\tau^{\frac{1}{2}\delta^2_{+}}e^{-\beta\delta^1_{-}}]~,
\end{equation}
where,
\begin{equation}
    \begin{split}
        &\delta^1_{-}=E-2R+r-2j~, \\
        &\tilde{\delta}^1_{\dot{+}}=E-2R-r+2\bar{j} ~,\\
        &\tilde{\delta}^1_{\dot{-}}=E-2R-r-2\bar{j}~, \\
        &\delta^2_{+}=E+2R+r+2j~. \\
    \end{split}
\end{equation}
We can change to a different set of variables as follows,
\begin{equation}
    \begin{split}
        &t:=\tau^2~, \\
        &p:=\sigma\tau~, \\
        &q:=\rho\tau~. \\
    \end{split}
\end{equation}
In these new variables, we have,
\begin{equation}
    \CI_{\mr{SCI}}(p,q,t)=\text{Tr}_{\mathcal{H}(S^3)}[(-1)^{F}p^{\frac{1}{2}\tilde{\delta}^1_{\dot{-}}}q^{\frac{1}{2}\tilde{\delta}^1_{\dot{+}}}t^{\frac{1}{4}(\delta^2_+-\tilde{\delta}^1_{\dot{+}}-\tilde{\delta}^1_{\dot{-}})}e^{-\beta\delta^1_-}]~.
\end{equation}
The index is automatically independent of $\beta$, i.e., it counts $1/8$-BPS operators that satisfy,
\begin{equation}
    E=2R-r+2j~.
\end{equation}
Manually substituting this condition into the index, we can write it in a simplified form,
\begin{equation}
    \CI_{\mr{SCI}}(p,q,t)=\text{Tr}_{\mathcal{H}(S^3)}[(-1)^Fp^{j-\bar{j}-r}q^{j+\bar{j}-r}t^{R+r}]~.
\end{equation}
Now, the superconformal index admits various special limits that give more refined counts of BPS operators. We describe some of these limits below, 
\paragraph{MacDonald limit ($p\to 0$).} This implements the additional condition that $j-\bar{j}-r=0$, which leads to the following conditions,
\begin{equation}
    \begin{split}
        &E-2R=j+\bar{j}~, \\
        &r=j-\bar{j}~. \\
    \end{split}
\end{equation}
Only $\hat{\mathcal{B}}_R$, $\bar{\mathcal{D}}_{R(j,0)}$ and $\hat{\mathcal{C}}_{R(j,\bar{j})}$ multiplets\footnote{We are using the notation of \cite{Dolan:2002zh}. See \cite{Cordova:2016emh} for another popular choice.} contribute to the MacDonald index $\CI_{\mr{M}}(q,t)$. We call this the Schur ring of operators, for reasons discussed next.

\paragraph{Schur limit ($t\to q$).} In this limit, the index automatically becomes independent of $p$, so it counts operators satisfying $E=2R-r+2j$ and $r=j-\bar{j}$. Written differently, the constraints are,
\begin{equation}
    \begin{split}
        &E-2R=j+\bar{j}~, \\
        &r=j-\bar{j}~. \\
    \end{split}
\end{equation}
so we see that we count the same operators as we did in the MacDonald limit, but the count is less refined. The contribution of these operators simplifies to,
\begin{equation}
    \CI_{\mr{S}}(q)=\text{Tr}_{\mathcal{H}_{\text{Schur}}}[(-1)^Fq^{E-R}]~.
\end{equation}
Similarly, the MacDonald index may be expressed as,
\begin{equation}
    \CI_{\mr{M}}(q,t)=\text{Tr}_{\mathcal{H}_{\text{Schur}}}[(-1)^Fq^{E-2R-r}t^{R+r}]=\text{Tr}_{\mathcal{H}_{\text{Schur}}}[(-1)^Fq^{2\bar{j}}t^{R+j-\bar{j}}]~.
\end{equation}
This last form is particularly convenient in studying the $(A_1,D_{2n+3})$ and $(A_1,A_{2n})$ theories, where all the Schur operators have $j=\bar{j}$, so that $r=0$. 
\paragraph{Hall-Littlewood limit ($p\to 0$, $q\to 0$).} This limit implements the additional conditions $j-\bar{j}-r=0$ and $j+\bar{j}-r=0$, which leads to the following conditions on quantum numbers,
\begin{equation}
    \begin{split}
        &E=2R+r~, \\
        &j=r~, \\
        &\bar{j}=0~. \\
    \end{split}
\end{equation}
Only $\hat{\mathcal{B}}_R$ and $\bar{\mathcal{D}}_{R(j,0)}$ contribute to this limit. 
%For example, one component of the gaugino (a level one descendant in a $\bar{\mathcal{D}}_{0(0,0)}$ multiplet) $\lambda^1_{+}$ with $R=r=j=1/2$, $\bar{j}=0$ and $E=3/2$ contributes to the Hall-Littlewood index.%
The contribution of these operators may be expressed as,
\begin{equation}
    \CI_{\mr{HL}}(t)=\text{Tr}_{\mathcal{H}_{\mr{HL}}}[(-1)^Ft^{E-R}]~.
\end{equation}
% \paragraph{Coulomb limit ($p,q,t\to 0$ with $u:=pq/t$ held constant).} The operators that contribute to this index live in $1/2$-BPS miltiplets annihilated by $Q^I_{\alpha}$ for $I\in\{1,2\}$. Their quantum numbers therefore satisfy the following constraints,
% \begin{equation}
%     \begin{split}
%         &R=0 \\
%         &j=0 \\
%         &E=-r \\
%     \end{split}
% \end{equation}
% The only operators that contribute in this limit are the superconformal primaries of $\mathcal{E}_{r(0,\bar{j})}$ multiplets. These multiplets form the $\mathcal{N}=2$ chiral ring.

\paragraph{Higgs branch Hilbert series.} This quantity is related to the Hall-Littlewood index, although it is not directly a limit of the $\mathcal{N}=2$ superconformal index. The operators that contribute to this quantity are $1/2$-BPS operators annihilated by $Q^1_{\alpha}$ and $\tilde{Q}^1_{\dot{\alpha}}$. This means that their quantum numbers satisfy,
\begin{equation}
    \begin{split}
        &j=\bar{j}=0~, \\
        &r=0~, \\
        &E=2R~.
    \end{split}
\end{equation}
Only the superconformal primaries of $\hat{\mathcal{B}}_R$ multiplets contribute to the Higgs branch Hilbert series. 
\section{Gr\"{o}bner Basis}\label{app:Grob}
In this appendix, we record some definitions from commutative algebra, specifically the Gr\"{o}bner basis of an ideal of a polynomial ring following \cite{loustaunau1994introduction}. This terminology and some standard results recorded here are used in Section \ref{sec:comp_Hilb_ser}.

Let $\mathds{F}$ be a field and $P:=\IF[x_1,\dots,x_n]$ be the polynomial ring in $n$-variables with coefficients in $\IF$. 
Recall that given $s$ polynomials $f_1,\dots,f_s$, the ideal generated by $I$ is given by 
\begin{equation}
    I=\langle f_1,\dots,f_s\rangle:=\left\{\sum_{i=1}^sg_if_i~:g_i\in P\right\}~.
\end{equation}
Our goal is to get a basis (or at least the dimension) of the vector space $P/I$. This is achieved by finding \textit{the} Gr\"{o}bner basis for $I$ which we now explain. The first step is to choose a \textit{monomial ordering} on the monomials in $P$. A monomial in $P$ is of the form 
\begin{equation}
    X^{\vec{a}}:=x_1^{a^1}\dots x_n^{a^n}~,\quad \vec{a}=(a^1,\dots,a^n)\in (\IZ_{\geq 0})^n~.
\end{equation}
A monomial ordering on $P$ is a total order $<$ on the set of monomials such that  
\begin{enumerate}[(1)]
    \item $X<Y\implies XZ<YZ$ for any monomials $X,Y,Z$.
    \item $1<X$ for any monomial $X$.
\end{enumerate}
There are several standard choices of monomial ordering, lexicographic (lex), degree lexicographic (deglex), degree reverse lexicographic (degrevlex) and so on. We will use the deglex ordering, which we now define.
\paragraph{Lexicographic ordering.} Choose an ordering on the variables. For example:
    \begin{equation}
        x_1>x_2>\dots x_n~.
    \end{equation}
Then to compare two monomials, lex first compares the exponent of $x_1$ (i.e., the largest variables), if equal then compares the exponent of $x_2$ (i.e., the second largest variables) and so forth. For example  \begin{equation}
    x_1^2>x_1 x_2>x_1>x_2^2>x_2>1~.
\end{equation}   
\paragraph{Degree lexicographic ordering.}
Deglex first compares the degree of the monomials. If the degrees are equal then, deglex compares the monomials using lex. For example
\begin{equation}
    x_1^2>x_1 x_2>x_2^2>x_1>x_2>1~.
\end{equation}
Once we have chosen a monomial ordering, we can rearrange any polynomial in the order of decreasing order. 
\begin{defn}
For any $0\neq f\in P$ of degree $r$ given by 
\begin{equation}
    f=\sum_{i=1}^rc_iX^{\vec{a}_i}~,
\end{equation}
such that $X^{\vec{a}_1}>X^{\vec{a}_2}>\dots>X^{\vec{a}_r}$, we define 
\begin{enumerate}[(1)]
    \item $\mathsf{LM}(f)=X^{a_1}$, the leading monomial of $f$,
\item $\mathsf{LC}(f)=c_1$, the leading coefficient of $f$,
\item $\mathsf{LT}(f)=c_1 X^{\alpha_1}$, the leading term of $f$.
\end{enumerate}
\end{defn}
\begin{defn}
\begin{enumerate}[(1)]
    \item Given polynomials $f,g,h\in P$, we say that $f$ is reducible to $h$ by $g$ if $\mr{LM}(g)$ divides a nonzero term $X$ of $f$ and 
    \begin{equation}
        h=f-\frac{X}{\mr{LM}(g)}g~.
    \end{equation}
    We symbolically write the reduction as 
    \begin{equation}
        f\overset{g}{\longrightarrow} h~.
    \end{equation}
    \item Given polynomials $f, h \in R$ and a set $G=\left\{g_1, \ldots, g_s\right\} \subset R$ of non-zero polynomials, we say that $f$ is reducible to $h$ modulo $G$, denoted
\begin{equation}
f \xrightarrow{G} h~,
\end{equation}
if there exists a subset $\{g_{i_1}, \ldots, g_{i_k}\} \subseteq G$ and a sequence of polynomials $h_1, \ldots, h_{k-1}$ such that
\begin{equation}
f \xrightarrow{g_{i_1}} h_1 \xrightarrow{g_{i_2}} h_2 \xrightarrow{g_{i_3}} \cdots \xrightarrow{g_{i_{k-1}}} h_{t-1} \xrightarrow{g_{i_k}} h~ .
\end{equation}
\item A polynomial $f$ is called reduced with respect to $G$ if it cannot be reduced modulo $G$. That is, no term of $f$ is divisible by any one of the $\mr{LM}\left(g_i\right)$.
\item The remainder of a polynomial $f$ with respect to $G$ is a polynomial $h$ such that $f$ is reducible to $h$ modulo $G$ and $h$ is reduced with respect to $G$. 
\end{enumerate}    
\end{defn}
\begin{defn}
A set of non-zero polynomials $G=\left\{g_1, \ldots, g_s\right\}$ is called a Gr\"{o}bner basis for an ideal $I\subset P$ if for all nonzero $f \in I$, $\mr{LM}\left(g_i\right)$ divides $\mr{LM}(f)$ for some $g_i \in G$.
\end{defn} 
\begin{thm}\emph{\cite[Theorem 1.6.7]{loustaunau1994introduction}}\label{prop:remain_grob}
$G$ is a Gr\"{o}bner basis if and only if for all $f \in P$ the remainder of reduction of $f$ by $G$ is unique.    
\end{thm}
Let $r_G(f)\in P$ denote the remainder of the reduction of a polynomial $f$ by $G$. Then Theorem \ref{prop:remain_grob} implies that $f\equiv g\bmod I$ if and only if $g_G(f)=r_G(g)$. This gives us a basis for the $\IF$-vector space $P/I$:
\begin{equation}
    \CB_{P/I}:=\{\text{monomial}~~X^{\vec{a}}:\mr{LM}(g_i)\nmid X^{\vec{a}}\text{ for all }g_i\in G\}~.
\end{equation}
We now discuss an algorithm to construct a Gr\"{o}bner basis for an ideal generated by a set of polynomials. 
\begin{defn}
Given non-zero $f, g \in P$, the $\CS$-polynomial of $f$ and $g$ is defined as
\begin{equation}
\CS(f, g)=\frac{L}{\mathsf{LT}(f)} f-\frac{L}{\mathsf{LT}(g)} g~,
\end{equation}
where $L=\mr{LCM}(\mr{LM}(f), \mr{LM}(g))$.    
\end{defn}
The following theorem is due to Buchberger \cite{Buchberger1985}:
\begin{thm}\emph{\cite[Theorem 1.7.4]{loustaunau1994introduction}}\label{thm:S-pol_grob}
A set of non-zero polynomials $G=\left\{g_1, \ldots, g_s\right\}$ is a Gr\"{o}bner basis for an ideal $I$ if and only if for all $i \neq j$,
\begin{equation}
\CS\left(g_i, g_j\right) \xrightarrow{G} 0~ .
\end{equation}    
\end{thm}
Theorem \ref{thm:S-pol_grob} gives us an algorithm to construct a Gr\"{o}bner basis for $I$. 
\paragraph{Buchberger's algorithm for Gr\"{o}bner basis.} Let $I=\langle f_1,\dots,f_s\rangle\subset P$ where each $f_i\neq 0$. To construct a Gr\"{o}bner basis for $I$, we follow the steps below:
\begin{enumerate}[(1)]
    \item Start with $G:=\{f_1,\dots,f_s\}$. If $s=1$, then $G$ is a Gr\"{o}bner basis. If $s>1$, define $A:=\{(f_i,f_j):i<j\}$.
    \item  For $(f,g)\in A$, calculate $\CS(f,g)$ and reduce modulo $G$: $\CS\left(f, g\right) \xrightarrow{G} r$. If $r\neq 0$ then add $f_{s+1}:=r$ to the set $G$ and the pairs $(f_i,f_{s+1})$ to $A$.
    \item Repeat steps (1) and (2) with the new $G$ and $A$ until $\CS\left(f_i, f_j\right) \xrightarrow{G} 0$ for all $(f_i,f_j)\in A$.
\end{enumerate}
It is a consequence of the Hilbert basis theorem that this algorithm terminates for polynomial rings. We refer the reader to \cite{loustaunau1994introduction} for further details on this point.
\section{Proof Of Formula \eqref{eq: hilbert_ser_Rinfty_gen} For Hilbert Series}\label{app:Hilb_ser_count}
In this appendix, we prove the formula \eqref{eq: hilbert_ser_Rn_gen} and \eqref{eq: hilbert_ser_Rinfty_gen} using explicit counting of monomials.  

Recall that given a Gr\"{o}bner basis of $I_n$, a canonical basis for $R_n/I_n$ is given by 
\begin{equation}
    \CB_{R_n}:=\{\text{monomial}~~X_1^{\vec{a}_1}X_2^{\vec{a}_2}\cdots X_N^{\vec{a}_N}:g\nmid X_1^{\vec{a}_1}X_2^{\vec{a}_2}\cdots X_N^{\vec{a}_N}\text{ for all }g\in \mr{LM}(G_n)\}~.
\end{equation}
This gives us the following algorithm for the computation of $\mr{dim}\left(R_n\right)_{\ell, m}$:
\begin{enumerate}
    \item Count the number of monomials in 
    \begin{equation}
        \C[x^{(1)}_{d_1-1},\dots,x^{(1)}_n,x^{(2)}_{d_2-1},\dots,x^{(2)}_n,\dots,x^{(N)}_{d_N-1},\dots,x^{(N)}_n]
    \end{equation}
    of weight $(\ell,m)$. 
    \item Subtract the number of monomials of weight $(\ell,m)$ that can be obtained by multiplying another monomial (necessarily of weight $(\ell',m')$ with $\ell'<\ell$ and $m'<m$) to any monomial in $G_n$.
\end{enumerate}
Let $M[\ell,m]$ be the number of monomials in variables $\{x^{(1)}_{i_1},\dots,x^{(N)}_{i_N}:i_k\geq d_{k}-1\}$ of weight $(\ell,m)$. We call such monomials \textit{free words}. Then the generating function for $M[\ell,m]$ is given by
\begin{equation}\label{eq:gen_free_words}
    \sum_{\ell,m\geq0}M[\ell,m]q^\ell t^m=\prod_{i=1}^N\prod_{j\geq 0}\left(1-q^{(d_i-1)+j}t\right)^{-1}~.
\end{equation}
Then we can write 
\begin{equation}
\mr{dim}\,(R_n)_{\ell,m}=M[\ell,m]-C[\ell,m]~,    
\end{equation}
where $C[\ell,m]$ is the number of monomials we get from step 2. 
To implement step 2, we count all possible ways of constructing a monomial of weight $(\ell,m)$ by multiplying free words of weight $(\ell',m')~,\ell'<\ell,m'<m$ with monomials in $\mr{LM}(G_n)$ of weight $(\ell-\ell',m-m')$. Let us define 
\begin{equation}
\mr{LM}(G_n)_{\ell,m}:=\{g\in\mr{LM}(G_n):\mr{wt}_1(g)\leq \ell,\mr{wt}_2(g)\leq m\}~.    
\end{equation}
For $g\in \mr{LM}(G_n)_{\ell,m}$, define  
\begin{equation}
C_g[\ell,m]:=\{g\cdot X_1^{\vec{a}_1}X_2^{\vec{a}_2}\cdots X_N^{\vec{a}_N}:\mr{wt}^{(2)}(g\cdot X_1^{\vec{a}_1}X_2^{\vec{a}_2}\cdots X_N^{\vec{a}_N})=(\ell,m) \}~.
\end{equation}
If $\mr{wt}^{(2)}(g)=(\ell',m')$, then clearly 
\begin{equation}
    |C_g[\ell,m]|=M[\ell-\ell',m-m']~.
\end{equation}
Then we can write 
\begin{equation}
    C[\ell,m]=\left|\bigcup_{g\in\mr{LM}(G_n)}C_g[\ell,m]\right|
\end{equation}
Note that 
\begin{equation}
    \bigcap_{i=1}^kC_{g_i}[\ell,m]\neq \emptyset\iff \mr{wt}^{(2)}(\mr{LCM}(g_1,\dots,g_k))=(\ell',m')\text{ with }\ell'<\ell,m'<m~.
\end{equation}
In this case
\begin{equation}
\left|\bigcap_{i=1}^kC_{g_i}[\ell,m]\right|=M[\ell-\mr{wt}_1(\mr{LCM}(g_1,\dots,g_k)),m-\mr{wt}_2(\mr{LCM}(g_1,\dots,g_k))]~.    
\end{equation}
Then using the fact that for sets $A_1,\dots,A_M$, we get
\begin{equation}
    \left|\bigcup_{i=1}^MA_i\right|=\sum_{k=1}^M(-1)^{k+1} \sum_{1 \leq i_1<\cdots<i_k \leq M}\left|A_{i_1} \cap \cdots \cap A_{i_k}\right|~,
\end{equation}
we get 
\begin{equation}
\begin{split}
C[\ell,m]&= \sum_{S\subseteq \mr{LM}(G_n)_{\ell,m}}(-1)^{|S|+1}\left|\bigcap_{g\in S}C_{g}[\ell,m]\right|
\\
&=\sum_{S\subseteq \mr{LM}(G_n)_{\ell,m}}(-1)^{|S|+1}M[\ell-\mr{wt}_1(\mr{LCM}(S)),m-\mr{wt}_2(\mr{LCM}(S))]~,
\end{split}
\end{equation}
where $|S|$ denotes the number of monomials in $S$, $\mr{LCM}(S)$ is least common multiple of the monomials in $S$. This gives us the final formula for $\mr{dim}\,(R_n)_{\ell,m}$:
\begin{equation}
    \mr{dim}\,(R_n)_{\ell,m}=M[\ell,m]+\sum_{S\subseteq \mr{LM}(G_n)_{\ell,m}}(-1)^{|S|}M[\ell-\mr{wt}_1(\mr{LCM}(S)),m-\mr{wt}_2(\mr{LCM}(S))]~.
\end{equation}
% Then a monomial $g\in \mr{LM}(G_n)_{\ell,m}$ of weight $(\ell',m')$ contributes $M[\ell-\ell',m-m']$ to $C[\ell,m]$. But note that if $\mr{LCM}(g_1,g_2)\neq 0$ then there is an overlap in the contributions due to $g_1,g_2$.  
% $\mr{dim}\,(R_n)_{\ell,m}$, we 
Substituting this in the formula
\begin{equation}
H_{R_n}(q, t)=\sum_{\ell=0}^n \sum_{m \geq 0} \mr{dim}\left(R_n\right)_{\ell, m} q^{\ell} t^m~,
\end{equation}
and using \eqref{eq:gen_free_words}, we obtain the formula 
\begin{equation}
\label{eq: hilbert_ser_Rn'}
H_{R_n}(q,t) = \mr{Ser}_{q=0,t=0}\frac{1+\sum_{S\subseteq \mr{LM}(G_n)}(-1)^{|S|}(q,t)^{\mr{wt}^{(2)}(\mr{LCM}(S))}}{\prod_{i=1}^N\prod_{j\geq 0}\left(1-q^{(d_i-1)+j}t\right)} \quad (\mr{mod} \;q^{n+1})~.
\end{equation}
The full Hilbert series is then given by taking $n\to \infty$:
\begin{equation}
\label{eq: hilbert_ser_Rinfty'}
H_{R_\infty}(q,t) = \mr{Ser}_{q=0,t=0}\frac{1+\sum_{S\subseteq \mr{LM}(G_\infty)}(-1)^{|S|}(q,t)^{\mr{wt}^{(2)}(\mr{LCM}(S))}}{\prod_{i=1}^N\prod_{j\geq 0}\left(1-q^{(d_i-1)+j}t\right)}~,
\end{equation}
where 
\begin{equation}
    \mr{LM}(G_\infty):=\lim _{\leftarrow } \mr{LM}(G_n)~.
\end{equation}
\section{Hilbert Series For The Arc Space Of $A(\CV_{(A_1,A_4)})$}\label{app:Hilb_A1A4}
The Zhu algebra for the chiral algebra of $(A_1,A_4)$ Argyres-Douglas theory is given by \eqref{eq:zhuA1A2}:
\begin{equation}
R_{(A_1,A_4)}:=A\left(\mathcal{V}_{\left(A_1, A_4\right)}\right) \cong \mathds{C}[J]/\langle J^3\rangle~.
\end{equation}
We compute the Hilbert series of its arc space 
$H_{(A(\mathcal{V}_{(A_1,A_4)})_\infty}(q,t)$ using technique discussed in Section \ref{sec:comp_Hilb_ser}. Let us  order the variables in the deglex ordering:
\begin{equation}
    J_1>J_2>\cdots
\end{equation}
Then from Mathematica computations of Gr\"{o}bner basis, we have the following conjecture: the
set of monomials of weight $(\ell,m)$ constituting a Gr\"obner basis for $\langle J(x)^2\rangle$ is given by
\begin{equation}
    \mathcal{G} (\ell,m) := \begin{cases}
\bigcup_{i = 1}^{(m-2)/2} G^{2i, m/2 - i}(\ell,m), & m \text{ is even},\\
\bigcup_{i = 1}^{(m-1)/2} G^{2i -1 , (m+1)/2 -i}(\ell,m), & m \text{ is odd},
\end{cases}
\end{equation}
where $G^{c,s}(\ell,m)$ is the set of monomials of weight $(\ell,m)$ given recursively by
\begin{equation}\label{eq:Grob_recusive_A1A4}
    G^{c,s}(\ell,m) = \begin{cases}
\{J_s^2\}\cdot F_s(\ell-2s, m -2), & c = 1, \\
\{J_s J_{s+1}\}\cdot F_{s+1}(\ell -2s -1, m -2), & c= 2,3,\\
\bigcup_{i = 1}^{c/2 -1} \{J_s\}\cdot G^{2i -1, s + c/2 +1 -i}(\ell-s , m-1), & c>3 \text{ and even}, \\
\bigcup_{i = 1}^{(c-3)/2}\{J_s\}\cdot G^{2i , s + (c+1)/2 -i}(\ell-s, m-1), & c>3 \text{ and odd,}
\end{cases}
\end{equation}
where $F_s(\ell,m)$ is defined as the set of monomials of weight $(\ell,m)$ freely generated by $\{J_i: \text{for } i\geq s \}$.
Let us write the Hilbert series in the form:
\begin{equation}
    H_{(A(\mathcal{V}_{(A_1,A_4)})_\infty}(q,t)
    \sum_{m = 0}\mathcal{I}_m^{(A_1, A_4)}(q)t^m.
\end{equation}
Let us analyze the Hilbert series to the first orders in $t$:
\begin{enumerate}
    \item $m<3$:
    Since there is no constraints till $O(t^3)$, words are freely generated by $j_i$ for $i \geq 1.$ Then using \eqref{eq:free_gen_ser_A1A4}, the Hilbert series up to $O(t^3)$ is thus given by 
    \begin{equation}
        H_{(A(\mathcal{V}_{(A_1,A_4)})_\infty}(q,t)=\mr{Ser}_{q=0,t=0}\prod_{j=1}^\infty\frac{1}{(1-q^it)}~~(\mr{mod}~t^3)= 1 + \mathcal{I}_1^1(q)t + \mathcal{I}^1_2 (q)t^2 + \mathcal{O}(t^3)~,
    \end{equation}
    where $\mathcal{I}_m^s$ is defined in \eqref{def:Ism} with its explicit expression given in \eqref{eq:Ism}.
    \item $m=3$:
    The set of constraints from Gr\"{o}bner basis is given by:
    \begin{equation}
        \mathcal{G}(\ell,3) =G^{1,1}(\ell,3) = \{j_1^2\}\cdot F_1(\ell-2,1)~.
    \end{equation}
    Thus the monomials which are freely generated by $\{j_\ell: \ell\geq1\}$ and not divisible by elements inside $\mathcal{G}(\ell,3)$ are
    \begin{equation}
        \{j_1\}\cdot F_2(\ell-1,2) \bigcup F_2(\ell,3)~,
    \end{equation}
   This set generates the homogeneous subspace in $(A(\mathcal{V}_{(A_1,A_4)})_\infty$ of weight $(\ell,m =3)$. Therefore, we have
    \begin{equation}
        \mathcal{I}^{(A_1,A_4)}_3 = q \mathcal{I}_{2}^{2} + \mathcal{I}_3^2~.
    \end{equation}
    \item $m =4$:
    The set of constraints from Gr\"{o}bner basis reads
    \begin{equation}
        \mathcal{G}(\ell,4) =G^{2,1}(\ell,4) = \{j_1j_2\}\cdot F_2(\ell-3,2)~.
    \end{equation}
    Thus the words which generate the homogeneous subspace in $(A(\mathcal{V}_{(A_1,A_4)})_\infty$ of weight $(\ell,m =4)$ are given by 
    \begin{equation}
        \{j_1\}\cdot F_3(\ell-1,3) \bigcup F_2(\ell,4)~.
    \end{equation}
   This gives us
\begin{equation}
\mathcal{I}^{(A_1,A_4)}_4 = q \mathcal{I}_{3}^{3} + \mathcal{I}_4^2~.
\end{equation}
\item
The set of constraints from Gr\"{o}bner basis reads
\begin{equation}
    \mathcal{G}(\ell,5) =G^{1,2}(\ell,5)\bigcup G^{3,1}(\ell,5) = \{j_2^2\}\cdot F_2(\ell-4,3)\bigcup\{j_1j_2\}\cdot F_2(\ell-3,3)~.
\end{equation}
The words which generates the homogeneous subspace in $(A(\mathcal{V}_{(A_1,A_2)})_\infty$ of weight $(\ell,m =5)$ are thus given by
\begin{equation}
    \{j_2\}\cdot F_3(\ell-2,4)\bigcup \{j_1\}\cdot F_3(\ell-1,4) \bigcup F_3(\ell,5)~.
\end{equation}
Thus we have
\begin{equation}
    \mathcal{I}_5^{(A_1,A_4)} = q^2 \mathcal{I}_4^3 + q \mathcal{I}_4^3 + \mathcal{I}_5^3~.
\end{equation}
\end{enumerate}
For $m\geq 6$, similar analysis can be performed using the general recursive form of the Gr\"{o}bner basis given in \eqref{eq:Grob_recusive_A1A4}. The result is that we can write
\begin{equation}
    \mathcal{I}_m^{(A_1,A_4)} = \begin{cases}\mathcal{I}_m^{m/2}+
\sum_{i=1}^{(m-2)/2}I^{2i, m/2 -i}_m, &m\text{  is even,}\\
\mathcal{I}_m^{(m+1)/2} +\sum_{i =1}^{(m-1)/2} I_m^{2i -1,(m+1)/2 -i}, & m \text{ is odd},
\end{cases}
\end{equation}
where  $I^{c,s}_m$ is determined by the following recursion relation:
\begin{equation}
I_m^{c, s}= \begin{cases}q^s \mathcal{I}_{m-1}^{s+1}, & c=1, \\ q^s \mathcal{I}_{m-1}^{s+2}, & c=2,3, \\ q^s \mathcal{I}_{m-1}^{s+c+1}+q^s \sum_{i=1}^{c / 2-1} I_{m-1}^{2 i-1, s+c / 2+1-i}, & c>3 \text { and even }, \\ q^s \mathcal{I}_{m-1}^{s+(c+1) / 2}+q^s \sum_{i=1}^{(c-3) / 2} I_{m-1}^{2 i, s+(c+1) / 2-i}, & c>3 \text { and odd~. }\end{cases}    
\end{equation}
\section{Evaluation Of Hilbert Series \eqref{eq: hilbert_ser_Rinf} At $t^3$}\label{app:Hilb_ser_t3}
In this appendix, we give the detailed calculation of the Hilbert series \eqref{eq: hilbert_ser_Rinf} at $t^3$. We want to compute the Hilbert series 
\begin{equation}
\label{eq: hilbert_ser_Rinf_app}
H_{R_\infty}(q,t) = \mr{Ser}_{q=0,t=0}\frac{1+\sum_{S\subseteq \mr{LM}(G_\infty)}(-1)^{|S|}(q,t)^{\mr{wt}^{(2)}(\mr{LCM}(S))}}{(1-t)^3 \prod_{i=1}^{\infty}\left(1-q^i t\right)^4}~,
\end{equation}
where $\mr{LM}(G_n)$ is given by
\begin{equation}
\label{eq:Gr\"{o}bner_basis_Gn_app}
\mr{LM}\left(G_n\right)=\bigcup_{\substack{0 \leq \ell \leq n \\ 2 \leq m \leq \ell+2}} I_{\ell, m}~,
\end{equation}
where $I_{\ell,m}$ is a homogeneous set of weight $(\ell,m)$ given by
\begin{equation}
I_{\ell,m} := \left(\{X_{m-2}\}\cdot F_z(\ell-m+2, m-1)\right)\bigcup \left(\{J_{m-1}\}\cdot F(\ell-m+1, m-1)\right)~.
\end{equation}
Let us write the numerator of \eqref{eq: hilbert_ser_Rinf_app} as
\begin{equation}
\sum_{S \subseteq \mr{L M}\left(G_{\infty}\right)}(-1)^{|S|}(q, t)^{\mr{wt}^{(2)}(\mr{LCM}(S))}=: \sum_{\substack{\ell \geq 0\\ m\geq2}} C[\ell,m]q^\ell t^m~.
\end{equation}
there are several non-trivial contributions to $C[\ell, m]$. The first contribution comes from the subsets $S\subset I_{\ell,m}\subseteq \mr{LM}(G_\ell)$. Note that the sets $S\subset I_{\ell,m}$ contributing to $C[\ell,m]$ are necessarily singletons since \tsf{LCM} of any two elements of $I_{\ell,m}$ has weight $(\ell',m')$ with $m'>m$. For this reason, the constraints obtained from $S\subset I_{\ell,m}$ are called \textit{first-depth constraints}. Other contributions come from subsets of 
\begin{equation}
\bigcup_{\substack{\ell'\leq \ell\\m'<m}}I_{\ell',m'}~.    
\end{equation}
If $|S|=k$, then the constraints obtained from $S$ are called \textit{depth-k constraints}. A good way to organize these constributions is as follows: given $k$ sets $S_1,S_2,\dots S_k$ of monomials, define
\begin{equation}
    \mr{LCM}(S_1,S_2,\dots S_k) := \sum_{\substack{s_1 \in S_1,\dots, s_k\in S_k\\
    s_i \neq s_j}} 
    \mr{LCM}(s_1,s_2,\dots, s_k)~,
\end{equation}
and define $\mr{LCM}^{(\ell,m)}(S_1,S_2,\dots S_k)$ to be the maximal homogeneous polynomial which is a summand in $\mr{LCM}(S_1,S_2,\dots S_k)$ of weight $(\ell,m)$. Then the number of depth-$k$ contributions to $C[\ell,m]$ is given by 
\begin{equation}
    \sum_{S_1,\dots,S_k\subseteq \mr{LM}(G_\ell)}(-1)^k|\mr{LCM}^{(\ell,m)}(S_1,S_2,\dots, S_k)|~,
\end{equation}
where for a homogeneous polynomial $f$, $|f|$ is the number of terms in the polynomial counted with multiplicity\footnote{This means that we write the polynomial such that each monomial has coefficient 1 and then count the number of terms.}. Thus, we can write 
%\myRed{Can the sum over $k$ be bounded above? Runkai: there should be a bound, but don't know what.}
\begin{equation}
    C[\ell,m]=\sum_{k=1}^{|\mr{LM}(G_\ell)|}\sum_{S_1,\dots,S_k\subseteq \mr{LM}(G_\ell)}(-1)^k|\mr{LCM}^{(\ell,m)}(S_1,S_2,\dots S_k)|~.
\end{equation}
\par
We need to obtain all the coefficients $C[\ell,m=3]$ for $\ell\geq 0$. From \eqref{eq:Gr\"{o}bner_basis_Gn_app} with $n=0$, it is not hard to see  that
\begin{equation}
    C[\ell=0, m =3] = 0~.
\end{equation}
For $\ell>0$, we list the constraints contributing to $C[\ell, m =3]$ below. 

\begin{enumerate}
\item $\boldsymbol{C[\ell=1, m=3].}$

From \eqref{eq:Gr\"{o}bner_basis_Gn_app} with $n=1$, the depth-1 constraint is given by
\begin{equation}
    \mr{LCM}^{(1,3)}(\mr{LM}(G_1)) = X_1 Z_0^2~.
\end{equation}
The non-trivial contributions for the depth-2 constraints
are given by 
\begin{equation}
\begin{aligned}
    \mr{LCM}^{(1,3)}(\{X_0 Z_0\}, \{X_0 Z_1\}) &= X_0 Z_0 Z_1~,\\
    \mr{LCM}^{(1,3)}(\{X_0 Z_0\}, \{J_1\}\cdot F(0,1))&= 2 J_1 X_0 Z_0~, \\
    \mr{LCM}^{(1,3)}(\{J_1\}\cdot F(0,1), \{J_1\} \cdot F(0,1)) &= 
    J_1(X_0Y_0+Y_0Z_0+X_0Z_0)~.
\end{aligned}
\end{equation}
Depth-3 contribution are given by
\begin{equation}
    \mr{LCM}^{(1,3)}(\{X_0Z_0\},\{J_1\}\cdot F(0,1), \{J_1\}\cdot F(0,1)) = 
    J_1X_0Z_0~.
\end{equation}
Then we have
\begin{equation}
    C[\ell=1,m=3]=-1+(1+2+3)-1=4~.
\end{equation}
\item $\boldsymbol{C[\ell=2, m = 3].}$

The depth-1 contribution comes from
\begin{equation}
\mr{LCM}^{(2,3)}\left(\mr{L M}\left(G_2\right)\right)=X_1\overline{ F_z(1,2)}+J_2 \overline{F(0,2)}~,
\end{equation}
where given a set $S$, we have defined 
\begin{equation}
    \overline{S}:=\sum_{s\in S}s~.
\end{equation}

The depth-2 constraints are given by:
\begin{equation}
\begin{aligned}
\mr{LCM}^{(2,3)}(\{X_0 Z_0\},\{X_0 Z_2\}) &= X_0Z_0 Z_2~,\\
\mr{LCM}^{(2,3)}(\{X_0 Z_0\},\{J_2\}\cdot F(0,2)) &= J_2 X_0 Z_0~,\\
\mr{LCM}^{(2,3)}(\{X_0 Z_1\},\{J_1\}\cdot F(0,1)) &= J_2 X_0 Z_0~,\\
\mr{LCM}^{(2,3)}(\{X_0 Z_1\},\{J_1\}\cdot F(1,1)) &= J_1 X_0 Z_1~,\\
\mr{LCM}^{(2,3)}(J_1F(0,1),\{J_1\}\cdot F(1,1)) &=J_1 \overline{F(1,2)}~.
    \end{aligned}
\end{equation}
The depth-3 constraints are given by:
\begin{equation}
    \mr{LCM}^{(2,3)}
    (\{X_0 Z_1\}, \{J_1\}\cdot F(0,1), \{J_1\}\cdot F(1,1)) =
    X_0 Z_1 J_1~.
\end{equation}
Then we have 
\begin{equation}
    C[\ell=2, m=3] = -(|F_z(1,2)|+ |F(0,2)|)+ (4 + |F(1,2)|) -1 = 8~, 
\end{equation}
where we used  \eqref{eq:num_Fl2} and 
\begin{equation}
\label{eq:num_Fzl2}
    |F_z(\ell,2)| =\begin{cases}
        \ell/2 + 1, & \ell \text{ even},\\
        (\ell+1)/2, 
        & \ell \text{ odd}~.
    \end{cases}
\end{equation}
\item $\boldsymbol{C[\ell, m=3]}$ \textbf{for} $\boldsymbol{\ell>2.}$

The depth-1 constraints are given by 
\begin{equation}
    \mr{LCM}^{(\ell,3)}(\mr{LM}(G_\ell)) = X_1 \overline{F_z(\ell-1,2)}+ J_2 \overline{F(\ell-2,2)}~.
\end{equation}
The depth-2 constraints are given by
\begin{equation}
\begin{aligned}
\mr{LCM}^{(\ell,3)}(\{X_0 Z_{\ell-1}\},\{J_1 Z_{\ell-1}\}) &= J_1 Z_{\ell-1}X_0~, \\
    \mr{LCM}^{(\ell,3)}(\{X_0 Z_{\ell-1}\}, \{J_1 X_0\}) &= J_1 Z_{\ell-1} X_0~, \\
    \mr{LCM}^{(\ell,3)}(\{X_0 Z_{\ell-2}\}, \{J_2 X_0Z_{\ell-2}\}) &= J_2 Z_{\ell-2}X_0~ ,\\
    \sum_{i=0}^{\lfloor \ell/2\rfloor}\mr{LCM}^{(\ell,3)}(\{X_0 Z_i\},\{X_0 Z_{\ell-i}\}) &= \sum_{i=0}^{\lfloor \ell/2\rfloor} X_0 Z_i Z_{\ell-i}~,\\
    \sum_{i=0}^{\lfloor \ell/2\rfloor} \mr{LCM}^{(\ell,3)}(J_1 F(i,1), J_1F(\ell-1-i,1)) &= 
    \begin{cases}
    J_1 \overline{F(\ell-1,2)}, & \ell \text{ even}~,\\
    J_1 \overline{F(\ell-1,2)}-J_1\big(X^2_{\frac{\ell-1}{2}}+Y^2_{\frac{\ell-1}{2}}+Z^2_{\frac{\ell-1}{2}}+ J^2_{\frac{\ell-1}{2}}\big), & \ell \text{ odd}~,
    \end{cases}\\
    \mr{LCM}^{(\ell,3)}(J_1 F(\ell-3,1), J_2 F(\ell-2,2)) &= J_1 J_2 \overline{F(\ell-3,1)}~,\\
    \mr{LCM}^{(\ell,3)}(J_1 F(\ell-1,1), J_2 F(\ell-2,2)) &= J_1 J_2 \overline{F(\ell-3,1)}~.
\end{aligned}
\end{equation}
The depth-3 constraints are given by:
\begin{equation}
    \begin{aligned}
        \mr{LCM}^{(\ell,3)}(\{X_0 Z_{\ell-1}\},\{J_1 Z_{\ell-1}\},\{J_1 X_0\}) &= J_1 Z_{\ell-1}X_0~,\\
        \mr{LCM}^{(\ell,3)}(J_1 F(\ell-1,1), J_2 F(\ell-2,2), J_1 F(\ell-3,1)) &= J_1 J_2 \overline{F(\ell-3,1)}~.
    \end{aligned}
\end{equation}
Then 
\begin{equation}
    \begin{aligned}C[\ell,m=3]&=\begin{cases}
        -|F_z(\ell-1,2)|-|F(\ell-2,2)| + |F(\ell-3,1)| & l \text{ even},\\
        \quad+|F_z(\ell,2)|+|F(\ell-1,2)| + 1,\\
        -|F_z(\ell-1,2)|-|F(\ell-2,2)| + |F(\ell-3,1)| & l \text{ odd},\\
        \quad+|F_z(\ell,2)|+|F(\ell-1,2)| -2,
    \end{cases}\\
    &= \begin{cases}
        11, & \ell = 3~,\\
        12, & \ell > 3~,
    \end{cases}
\end{aligned}
\end{equation}
where we used \eqref{eq:num_Fl2} and 
\eqref{eq:num_Fzl2}.
\end{enumerate}
Thus we get
\begin{equation}
\sum_{\substack{S \subseteq \mr{L M}\left(G_{\infty}\right) \\ \mr{wt}_2(\mr{LCM}(S))=3}}(-1)^{|S|}(q, t)^{\mathsf{wt}^{(2)}(\mr{LCM}(S))}=t^3\left(4 q+8 q^2+11 q^3+12 \sum_{\ell \geq 4} q^\ell\right)~.
\end{equation}
The Hilbert series up to $t^3$ is then given by
\begin{equation}\label{eq:Hil_ser_modt4}
\begin{aligned}
H_{R_\infty}(q,t) &=\mr{Ser}_{q=0,t=0} \frac{1-t^2(1+4q + 5\sum_{\ell=2}^\infty q^\ell)+t^3(4q + 8 q^2 + 11 q^3 + 12 \sum_{\ell\geq 4}q^\ell)}{(1-t)^3 \prod_{i=1}^{\infty}\left(1-q^i t\right)^4}\quad (\mr{mod} \;t^4)~.
\end{aligned}
\end{equation}
Let us focus on $t^3$ term
\begin{equation}
H_{R_\infty}(q,t)|_{t^3} = \sum_{\ell=0}^\infty \mr{dim}(R_\infty)_{\ell,3} q^\ell~.
\end{equation}
From \eqref{eq:Hil_ser_modt4}, we see that the coefficients are given by
\begin{equation}
\label{eq:coefficients_t3}
\begin{aligned}
\mr{dim}(R_\infty)_{0,3} &= |F(0,3)| - |F(0,1)|~, \\
\mr{dim}(R_\infty)_{1,3} &= |F(1,3)| - |F(1,1)| - 4 |F(0,1)|+4|F(0,0)|~,\\
\mr{dim}(R_\infty)_{2,3} &= |F(2,3)|-|F(2,1)|-4|F(1,1)|-5|F(0,1)|+8 |F(1,1)|~,\\
\mr{dim}(R_\infty)_{3,3} & = |F(3,3)|-|F(3,1)|-4 |F(2,1)|-5 |F(1,1)|-5|F(0,1)|+ 11 |F(0,0)|~,\\
\mr{dim}(R_\infty)_{\ell,3} &= |F(\ell,3)|-|F(\ell,1)|-4|F(\ell-1,1)|-5\sum_{i=0}^{\ell -2} |F(i,1)|+ 12 |F(0,0)| \quad 
\ell >3~.
\end{aligned}
\end{equation}
Note that 
\begin{equation}
F(\ell,1) = \begin{cases}
\{X_0,Y_0,Z_0\} & \ell = 0~,\\
 \{X_\ell, Y_\ell,Z_\ell, J_\ell\} & \ell >0~.
\end{cases}    
\end{equation}
Thus, we have
\begin{equation}
\label{eq:Fl1}
|F(\ell,1)| = \begin{cases}
3 & \ell = 0~,\\
4 & \ell >0~.
\end{cases}
\end{equation}
Next, $|F(0,3)| = 10$, since the elements of $F(0,3)$ is constructed using the variables $\{X_0, Y_0, Z_0\}$.
For $\ell>0$, the number $|F(\ell,3)|$ is counted in six cases below:
\begin{itemize}
\item
The number of monomials constructed  using one variable from $\{X_\ell,Y_\ell,Z_\ell,J_\ell\}$ and two variables from $\{X_0,Y_0,Z_0\}$ is $4\times(3\times 2)=24$.
\item
The number of monomials constructed using one variable from $\{X_i, Y_i,Z_i,J_i\}$, one from $\{X_j, Y_j,Z_j,J_j\}$ and one from $\{X_0,Y_0,Z_0\}$, such that $i>j>0$ and $i+j = \ell$, is 
\begin{equation}
    4\times 4 \times 3 \bigg\lfloor\frac{\ell-1}{2}\bigg\rfloor~.
\end{equation}
\item 
The number of monomials constructed using one variable from $\{X_i, Y_i,Z_i,J_i\}$ and two variables from $\{X_j, Y_j,Z_j,J_j\}$, such that $i >0, j>0,i\neq j, i + 2j = \ell$, is  
\begin{equation}
4 \times 10 \left(\bigg\lfloor \frac{\ell-1}{2}\bigg\rfloor+\bigg\lfloor \frac{\ell-1}{3}\bigg\rfloor - \bigg\lfloor \frac{\ell}{3}\bigg\rfloor\right)~.
\end{equation}
\item 
The number of monomials constructed using one variable from $\{X_0,Y_0,Z_0\}$ and two variables from  $\{X_i, Y_i,Z_i,J_i\}$ , such that $i>0, 2i = \ell$, is 
\begin{equation}
    3 \times 10\left(\bigg\lfloor \frac{\ell}{2}\bigg\rfloor - \bigg\lfloor \frac{\ell-1}{2}\bigg\rfloor \right)~.
\end{equation}
\item 
The number of monomials constructed using three variables from  $\{X_i, Y_i,Z_i,J_i\}$ , such that $i>0, 3i = \ell$, is 
\begin{equation}
    20\left(\bigg\lfloor \frac{\ell}{3}\bigg\rfloor - \bigg\lfloor \frac{\ell-1}{3}\bigg\rfloor \right)~.
\end{equation}
\item 
The number of monomials constructed using three variables from three different sets $\{X_i,Y_i,Z_i,J_i\}$, $\{X_j,Y_j,Z_j,J_j\}$ and $\{X_k, Y_k, Z_k, J_k\}$, such that $i>j>k>0, i+j+k = \ell$, is
\begin{equation}
4\times 4 \times 4 \bigg\lfloor \frac{(\ell-3)^2 + 6}{12}\bigg\rfloor~.    
\end{equation}
\end{itemize}
Therefore, for $\ell > 0$, we get 
\begin{equation}
\label{eq:Fl3}
|F(\ell,3)| = 24 + 64\bigg\lfloor \frac{(\ell-3)^2 + 6}{12}\bigg\rfloor +20\bigg\lfloor \frac{\ell-1}{3}\bigg\rfloor + 58\bigg\lfloor \frac{\ell-1}{2}\bigg\rfloor -20 \bigg\lfloor \frac{\ell}{3}\bigg\rfloor+ 30\bigg\lfloor \frac{\ell}{2}\bigg\rfloor~.
\end{equation}
Combining \eqref{eq:coefficients_t3},\eqref{eq:Fl1} and \eqref{eq:Fl3}, then the coefficients are
\begin{equation}
\label{eq:coeff_t3}
\begin{aligned}
\mr{dim}(R_\infty)_{0,3} &= 7~, \\
\mr{dim}(R_\infty)_{1,3} &= 12~,\\
\mr{dim}(R_\infty)_{2,3} &= 27~,\\
\mr{dim}(R_\infty)_{3,3} & = 48~,\\
\mr{dim}(R_\infty)_{\ell,3} &= 41-20\ell + 64\bigg\lfloor \frac{(\ell-3)^2 + 6}{12}\bigg\rfloor +20\bigg\lfloor \frac{\ell-1}{3}\bigg\rfloor + 58\bigg\lfloor \frac{\ell-1}{2}\bigg\rfloor -20 \bigg\lfloor \frac{\ell}{3}\bigg\rfloor+ 30\bigg\lfloor \frac{\ell}{2}\bigg\rfloor~.
\end{aligned}
\end{equation}
We separate the above coefficients into 12 groups by taking $\ell=12 k+i$, where $i=4,5, \ldots 15,$ and $k\in\IZ_{\geq 0}$. All floor functions can then be removed:
\begin{equation}
\begin{aligned}
\mr{dim}(R_\infty)_{12k+4,3} & =768 k^2+416 k+79~ ,\\
\mr{dim}(R_\infty)_{12k+5,3} & =768 k^2+544 k+117~ ,\\
\mr{dim}(R_\infty)_{12k+6,3} & =768 k^2+672 k+171~, \\
\mr{dim}(R_\infty)_{12k+7,3} & =768 k^2+800 k+229~, \\
\mr{dim}(R_\infty)_{12k+8,3} & =768 k^2+928 k+303~, \\
\mr{dim}(R_\infty)_{12k+9,3} & =768 k^2+1056 k+385~, \\
\mr{dim}(R_\infty)_{12k+10,3} & =768 k^2+1184 k+479~, \\
\mr{dim}(R_\infty)_{12k+11,3} & =768 k^2+1312 k+581~, \\
\mr{dim}(R_\infty)_{12k+12,3}& =768 k^2+1440 k+699~, \\
\mr{dim}(R_\infty)_{12k+13,3} & =768 k^2+1568 k+821~, \\
\mr{dim}(R_\infty)_{12k+14,3} & =768 k^2+1696 k+959~, \\
\mr{dim}(R_\infty)_{12k+15,3} & =768 k^2+1824 k+1105~.
\end{aligned}
\end{equation}
Let us denote
\begin{equation}
\mr{dim}(R_\infty)_{12k+i,3}=a_i k^2+b_i k+c_i~.
\end{equation}
The $i$-th sector can be summed over as follows:
\begin{equation}
\sum_{k=0}^{\infty} \mr{dim}(R_\infty)_{12k+i,3} q^{12 k+i}=q^i \frac{a_i q^{12}\left(1+q^{12}\right)+b_i q^{12}\left(1-q^{12}\right)+c_i\left(1-q^{12}\right)^2}{\left(1-q^{12}\right)^3}~.
\end{equation}
Combing \eqref{eq:coeff_t3}, by straightforward computation, we can show that
\begin{equation}\label{eq:I-3_ser}
\begin{aligned}
\sum_{\ell=0}^\infty \mr{dim}(R_\infty)_{\ell,3} q^\ell&= 7+12 q+27 q^2+48 q^3+\sum_{i=4}^{15} q^i \frac{a_i q^{12}\left(1+q^{12}\right)+b_i q^{12}\left(1-q^{12}\right)+c_i\left(1-q^{12}\right)^2}{\left(1-q^{12}\right)^3} \\
& =\left(1-q^4\right)\left(1-q^5\right)\left(1-q^6\right) \sum_{k=0}^6 \frac{1}{(q ; q)_k(q ; q)_{6-k}}
\\
&=\frac{1}{(q;q)_3} \sum_{k=0}^{6}\binom{6}{k}_q~,
\end{aligned}
\end{equation}
where the second last equality follows from explicit Mathematica calculation. 
This matches with the coefficient of $t^3$ in our formula \eqref{eq:Mac_ind_A1D3}.
\section{Numerical Computation Of $f^n_{m,l}(q)$}\label{app:data_fmnl}
In this section, we present the computation of the polynomials $f^n_{m,l}(q)$ for $n\geq 2$. These are consistent with the general formula \eqref{eq:fnml_closed_form_sec5}. 
In our computation, we noticed a curious recursion relation for 
 $f^n_{m,l}(q)$:
 \begin{equation}\label{eq:rec_fnml}
     f^{n}_{m,l} = \begin{cases}
         0, \quad m \leq 0\; \text{or}\; l \leq 0\;, \text{or}\; m > n l~,\\
         1,\quad m=l> 0~, \\
         q^{(n-1) l^2}, \quad m = n l~,\\
         q^l f^{n}_{m-1,l} + f^{n}_{m-1,l-1}-q^{2m-l-n-1} f^{n}_{m-n-1,l-1}, \quad m/n< l < m~.
     \end{cases}
 \end{equation}
We show in Appendix \ref{app:new_id} that the formula for $f^n_{m,l}$ satisfies this recursion relation for $n=2$ and leave the general proof for future work. We do not yet have a physical interpretation of these recursion relations. 
\begin{table}[H]
    \centering
    \begin{tabular}{|c|l|}
\hline
$(n,m)$ & $f^n_{m,l}(q)$ for $l=1,\dots,m$
\\
\hline
$(2,1)$ & 
$f^2_{1,1} = 1$ \\
\hline
$(2,2)$ & 
$f^2_{2,1} = q$ \\
        & $f^2_{2,2} = 1$ \\
\hline
$(2,3)$ & 
$f^2_{3,1} = 0$ \\
        & $f^2_{3,2} = q + q^2$ \\
        & $f^2_{3,3} = 1$ \\
\hline
$(2,4)$ & 
$f^2_{4,1} = 0$ \\
        & $f^2_{4,2} = q^4$ \\
        & $f^2_{4,3} = q + q^2 + q^3$ \\
        & $f^2_{4,4} = 1$ \\
\hline
$(2,5)$ & 
$f^2_{5,1} = f^2_{5,2} = 0$ \\
        & $f^2_{5,3} = q^4 + q^5 + q^6$ \\
        & $f^2_{5,4} = q + q^2 + q^3 + q^4$ \\
        & $f^2_{5,5} = 1$ \\
\hline
$(2,6)$ & 
$f^2_{6,1} = f^2_{6,2} = 0$ \\
        & $f^2_{6,3} = q^9$ \\
        & $f^2_{6,4} = q^4 + q^5 + 2q^6 + q^7 + q^8$ \\
        & $f^2_{6,5} = q + q^2 + q^3 + q^4 + q^5$ \\
        & $f^2_{6,6} = 1$ \\
\hline
$(2,7)$ & 
$f^2_{7,1} = f^2_{7,2} = f^2_{7,3} = 0$ \\
        & $f^2_{7,4} = q^9 + q^{10} + q^{11} + q^{12}$ \\
        & $f^2_{7,5} = q^4 + q^5 + 2q^6 + 2q^7 + 2q^8 + q^9 + q^{10}$ \\
        & $f^2_{7,6} = q + q^2 + q^3 + q^4 + q^5 + q^6$ \\
        & $f^2_{7,7} = 1$ \\
\hline
$(2,8)$ & 
$f^2_{8,1} = f^2_{8,2} = f^2_{8,3} = 0$ \\
        & $f^2_{8,4} = q^{16}$ \\
        & $f^2_{8,5} = q^9 + q^{10} + 2q^{11} + 2q^{12} + 2q^{13} + q^{14} + q^{15}$ \\
        & $f^2_{8,6} = q^4 + q^5 + 2q^6 + 2q^7 + 3q^8 + 2q^9 + 2q^{10} + q^{11} + q^{12}$ \\
        & $f^2_{8,7} = q + q^2 + q^3 + q^4 + q^5 + q^6 + q^7$ \\
        & $f^2_{8,8} = 1$ \\
\hline
$(2,9)$ & 
$f^2_{9,1} = f^2_{9,2} = f^2_{9,3} = f^2_{9,4} = 0$ \\
        & $f^2_{9,5} = q^{16} + q^{17} + q^{18} + q^{19} + q^{20}$ \\
        & $f^2_{9,6} = q^9 + q^{10} + 2q^{11} + 3q^{12} + 3q^{13} + 3q^{14} + 3q^{15} + 2q^{16} + q^{17} + q^{18}$ \\
        & $f^2_{9,7} = q^4 + q^5 + 2q^6 + 2q^7 + 3q^8 + 3q^9 + 3q^{10} + 2q^{11} + 2q^{12} + q^{13} + q^{14}$ \\
        & $f^2_{9,8} = q + q^2 + q^3 + q^4 + q^5 + q^6 + q^7 + q^8$ \\
        & $f^2_{9,9} = 1$ \\
\hline
\end{tabular}
    \caption{Mathematica computations of the polynomials $f^2_{m,l}$.}
    \label{tab:my_label}
\end{table}
\begin{table}[H]
    \centering
    \begin{tabular}{|c|l|}
    \hline
    $(n,m)$ & $f^n_{m,l}(q)$ for $l=1,\dots,m$
\\
    \hline
$(2,10)$ & 
$f^2_{10,1} = f^2_{10,2} = f^2_{10,3} = f^2_{10,4} = 0$ \\
         & $f^2_{10,5} = q^{25}$ \\
         & $f^2_{10,6} = q^{16} + q^{17} + 2q^{18} + 2q^{19} + 3q^{20} + 2q^{21} + 2q^{22} + q^{23} + q^{24}$ \\
         & $f^2_{10,7} = q^9 + q^{10} + 2q^{11} + 3q^{12} + 4q^{13} + 4q^{14} + 5q^{15} + 4q^{16} + 4q^{17} + 3q^{18} + 2q^{19} + q^{20} + q^{21}$ \\
         & $f^2_{10,8} = q^4 + q^5 + 2q^6 + 2q^7 + 3q^8 + 3q^9 + 4q^{10} + 3q^{11} + 3q^{12} + 2q^{13} + 2q^{14} + q^{15} + q^{16}$ \\
         & $f^2_{10,9} = q + q^2 + q^3 + q^4 + q^5 + q^6 + q^7 + q^8 + q^9$ \\
         & $f^2_{10,10} = 1$ \\
\hline
$(2,11)$ & 
$f^2_{11,1} = f^2_{11,2} = f^2_{11,3} = f^2_{11,4} = f^2_{11,5} = 0$ \\
         & $f^2_{11,6} = q^{25} + q^{26} + q^{27} + q^{28} + q^{29} + q^{30}$ \\
         & \begin{tabular}[t]{@{}l@{}}$f^2_{11,7} = q^{16} + q^{17} + 2q^{18} + 3q^{19} + 4q^{20} + 4q^{21} + 5q^{22}+ 4q^{23} + 4q^{24}+ 3q^{25} + 2q^{26} + q^{27} + q^{28} $\end{tabular} \\
         & \begin{tabular}[t]{@{}l@{}}$f^2_{11,8} = q^9 + q^{10} + 2q^{11} + 3q^{12} + 4q^{13} + 5q^{14} + 6q^{15}+ 6q^{16} + 6q^{17} + 6q^{18} + 5q^{19}+ 4q^{20} + 3q^{21}$ \\
         \quad\quad~$+ 2q^{22} + q^{23} + q^{24} $\end{tabular} \\
         & \begin{tabular}[t]{@{}l@{}}$f^2_{11,9} = q^4 + q^5 + 2q^6 + 2q^7 + 3q^8 + 3q^9 + 4q^{10} + 4q^{11}+ 4q^{12} + 3q^{13} + 3q^{14} + 2q^{15}+ 2q^{16}  $ \\
         \quad\quad~$+ q^{17} + q^{18}$\end{tabular} \\
         & $f^2_{11,10} = q + q^2 + q^3 + q^4 + q^5 + q^6 + q^7 + q^8 + q^9 + q^{10}$ \\
         & $f^2_{11,11} = 1$ \\
\hline
$(2,12)$ & 
$f^2_{12,1} = f^2_{12,2} = f^2_{12,3} = f^2_{12,4} = f^2_{12,5} = 0$ \\
         & $f^2_{12,6} = q^{36}$ \\
         & \begin{tabular}[t]{@{}l@{}}$f^2_{12,7} = q^{25} + q^{26} + 2q^{27} + 2q^{28} + 3q^{29} + 3q^{30} + 3q^{31}+ 2q^{32} + 2q^{33} + q^{34} + q^{35}$ \end{tabular} \\
         & \begin{tabular}[t]{@{}l@{}}$f^2_{12,8} = q^{16} + q^{17} + 2q^{18} + 3q^{19} + 5q^{20} + 5q^{21} + 7q^{22}+ 7q^{23} + 8q^{24} + 7q^{25} + 7q^{26}+ 5q^{27} + 5q^{28} $ \\
         \quad\quad~$+ 3q^{29} + 2q^{30} + q^{31} + q^{32}$\end{tabular} \\
         & \begin{tabular}[t]{@{}l@{}}$f^2_{12,9} = q^9 + q^{10} + 2q^{11} + 3q^{12} + 4q^{13} + 5q^{14} + 7q^{15}+ 7q^{16} + 8q^{17} + 8q^{18} + 8q^{19}+ 7q^{20} + 7q^{21} $ \\
         \quad\quad~$+ 5q^{22} + 4q^{23} + 3q^{24} + 2q^{25} + q^{26} + q^{27}$\end{tabular} \\
         & \begin{tabular}[t]{@{}l@{}}$f^2_{12,10} = q^4 + q^5 + 2q^6 + 2q^7 + 3q^8 + 3q^9 + 4q^{10} + 4q^{11}+ 5q^{12} + 4q^{13} + 4q^{14} + 3q^{15}+ 3q^{16}  $ \\
         \quad\quad ~~$+ 2q^{17} + 2q^{18}+ q^{19} + q^{20}$\end{tabular} \\
         & $f^2_{12,11} = q + q^2 + q^3 + q^4 + q^5 + q^6 + q^7 + q^8 + q^9 + q^{10} + q^{11}$ \\
         & $f^2_{12,12} = 1$ \\
\hline
\end{tabular}
    \caption{Mathematica computations of the polynomials $f^2_{m,l}$ continued.}
    \label{tab:my_label}
\end{table}
\begin{table}[H]
    \centering
    \begin{tabular}{|c|l|}
\hline
$(n,m)$ & $f^n_{m,l}(q)$ for $l=1,\dots,m$ \\
\hline
$(3,1)$ & $f^3_{1,1} = 1$ \\
\hline
$(3,2)$ & $f^3_{2,1} = q$ \\
        & $f^3_{2,2} = 1$ \\
\hline
$(3,3)$ & $f^3_{3,1} = q^2$ \\
        & $f^3_{3,2} = q + q^2$ \\
        & $f^3_{3,3} = 1$ \\
\hline
$(3,4)$ & $f^3_{4,1} = 0$ \\
        & $f^3_{4,2} = q^2 + q^3 + q^4$ \\
        & $f^3_{4,3} = q + q^2 + q^3$ \\
        & $f^3_{4,4} = 1$ \\
\hline
$(3,5)$ & $f^3_{5,1} = 0$ \\
        & $f^3_{5,2} = q^5 + q^6$ \\
        & $f^3_{5,3} = q^2 + q^3 + 2q^4 + q^5 + q^6$ \\
        & $f^3_{5,4} = q + q^2 + q^3 + q^4$ \\
        & $f^3_{5,5} = 1$ \\
\hline
$(3,6)$ & $f^3_{6,1} = 0$ \\
        & $f^3_{6,2} = q^8$ \\
        & $f^3_{6,3} = q^5 + 2q^6 + 2q^7 + q^8 + q^9$ \\
        & $f^3_{6,4} = q^2 + q^3 + 2q^4 + 2q^5 + 2q^6 + q^7 + q^8$ \\
        & $f^3_{6,5} = q + q^2 + q^3 + q^4 + q^5$ \\
        & $f^3_{6,6} = 1$ \\
\hline
$(3,7)$ & $f^3_{7,1} = f^3_{7,2} = 0$ \\
        & $f^3_{7,3} = q^8 + q^9 + 2q^{10} + q^{11} + q^{12}$ \\
        & $f^3_{7,4} = q^5 + 2q^6 + 3q^7 + 3q^8 + 3q^9 + 2q^{10} + q^{11} + q^{12}$ \\
        & $f^3_{7,5} = q^2 + q^3 + 2q^4 + 2q^5 + 3q^6 + 2q^7 + 2q^8 + q^9 + q^{10}$ \\
        & $f^3_{7,6} = q + q^2 + q^3 + q^4 + q^5 + q^6$ \\
        & $f^3_{7,7} = 1$ \\
\hline
$(3,8)$ & $f^3_{8,1} = f^3_{8,2} = 0$ \\
        & $f^3_{8,3} = q^{13} + q^{14} + q^{15}$ \\
        & $f^3_{8,4} = q^8 + q^9 + 3q^{10} + 3q^{11} + 4q^{12} + 3q^{13} + 2q^{14} + q^{15} + q^{16}$ \\
        & $f^3_{8,5} = q^5 + 2q^6 + 3q^7 + 4q^8 + 5q^9 + 4q^{10} + 4q^{11} + 3q^{12} + 2q^{13} + q^{14} + q^{15}$ \\
        & $f^3_{8,6} = q^2 + q^3 + 2q^4 + 2q^5 + 3q^6 + 3q^7 + 3q^8 + 2q^9 + 2q^{10} + q^{11} + q^{12}$ \\
        & $f^3_{8,7} = q + q^2 + q^3 + q^4 + q^5 + q^6 + q^7$ \\
        & $f^3_{8,8} = 1$ \\
\hline
\end{tabular}
    \caption{Mathematica computations of the polynomials $f^3_{m,l}$.}
    \label{tab:my_label}
\end{table}
\begin{table}[H]
    \centering
    \begin{tabular}{|c|l|}
    \hline
$(n,m)$ & $f^n_{m,l}(q)$ for $l=1,\dots,m$ \\
\hline
$(3,9)$ & $f^3_{9,1} = f^3_{9,2} = 0$ \\
&$f^3_{9,3} = q^{18}$
\\
       & $f^3_{9,4} = q^{13}+2q^{14}+3q^{15}+3q^{16}+3q^{17}+2q^{18}+q^{19}+q^{20}$ \\
       & \begin{tabular}[t]{@{}l@{}}$f^3_{9,5} = q^8+q^9+3q^{10}+4q^{11}+6q^{12}+6q^{13}+7q^{14}+5q^{15}+5q^{16}+3q^{17}+2q^{18}+q^{19}+q^{20}$\end{tabular} \\
       & \begin{tabular}[t]{@{}l@{}}$f^3_{9,6} = q^5 + 2q^6 + 3q^7 + 4q^8 + 6q^9 + 6q^{10} + 6q^{11}+ 6q^{12} + 5q^{13} + 4q^{14} + 3q^{15}+ 2q^{16} + q^{17}+q^{18} $\end{tabular} \\
       & \begin{tabular}[t]{@{}l@{}}$f^3_{9,7} = q^2 + q^3 + 2q^4 + 2q^5 + 3q^6 + 3q^7 + 4q^8+ 3q^9 + 3q^{10} + 2q^{11} + 2q^{12} + q^{13} + q^{14}$\end{tabular} \\
       & $f^3_{9,8} = q + q^2 + q^3 + q^4 + q^5 + q^6 + q^7 + q^8$ \\
       & $f^3_{9,9} = 1$ \\
\hline
$(3,10)$ & $f^3_{10,1} = f^3_{10,2} = f^3_{10,3} = 0$ \\
&$f^3_{10,4} =q^{18}+q^{19}+2q^{20}+2q^{21}+2q^{22}+q^{23}+q^{24}$
\\
        & $f^3_{10,5} = q^{13}+2q^{14}+4q^{15}+5q^{16}+7q^{17}+7q^{18}+7q^{19}+6q^{20}+5q^{21}+3q^{22}+2q^{23}+q^{24}+q^{25}$ \\
        & \begin{tabular}[t]{@{}l@{}}$f^3_{10,6} = q^8+q^9+3q^{10}+4q^{11}+7q^{12}+8q^{13}+10q^{14}+10q^{15}+11q^{16}+9q^{17}+8q^{18}+6q^{19}+5q^{20}$
\\\quad\quad~$+3q^{21}+2q^{22}+q^{23}+q^{24}$\end{tabular} \\
        & \begin{tabular}[t]{@{}l@{}}$f^3_{10,7} = q^5+2q^6+3q^7+4q^8+6q^9+7q^{10}+8q^{11}+8q^{12}+8q^{13}+7q^{14}+7q^{15}+5q^{16}+4q^{17}$
\\\quad\quad~$+3q^{18}+2q^{19}+q^{20}+q^{21}$
\end{tabular} \\
        & \begin{tabular}[t]{@{}l@{}}$f^3_{10,8} = q^2 + q^3 + 2q^4 + 2q^5 + 3q^6 + 3q^7 + 4q^8+ 4q^9 + 4q^{10} + 3q^{11} + 3q^{12} + 2q^{13} + 2q^{14} + q^{15} + q^{16}$\end{tabular} \\
        & $f^3_{10,9} = q + q^2 + q^3 + q^4 + q^5 + q^6 + q^7 + q^8 + q^9$ \\
        & $f^3_{10,10} = 1$ \\
\hline
$(3,11)$ & $f^3_{11,1} = f^3_{11,2} =f^3_{11,3} = 0$ \\
&$ f^3_{11,4} = q^{25}+q^{26}+q^{27}+q^{28}$
\\
& $f^3_{11,5}=q^{18}+q^{19}+3q^{20}+4q^{21}+6q^{22}+6q^{23}+7q^{24}+5q^{25}+5q^{26}+3q^{27}+2q^{28}+q^{29}+q^{30} $
\\
        &\begin{tabular}[t]{@{}l@{}} $f^3_{11,6} = q^{13}+2q^{14}+4q^{15}+6q^{16}+9q^{17}+11q^{18}+13q^{19}+14q^{20}+14q^{21}+13q^{22}+11q^{23}+9q^{24}$\\\quad\quad~$+7q^{25}+5q^{26}+3q^{27}+2q^{28}+q^{29}+q^{30}$\end{tabular} \\
        & \begin{tabular}[t]{@{}l@{}}$f^3_{11,7} = q^8+q^9+3q^{10}+4q^{11}+7q^{12}+9q^{13}+12q^{14}+13q^{15}+16q^{16}+15q^{17}+16q^{18}+14q^{19}$\\\quad\quad~$+13q^{20}+10q^{21}+9q^{22}+6q^{23}+5q^{24}+3q^{25}+2q^{26}+q^{27}+q^{28}$\end{tabular} \\
        & \begin{tabular}[t]{@{}l@{}}$f^3_{11,8} = q^5+2q^6+3q^7+4q^8+6q^9+7q^{10}+9q^{11}+10q^{12}+10q^{13}+10q^{14}+10q^{15}+9q^{16}+8q^{17}$\\\quad\quad~$
+7q^{18}+5q^{19}+4q^{20}+3q^{21}+2q^{22}+q^{23}+q^{24}$\end{tabular} \\
        & \begin{tabular}[t]{@{}l@{}}$f^3_{11,9} = q^2+q^3+2q^4+2q^5+3q^6+3q^7+4q^8+4q^9+5q^{10}+4q^{11}+4q^{12}+3q^{13}+3q^{14}+2q^{15}$\\\quad\quad~$+2q^{16}+q^{17}+q^{18}$\end{tabular} \\
        & $f^3_{11,10} = q + q^2 + q^3 + q^4 + q^5 + q^6 + q^7 + q^8 + q^9 + q^{10}$ \\
        & $f^3_{11,11} = 1$ \\
\hline
$(3,12)$ & $f^3_{12,1} = f^3_{12,2} = f^3_{12,3} =  0$
\\
&$f^3_{12,4} =q^{32} $
\\
&$f^3_{12,5} =q^{25}+2q^{26}+3q^{27}+4q^{28}+5q^{29}+4q^{30}+4q^{31}+3q^{32}+2q^{33}+q^{34}+q^{35}$
\\
&\begin{tabular}[t]{@{}l@{}} $f^3_{12,6} = q^{18}+q^{19}+3q^{20}+5q^{21}+8q^{22}+10q^{23}+14q^{24}+14q^{25}+16q^{26}+15q^{27}+14q^{28}+11q^{29}$\\\quad\quad~$+10q^{30}+7q^{31}+5q^{32}+3q^{33}+2q^{34}+q^{35}+q^{36}$\end{tabular}
\\
        &\begin{tabular}[t]{@{}l@{}} $f^3_{12,7} = q^{13}+2q^{14}+4q^{15}+6q^{16}+10q^{17}+13q^{18}+17q^{19}+20q^{20}+23q^{21}+24q^{22}+25q^{23}+23q^{24}$
\\\quad\quad~$+22q^{25}+19q^{26}+16q^{27}+12q^{28}+10q^{29}+7q^{30}+5q^{31}+3q^{32}+2q^{33}+q^{34}+q^{35}$\end{tabular} \\
        & \begin{tabular}[t]{@{}l@{}}$f^3_{12,8} = q^8+q^9+3q^{10}+4q^{11}+7q^{12}+9q^{13}+13q^{14}+15q^{15}+19q^{16}+20q^{17}+22q^{18}+22q^{19}+23q^{20}$\\\quad\quad~$+20q^{21}+19q^{22}+16q^{23}+14q^{24}+11q^{25}+9q^{26}+6q^{27}+5q^{28}+3q^{29}+2q^{30}+q^{31}+q^{32}$\end{tabular} \\
        & \begin{tabular}[t]{@{}l@{}}$f^3_{12,9} = q^5+2q^6+3q^7+4q^8+6q^9+7q^{10}+9q^{11}+11q^{12}+12q^{13}+12q^{14}+13q^{15}+12q^{16}+12q^{17}$\\\quad\quad~$+11q^{18}+10q^{19}+8q^{20}+7q^{21}+5q^{22}+4q^{23}+3q^{24}+2q^{25}+q^{26}+q^{27}$\end{tabular} \\
        & \begin{tabular}[t]{@{}l@{}}$f^3_{12,10} = q^2+q^3+2q^4+2q^5+3q^6+3q^7+4q^8+4q^9+5q^{10}+5q^{11}+5q^{12}+4q^{13}+4q^{14}+3q^{15}$\\\quad\quad~~$+3q^{16}+2q^{17}+2q^{18}+q^{19}+q^{20}$\end{tabular} \\
        & \begin{tabular}[t]{@{}l@{}}$f^3_{12,11} = q+q^2+q^3+q^4+q^5+q^6+q^7+q^8+q^9+q^{10}+q^{11}$\end{tabular} \\
        & $f^3_{12,12} = 1$ \\
\hline
\end{tabular}

    \caption{Mathematica computations of the polynomials $f^3_{m,l}$ continued.}
    \label{tab:my_label}
\end{table}

% \textcolor{red}{Is the following expression for the polynomials $f^2_{m,m-k}$ appearing in the $(A_1,D_5)$ index consistent with the above data?
% \begin{equation}
%     \begin{split}
%         &f^2_{m,m-k}=q^{k^2}\frac{(q^{m-2k+1};q)_k}{(q;q)_k}, \quad 0\leq 2k\leq m \\
%         &f^2_{m,m-k}=0, \quad 2k>m \\
%     \end{split}
% \end{equation}
% Also, looking at all three cases ($n=1,2,3$), it looks like $f^n_{m,\ell}=0$ for $n\ell<m$ and nonzero for $n\ell\geq m$. Is this true?}

% \subsection*{$n=3$}
\section{Explicit Formula For $f^n_{m,l}(q)$}\label{app:new_id}    
In this appendix, we find an explicit form of the polynomials $f_{m,l}^n(q)$.
Recall that  
\begin{equation}\label{eq:main_id_gen_app}
\frac{\left(z ; q^{2n+1}\right)_{\infty}\left(q^{2n+1} / z ; q^{2n+1}\right)_{\infty}\left(q^{2n+1} ; q^{2n+1}\right)_{\infty}}{(z ; q)_{\infty}(q / z ; q)_{\infty}(q ; q)_{\infty}}=\sum_{m=0}^{\infty}q^m\sum_{l=1}^m \frac{f^n_{m,l}(q)}{(q ; q)_l} \sum_{k=-l}^l{
2 l \choose
l+k}_q z^k~,    
\end{equation}
where $f^n_{m,l}$ is a polynomial satisfying 
\begin{equation}\label{eq:fmnl_def_app}
    1 + \sum_{m=1}^\infty q^m \sum_{l=1}^m 
    \frac{f^n_{m,l}(q)}{(q;q)_l} (q^N; q)_l 
    (q^{1-N};q)_l = q^{nN(N-1)}~.
\end{equation}
Our goal is to find a closed form expression for the polynomials $f_{m,l}^n(q)$. Let us begin by proving a lemma.
\begin{lemma}\label{lemma:S_kqN_rel}
For each $k\in\mathds{N}$, define the sum
\begin{equation}
\begin{split}
    S_k:=\sum_{i_k,i_{k-1},\dots,i_3,i_2,i_1\geq 0}&q^{i_k^2+(i_k+i_{k-1})^2+\dots+(i_k+\dots+i_3+i_2)^2}q^{ki_k+(k-1)i_{k-1}+\dots+2i_2+i_1}
    \\
    \times&\frac{(q^N;q)_{i_k+i_{k-1}+\dots +i_2+i_1}(q^{1-N};q)_{i_k+i_{k-1}+\dots +i_2+i_1}}{(q;q)_{i_k}(q;q)_{i_{k-1}}\dots (q;q)_{i_3}(q;q)_{i_2}(q;q)_{i_1}}~.
\end{split}    
\end{equation}
Then we have 
\begin{equation}
    S_k=q^{kN(N-1)}~.
\end{equation}
\end{lemma}
\begin{proof}
We prove this by induction on $k$. For $k=1$, we have 
\begin{equation}
    S_1=\sum_{i_1\geq 0}q^{i_1}\frac{(q^N;q)_{i_1}(q^{1-N};q)_{i_1}}{(q;q)_{i_1}}~.
\end{equation}
This is the hypergeometric sum in \eqref{eq:S_1_sum} and is equal to $q^{N(N-1)}$. We now assume that $S_{k-1}=q^{(k-1)N(N-1)}$ and show that $S_{k}=q^{kN(N-1)}$. We have 
\begin{equation}
\begin{split}
    S_k=\sum_{i_k,i_{k-1},\dots,i_3,i_2\geq 0}&q^{i_k^2+(i_k+i_{k-1})^2+\dots+(i_k+\dots+i_3+i_2)^2}q^{ki_k+(k-1)i_{k-1}+\dots+2i_2}
    \\
    \times&\frac{(q^N;q)_{i_k+i_{k-1}+\dots +i_2}(q^{1-N};q)_{i_k+i_{k-1}+\dots +i_2}}{(q;q)_{i_k}(q;q)_{i_{k-1}}\dots (q;q)_{i_3}(q;q)_{i_2}}
    \\
    \times&\sum_{i_1\geq 0}q^{i_1}\frac{(q^{N+i_k+i_{k-1}+\dots+i_2};q)_{i_1}(q^{1-N+i_k+i_{k-1}+\dots+i_2};q)_{i_1}}{(q;q)_{i_1}}~,
\end{split}    
\end{equation}
where we used 
\begin{equation}
    (a;q)_{n+k}=(a;q)_n(aq^n;q)_k~.
\end{equation}
The inner sum can again be evaluated using the hypergeometric sum \eqref{eq:phi_identity} to be
\begin{equation}
\begin{split}
\sum_{i_1\geq 0}q^{i_1}\frac{(q^{N+i_k+i_{k-1}+\dots+i_2};q)_{i_1}(q^{1-N+i_k+i_{k-1}+\dots+i_2};q)_{i_1}}{(q;q)_{i_1}}&=q^{(N+i_k+i_{k-1}+\dots+i_2)(N-i_k-i_{k-1}-\dots-i_2-1)}
\\
&=q^{N(N-1)-(i_k+\dots+i_2)^2-i_k-i_{k-1}-\dots-i_2}~.
\end{split}
\end{equation}
Plugging in $S_k$, we get 
\begin{equation}
\begin{split}
    S_k=\sum_{i_k,i_{k-1},\dots,i_3,i_2\geq 0}&q^{i_k^2+(i_k+i_{k-1})^2+\dots+(i_k+\dots+i_3+i_2)^2}q^{ki_k+(k-1)i_{k-1}+\dots+2i_2}
    \\
    \times&\frac{(q^N;q)_{i_k+i_{k-1}+\dots +i_2}(q^{1-N};q)_{i_k+i_{k-1}+\dots +i_2}}{(q;q)_{i_k}(q;q)_{i_{k-1}}\dots (q;q)_{i_3}(q;q)_{i_2}}
    \\
    \times&q^{N(N-1)-(i_k+\dots+i_2)^2-i_k-i_{k-1}-\dots-i_2}
    \\
    =q^{N(N-1)}&\sum_{i_k,i_{k-1},\dots,i_3,i_2\geq 0}q^{i_k^2+(i_k+i_{k-1})^2+\dots+(i_k+\dots+i_3)^2}q^{(k-1)i_k+(k-2)i_{k-1}+\dots+i_2}
    \\
    &\hspace{3cm}\times\frac{(q^N;q)_{i_k+i_{k-1}+\dots +i_2}(q^{1-N};q)_{i_k+i_{k-1}+\dots +i_2}}{(q;q)_{i_k}(q;q)_{i_{k-1}}\dots (q;q)_{i_3}(q;q)_{i_2}}~.
\end{split}    
\end{equation}
Renaming the dummy variables $i_j\to i_{j-1},~2\leq j\leq k$ in the last sum, we see that
\begin{equation}
    S_k=q^{N(N-1)}S_{k-1}=q^{(k-1)N(N-1)}q^{N(N-1)}=q^{kN(N-1)}~.
\end{equation}
\end{proof}
Now in $S_k$ we change variable to  
\begin{equation}
    i_1+i_2+\dots+i_k=:i_1'~.
\end{equation}
Then the sum becomes 
\begin{equation}
\begin{split}
    S_k=\sum_{i_k,i_{k-1},\dots,i_3,i_2,i_1'\geq 0}&q^{Q(i_2,i_3,\dots,i_k)}q^{(k-1)i_k+(k-2)i_{k-1}+\dots+i_2+i_1'}
    \\
    \times&\frac{(q^N;q)_{i'_1}(q^{1-N};q)_{i'_1}}{(q;q)_{i_k}(q;q)_{i_{k-1}}\dots (q;q)_{i_3}(q;q)_{i_2}(q;q)_{i'_1-i_2-\dots-i_k}}~,
\end{split}    
\end{equation}
where we have defined 
\begin{equation}\label{eq:Q_def}
    Q(i_2,i_3,\dots,i_k):=i_k^2+(i_k+i_{k-1})^2+\dots+(i_k+\dots+i_3+i_2)^2~.
\end{equation}
Next we change variables to 
\begin{equation}
    (k-1)i_k+(k-2)i_{k-1}+\dots+i_2+i_1'=:i_2'~.
\end{equation}
Noting that 
\begin{equation}
    \frac{1}{(q;q)_k}=\frac{(q^{k+1};q)_\infty}{(q;q)_\infty}=0,\quad \text{for}\quad k<0~,
\end{equation}
we get 
\begin{equation}
\begin{split}
    S_k&=\sum_{i_k,i_{k-1},\dots,i_3,i_2',i_1'\geq 0}\frac{q^{Q(i_2'-(k-1)i_k-(k-2)i_{k-1}-\dots-i_1',i_3,\dots,i_k)}q^{i_2'}}{(q;q)_{i_k}(q;q)_{i_{k-1}}\dots (q;q)_{i_3}}
    \\
    &\hspace{4cm}\times\frac{(q^N;q)_{i'_1}(q^{1-N};q)_{i'_1}}{(q;q)_{i_2'-(k-1)i_k-(k-2)i_{k-1}-\dots-i_1'}(q;q)_{2i'_1+i_3+2i_4+\dots+(k-2)i_k-i_2'}}
    \\
    &=\sum_{i_2'\geq 0}q^{i_2'}\sum_{i_1'\geq 0}\frac{(q^N;q)_{i'_1}(q^{1-N};q)_{i'_1}}{(q;q)_{i_1'}}
    \\
    &\times\sum_{i_k,i_{k-1},\dots,i_3\geq 0}\frac{q^{Q(i_2'-(k-1)i_k-(k-2)i_{k-1}-\dots-i_1',i_3,\dots,i_k)}(q;q)_{i_1'}}{(q;q)_{i_k}\dots (q;q)_{i_3}(q;q)_{i_2'-(k-1)i_k-(k-2)i_{k-1}-\dots-i_1'}(q;q)_{2i'_1+i_3+2i_4+\dots+(k-2)i_k-i_2'}} ~.
\end{split}      
\end{equation}
Note that the second sum is bounded above by $i_2'$ since 
\begin{equation}
    \frac{1}{(q;q)_{i_2'-(k-1)i_k-(k-2)i_{k-1}-\dots-i_1'}}=0\quad\text{unless}\quad i_1'\leq i_2'~.
\end{equation}
Thus we have 
\begin{equation}
    S_k=\sum_{m=0}^\infty q^{m}\sum_{l= 0}^{m}f^k_{m,l}(q)\frac{(q^N;q)_{l}(q^{1-N};q)_{l}}{(q;q)_{l}}~,
\end{equation}
where 
\begin{equation}\label{eq:fnml_closed_form}
f^k_{m,l}(q)=\sum_{i_k,i_{k-1},\dots,i_3\geq 0}\frac{q^{Q(m-(k-1)i_k-(k-2)i_{k-1}-\dots-2i_3-l,i_3,\dots,i_k)}(q;q)_{l}}{(q;q)_{i_k}\dots (q;q)_{i_3}(q;q)_{m-(k-1)i_k-(k-2)i_{k-1}-\dots-2i_3-l}(q;q)_{2l+i_3+2i_4+\dots+(k-2)i_k-m}}~.    
\end{equation}
By Lemma \ref{lemma:S_kqN_rel}, we see that $f^n_{m,l}$ in \eqref{eq:fnml_closed_form} satisfies \eqref{eq:fmnl_def_app}. 
\subsection{Proof Of Recursion Relation \eqref{eq:rec_fnml} For $k=2$}
In this appendix, we prove that the formula \eqref{eq:fnml_closed_form} satisfies the recursion relation \eqref{eq:rec_fnml} for $k=2$. First note that by definition $f^k_{m,l}=0$ for $m\leq 0$ or $l\leq 0$. For $m=l$, we see that only the term $i_3=i_4=\dots i_k=0$ is nonvanishing since 
\begin{equation}
    \frac{1}{(q;q)_n}=0 \quad\text{for}\quad n<0~.
\end{equation}
Clearly, the term $i_3=i_4=\dots i_k=0$ in the sum \eqref{eq:fnml_closed_form} is 1. For $m=kl$, the only nonvanishing term in the sum is $i_3=i_4=\dots i_{k-1}=0,i_k=l$. This term in the sum is given by 
\begin{equation}
    \frac{q^{(k-1)l^2}(q;q)_l}{(q;q)_l}=q^{(k-1)l^2}~.
\end{equation}
We now prove the final recursion relation
\begin{equation}
f^2_{m,l}=q^l f^{2}_{m-1,l} + f^{2}_{m-1,l-1}-q^{2m-l-3} f^{2}_{m-3,l-1}~.    
\end{equation}
For $k=2$, the sum is empty and we simply have 
\begin{equation}
f^2_{m,l}(q)=q^{(m-l)^2}\frac{(q;q)_{l}}{(q;q)_{m-l}(q;q)_{2l-m}}=q^{(m-l)^2}{l\choose m-l}_q~.    
\end{equation}
% Let us introduce the notation
% \begin{equation}
% \begin{split}
%     a&:=i_k+i_{k-1}+\dots+i_3+m-(k-1)i_k-(k-2)i_{k-1}-\dots-2i_3-l~,
%     \\
%     b&:=m-(k-1)i_k-(k-2)i_{k-1}-\dots-2i_3-l~,
%     \\
%     c&:=2l+i_3+2i_4+\dots+(k-2)i_k-m~.
% \end{split}    
% \end{equation}
% Then we can write 
% \begin{equation}
% \begin{split}
% &q^l f^{k}_{m-1,l} + f^{k}_{m-1,l-1}-q^{2m-l-k-1} f^{k}_{m-k-1,l-1}\\&=\sum_{i_k,i_{k-1},\dots,i_3\geq 0}\frac{q^{Q(i_3,i_4,\dots,i_k)}}{(q;q)_{i_k}\dots (q;q)_{i_3}}\left[\frac{q^{(a-1)^2+l}(q;q)_{l}}{(q;q)_{b-1}(q;q)_{c+1}}+\frac{q^{a^2}(q;q)_{l-1}}{(q;q)_b(q;q)_{c-1}}-\frac{q^{(a-k)^2+2m-l-k-1}(q;q)_{l-1}}{(q;q)_{b-k}(q;q)_{c+k-1}}\right]
% \\&=\sum_{i_k,i_{k-1},\dots,i_3\geq 0}\frac{q^{Q(b,i_3,i_4,\dots,i_k)}(q;q)_l}{(q;q)_{i_k}\dots (q;q)_{i_3}(q;q)_{b}(q;q)_{c}}
% \\&\hspace{3.25cm}\times\left[\frac{q^{1-2a+l}(1-q^b)}{(1-q^{c+1})}+\frac{(1-q^c)}{(1-q^l)}-\frac{q^{k^2-2ak+2m-l-k-1}(1-q^b)}{(1-q^l)}\prod_{i=1}^{k-1}\left(\frac{1-q^{b-k+i}}{1-q^{c+i}}\right)\right]
% \end{split}
% \end{equation}
% We claim that 
% \begin{equation}
% \left[\frac{q^{1-2a+l}(1-q^b)}{(1-q^{c+1})}+\frac{(1-q^c)}{(1-q^l)}-\frac{q^{k^2-2ak+2m-l-k-1}(1-q^b)}{(1-q^l)}\prod_{i=1}^{k-1}\left(\frac{1-q^{b-k+i}}{1-q^{c+i}}\right)\right]=1~.    
% \end{equation}
We have 
\begin{equation}
\begin{split}
q^l f^{2}_{m-1,l} &+ f^{2}_{m-1,l-1}-q^{2m-l-3} f^{2}_{m-3,l-1}\\&=\frac{q^{(m-l-1)^2+l}(q;q)_{l}}{(q;q)_{m-l-1}(q;q)_{2l-m+1}}+ \frac{q^{(m-l)^2}(q;q)_{l-1}}{(q;q)_{m-l}(q;q)_{2l-m-1}}-\frac{q^{(m-l-2)^2+2m-l-3}(q;q)_{l-1}}{(q;q)_{m-l-2}(q;q)_{2l-m+1}}
\\&=\frac{q^{(m-l)^2}(q;q)_{l}}{(q;q)_{m-l}(q;q)_{2l-m}}\left[q^{l-2(m-l)+1}\frac{1-q^{m-l}}{1-q^{2l-m+1}}+\frac{1-q^{m-2l}}{1-q^l}\right.\\&\left.\hspace{6cm}-q^{-2m+3l+1}\frac{(1-q^{m-l-1})(1-q^{m-l})}{(1-q^l)(1-q^{2l-m+1})}\right]~.
\end{split}
\end{equation}
Straightforward calculation gives
\begin{equation}
q^{l-2(m-l)+1}\frac{1-q^{m-l}}{1-q^{2l-m+1}}+\frac{1-q^{m-2l}}{1-q^l}-q^{-2m+3l+1}\frac{(1-q^{m-l-1})(1-q^{m-l})}{(1-q^l)(1-q^{2l-m+1})}= 1~.   
\end{equation}
\end{appendix}
% \section{The Hilbert Series At $t^2$}
% We want to show that the Hilbert series to order $t^2$ is given by 
% \begin{equation}\label{eq:Hil_ser_t2}
%     1+\frac{t}{(q;q)_1}\sum_{k=0}^2{2\choose k}_q+\frac{t^2}{(q;q)_2}\sum_{k=0}^4{4\choose k}_q~.
% \end{equation}
% The starting algebraic variety is 
% \begin{equation}
% \label{arcspacea1a3}
%     R=\C[X,Y,Z,J]/\langle f_1,f_2,f_3,f_4,f_5\rangle~,
% \end{equation}
% where 
% \begin{equation}
%     f_1:=XZ-Y^2\quad f_2:=JX,\quad f_3:=JY,\quad f_4:=JZ,\quad f_5:=J^2~.
% \end{equation}


\bibliography{mtc}
\end{document}

\section{The Hilbert Series from Arc space constraints counting}
A general formula for the Hilbert Series of Arc space
\eqref{arcspacea1a3} is
given by 
\begin{equation}
\label{gen Hilbert Series}
    H_R(q, t) = 
    \frac{1-\sum_{m=2}^\infty\sum_{n=0}^\infty \sum_{l=0}^\infty(-1)^lF^{(l)}_{n,m}q^nt^m}{(1-t)^3\prod_{i=1}^\infty(1-q^i t)^4},
\end{equation}
where $F^{(l)}_{n,m}$ is the number of generators in the $l$-th depth syzygy module of degree $(n,m)$ depends on the choice of ordering, here we choose

\begin{equation}
J_1>J_2>\cdots>X_0>X_1>\cdots>Y_0>Y_1>\cdots>Z_0>Z_1>\ldots
\end{equation}
but $ \sum_{l=0}^\infty(-1)^lF^{(l)}_{n,m}$ is independent of the choice of ordering.


For $m = 2$,
\begin{equation}
\begin{aligned}
    F^{(0)}_{n,m= 2} &= 
    C_z(n,1) + C_{xyzj}(n-1,1),\\
    F^{(l)}_{n,m= 2} &= 
    0,\quad \text{for}\quad l >0.
\end{aligned}
\end{equation}
These will give the result for $m = 2$
in Appendix C.

For $m=3$,
\begin{equation}
\label{m=3 case}
    \begin{aligned}
    F^{(0)}_{n,m= 3} =& 
    C_z(n-1,2) + C_{xyzj}(n-2,2),\\
    F^{(1)}_{n,m= 3} =& 
    C_z(n,2)+2C_z(n-1,1)+C_z(n-2,1)\\&+C_{xyzj}(n-1,2)+2C_{xyzj}(n-3,1)-i_n,\\
    F^{(2)}_{n,m= 3} =& 
    C_z(n-1,1) + C_{xyzj}(n-3,1), \\
    F^{(l)}_{n,m= 3} =& 
    0,\quad \text{for}\quad l >2,
\end{aligned}
\end{equation}
where $i_n = 1$ if $n$ is even, $i_n = 3$ if $n = 1$
and $i_n = 4$ otherwise.

For generic $m$
\begin{equation}
    F_{n, m}^{(0)}=C_z(n-m +2,m-1)+C_{x y z j}(n-m+1,m-1),
\end{equation}
\begin{equation}
    \dots
\end{equation}

We now focus on the case when $m = 3$. Let us denote by $N_n$, the coefficient of $q^n t^3$ in the expansion of \eqref{gen Hilbert Series} around $t=0, q=0$. From \eqref{m=3 case}, these coefficients can be expressed as follows
\begin{equation}
    \begin{aligned}
        N_0 =& 7, \\
        N_1 =& 12, \\
        N_2 =& 27, \\
        N_3 =& 48, \\
        N_n =& 41-20n +64 \bigg\lfloor\frac{6+(n-3)^2}{12}\bigg\rfloor +
        20 \bigg\lfloor
        \frac{-1+n}{3}
        \bigg\rfloor\\
         &+
        58 \bigg\lfloor
        \frac{-1+n}{2}
        \bigg\rfloor
        -20
        \bigg\lfloor
        \frac{n}{3}
        \bigg\rfloor
        +30 
        \bigg\lfloor
        \frac{n}{2}
        \bigg\rfloor\quad\quad\quad\quad\quad \text{for}\quad n\geq4.
    \end{aligned}
\end{equation}
We separate the coefficients into 12 groups by taking $n = 12 k + i$, where $i=4, 5,\dots 15,\quad k=0\dots$. All floor functions are then removed
\begin{equation}
\begin{aligned}
    N_{12k+4} &= 768 k^2+416 k+79\\
    N_{12k+5} &= 768 k^2+544 k+117\\
    N_{12k+6} &= 768 k^2+672 k+171\\
    N_{12k+7} &= 768 k^2+800 k+229\\
    N_{12k+8} &= 768 k^2+928 k+303\\
    N_{12k+9} &=768 k^2+1056 k+385\\
    N_{12k+10} & = 768 k^2+1184 k+479\\
    N_{12k+11} & =768 k^2+1312 k+581\\
    N_{12k+12} & = 768 k^2+1440 k+699\\
    N_{12k+13} & = 768 k^2+1568 k+821\\
    N_{12k+14} & = 768 k^2+1696 k+959\\ 
    N_{12k+15} & =768 k^2+1824 k+1105
\end{aligned}
\end{equation}
We denote 
\begin{equation}
   N_{12k +i} = a_i k^2 + b_i k +c_i, 
\end{equation}
The $i$-th sector can be summed over 
\begin{equation}
     \sum_{k=0}^\infty 
    N_{12k+i} q^{12k +i} = 
    q^i \frac{a_i q^{12}\left(1+q^{12}\right)+b_i q^{12}\left(1-q^{12}\right)+c_i\left(1-q^{12}\right)^2}{\left(1-q^{12}\right)^3}
\end{equation}
By straightforward computation, we can show
\begin{equation}
\begin{aligned}
    &7+12q + 27 q^2 + 48 q^3 + \sum_{i=4}^{15} q^i \frac{a_i q^{12}\left(1+q^{12}\right)+b_i q^{12}\left(1-q^{12}\right)+c_i\left(1-q^{12}\right)^2}{\left(1-q^{12}\right)^3}\\
&=(1-q^4)(1-q^5)(1-q^6)\sum_{k=0}^{6} \frac{1}{(q ; q)_k(q ; q)_{6-k}}.
\end{aligned}
\end{equation}
\section{Hilbert Series from Arc space $(A_1,A_3)$}
We start from the algebraic variety
\begin{equation}
R=\mathbb{C}[X, Y, Z, J] /\left\langle f_1, f_2, f_3, f_4, f_5\right\rangle,
\end{equation}
where
\begin{equation}
f_1:=X Z-Y^2 \quad f_2:=J X, \quad f_3:=J Y, \quad f_4:=J Z, \quad f_5:=J^2.
\end{equation}
To define the arc space, we start with the power series
\begin{equation}
X(x)=\sum_{n=0}^{\infty} X_n x^n, \quad Y(x)=\sum_{n=0}^{\infty} Y_n x^n, \quad Z(x)=\sum_{n=0}^{\infty} Z_n x^n, \quad J(x)=\sum_{n=1}^{\infty} J_n x^n.
\end{equation}
Note that the power series for $J$ starts from $n=1$.
We define the $n$-jet scheme of algebraic variety $R_n$ by
\begin{equation}
R_n:=\mathbb{C}\left[X_0, \ldots, X_n, Y_0, \ldots, Y_n, Z_0, \ldots, Z_n, J_1, \ldots, J_n\right] /\left\langle f_1(x), \ldots, f_5(x) \equiv 0 \bmod x^{n+1}\right\rangle,
\end{equation}
where
\begin{equation}
f_i(x):=f_i(X(x), Y(x), Z(x), J(x)), \quad i=1, \ldots, 5.
\end{equation}


\noindent{\textbf{The first order in $t$:}}


\noindent{\textbf{The second order in $t$:}}


\noindent{\textbf{The third order in $t$:}}

% \subsection*{$n=1$}
% \[
% \begin{aligned}
% &\text{For } m=1 \text{ case:} \\
% &f^1_{1,1}=1 \\
% &\text{For } m=2 \text{ case:} \\
% &f^1_{2,1}=0, \quad f^1_{2,2}=1 \\
% &\text{For } m=3 \text{ case:} \\
% &f^1_{3,1}=0, \quad f^1_{3,2}=0, \quad f^1_{3,3}=1 \\
% &\text{For } m=4 \text{ case:} \\
% &f^1_{4,1}=0, \quad f^1_{4,2}=0, \quad f^1_{4,3}=0, \quad f^1_{4,4}=1 \\
% &\text{For } m=5 \text{ case:} \\
% &f^1_{5,1}=0, \quad f^1_{5,2}=0, \quad f^1_{5,3}=0, \quad f^1_{5,4}=0, \quad f^1_{5,5}=1 \\
% &\text{For } m=6 \text{ case:} \\
% &f^1_{6,1}=0, \quad f^1_{6,2}=0, \quad f^1_{6,3}=0, \quad f^1_{6,4}=0, \quad f^1_{6,5}=0, \quad f^1_{6,6}=1 \\
% &\text{For } m=7 \text{ case:} \\
% &f^1_{7,1}=0, \quad f^1_{7,2}=0, \quad f^1_{7,3}=0, \quad f^1_{7,4}=0, \quad f^1_{7,5}=0, \quad f^1_{7,6}=0, \quad f^1_{7,7}=1 \\
% &\text{For } m=8 \text{ case:} \\
% &f^1_{8,1}=0, \quad f^1_{8,2}=0, \quad f^1_{8,3}=0, \quad f^1_{8,4}=0, \quad f^1_{8,5}=0, \quad f^1_{8,6}=0, \quad f^1_{8,7}=0, \quad f^1_{8,8}=1 \\
% &\text{For } m=9 \text{ case:} \\
% &f^1_{9,1}=0, \quad f^1_{9,2}=0, \quad f^1_{9,3}=0, \quad f^1_{9,4}=0, \quad f^1_{9,5}=0, \quad f^1_{9,6}=0, \quad f^1_{9,7}=0, \quad f^1_{9,8}=0,
% \\
% &f^1_{9,9}=1 \\
% &\text{For } m=10 \text{ case:} \\
% &f^1_{10,1}=0, \quad f^1_{10,2}=0, \quad f^1_{10,3}=0, \quad f^1_{10,4}=0, \quad f^1_{10,5}=0, \quad f^1_{10,6}=0, \quad f^1_{10,7}=0, \\ 
% &f^1_{10,8}=0, \quad f^1_{10,9}=0, \quad f^1_{10,10}=1 \\
% &\text{For } m=11 \text{ case:} \\
% &f^1_{11,1}=0, \quad f^1_{11,2}=0, \quad f^1_{11,3}=0, \quad f^1_{11,4}=0, \quad f^1_{11,5}=0, \quad f^1_{11,6}=0, \quad f^1_{11,7}=0, \\ 
% &f^1_{11,8}=0, \quad f^1_{11,9}=0, \quad f^1_{11,10}=0, \quad f^1_{11,11}=1 \\
% &\text{For } m=12 \text{ case:} \\
% &f^1_{12,1}=0, \quad f^1_{12,2}=0, \quad f^1_{12,3}=0, \quad f^1_{12,4}=0, \quad f^1_{12,5}=0, \quad f^1_{12,6}=0, \quad f^1_{12,7}=0, \\ &f^1_{12,8}=0, \quad f^1_{12,9}=0, \quad f^1_{12,10}=0, \quad f^1_{12,11}=0, \quad f^1_{12,12}=1 \\
% \end{aligned}
% \]

Table of functions

$\underline{\boldsymbol{n=2.}}$
\[
\begin{aligned}
&\text{For } m=1 \text{ case:} \\
&f^2_{1,1}=1 \\
&\text{For } m=2 \text{ case:} \\
&f^2_{2,1}=q, \quad f^2_{2,2}=1 \\
&\text{For } m=3 \text{ case:} \\
&f^2_{3,1}=0, \quad f^2_{3,2}=q+q^2, \quad f^2_{3,3}=1 \\
&\text{For } m=4 \text{ case:} \\
&f^2_{4,1}=0, \quad f^2_{4,2}=q^4, \quad f^2_{4,3}=q+q^2+q^3, \quad f^2_{4,4}=1 \\
&\text{For } m=5 \text{ case:} \\
&f^2_{5,1}=0, \quad f^2_{5,2}=0, \quad f^2_{5,3}=q^4+q^5+q^6, \quad f^2_{5,4}=q+q^2+q^3+q^4, \quad f^2_{5,5}=1 \\
&\text{For } m=6 \text{ case:} \\
&f^2_{6,1}=0, \quad f^2_{6,2}=0, \quad f^2_{6,3}=q^9, \\
&f^2_{6,4}=q^4+q^5+2q^6+q^7+q^8, \\
&f^2_{6,5}=q+q^2+q^3+q^4+q^5, \quad f^2_{6,6}=1 \\
&\text{For } m=7 \text{ case:} \\
&f^2_{7,1}=0, \quad f^2_{7,2}=0, \quad f^2_{7,3}=0, \\
&f^2_{7,4}=q^9+q^{10}+q^{11}+q^{12}, \\
&f^2_{7,5}=q^4+q^5+2q^6+2q^7+2q^8+q^9+q^{10}, \\
&f^2_{7,6}=q+q^2+q^3+q^4+q^5+q^6, \quad f^2_{7,7}=1 \\
&\text{For } m=8 \text{ case:} \\
&f^2_{8,1}=0, \quad f^2_{8,2}=0, \quad f^2_{8,3}=0, \\
&f^2_{8,4}=q^{16}, \\
&f^2_{8,5}=q^9+q^{10}+2q^{11}+2q^{12}+2q^{13}+q^{14}+q^{15}, \\
&f^2_{8,6}=q^4+q^5+2q^6+2q^7+3q^8+2q^9+2q^{10}+q^{11}+q^{12}, \\
&f^2_{8,7}=q+q^2+q^3+q^4+q^5+q^6+q^7, \quad f^2_{8,8}=1 \\
&\text{For } m=9 \text{ case:} \\
&f^2_{9,1}=0, \quad f^2_{9,2}=0, \quad f^2_{9,3}=0, \quad f^2_{9,4}=0, \\
&f^2_{9,5}=q^{16}+q^{17}+q^{18}+q^{19}+q^{20}, \\
&f^2_{9,6}=q^9+q^{10}+2q^{11}+3q^{12}+3q^{13}+3q^{14}+3q^{15}+2q^{16}+q^{17}+q^{18}, \\
&f^2_{9,7}=q^4+q^5+2q^6+2q^7+3q^8+3q^9+3q^{10}+2q^{11}+2q^{12}+q^{13}+q^{14}, \\
&f^2_{9,8}=q+q^2+q^3+q^4+q^5+q^6+q^7+q^8, \quad f^2_{9,9}=1 \\
&\text{For } m=10 \text{ case:} \\
&f^2_{10,1}=0, \quad f^2_{10,2}=0, \quad f^2_{10,3}=0, \quad f^2_{10,4}=0, \\
&f^2_{10,5}=q^{25}, \\
&f^2_{10,6}=q^{16}+q^{17}+2q^{18}+2q^{19}+3q^{20}+2q^{21}+2q^{22}+q^{23}+q^{24}, \\
&f^2_{10,7}=q^9+q^{10}+2q^{11}+3q^{12}+4q^{13}+4q^{14}+5q^{15}+4q^{16}+4q^{17}+3q^{18}+2q^{19}+q^{20}+q^{21}, \\
\end{aligned}
\]
\[
\begin{aligned}
&f^2_{10,8}=q^4+q^5+2q^6+2q^7+3q^8+3q^9+4q^{10}+3q^{11}+3q^{12}+2q^{13}+2q^{14}+q^{15}+q^{16},
\\
&f^2_{10,9}=q+q^2+q^3+q^4+q^5+q^6+q^7+q^8+q^9, \quad f^2_{10,10}=1~.
\end{aligned}
\]

\[
\begin{aligned}
&\text{For } m=11 \text{ case:} \\
&f^2_{11,1}=0, \quad f^2_{11,2}=0, \quad f^2_{11,3}=0, \quad f^2_{11,4}=0, \quad f^2_{11,5}=0, \\
&f^2_{11,6}=q^{25}+q^{26}+q^{27}+q^{28}+q^{29}+q^{30}, \\
&f^2_{11,7}=q^{16}+q^{17}+2q^{18}+3q^{19}+4q^{20}+4q^{21}+5q^{22}+4q^{23}+4q^{24}+3q^{25}+2q^{26}+q^{27}+q^{28}, \\
&f^2_{11,8}=q^9+q^{10}+2q^{11}+3q^{12}+4q^{13}+5q^{14}+6q^{15}+6q^{16}+6q^{17}+6q^{18}+5q^{19}
\\&\quad\quad+4q^{20}+3q^{21}+2q^{22}+q^{23}+q^{24}, \\
&f^2_{11,9}=q^4+q^5+2q^6+2q^7+3q^8+3q^9+4q^{10}+4q^{11}+4q^{12}+3q^{13}+3q^{14}+2q^{15}+2q^{16}+q^{17}+q^{18}, \\
&f^2_{11,10}=q+q^2+q^3+q^4+q^5+q^6+q^7+q^8+q^9+q^{10}, \quad f^2_{11,11}=1 \\
&\text{For } m=12 \text{ case:} \\
&f^2_{12,1}=0, \quad f^2_{12,2}=0, \quad f^2_{12,3}=0, \quad f^2_{12,4}=0, \quad f^2_{12,5}=0, \\
&f^2_{12,6}=q^{36}, \\
&f^2_{12,7}=q^{25}+q^{26}+2q^{27}+2q^{28}+3q^{29}+3q^{30}+3q^{31}+2q^{32}+2q^{33}+q^{34}+q^{35}, \\
&f^2_{12,8}=q^{16}+q^{17}+2q^{18}+3q^{19}+5q^{20}+5q^{21}+7q^{22}+7q^{23}+8q^{24}+7q^{25}+7q^{26}\\
&\quad\quad+5q^{27}+5q^{28}+3q^{29}+2q^{30}+q^{31}+q^{32}, \\
&f^2_{12,9}=q^9+q^{10}+2q^{11}+3q^{12}+4q^{13}+5q^{14}+7q^{15}+7q^{16}+8q^{17}+8q^{18}+8q^{19}+7q^{20}\\&\quad\quad+7q^{21}+5q^{22}+4q^{23}+3q^{24}+2q^{25}+q^{26}+q^{27}, \\
&f^2_{12,10}=q^4+q^5+2q^6+2q^7+3q^8+3q^9+4q^{10}+4q^{11}+5q^{12}+4q^{13}+4q^{14}+3q^{15}+3q^{16}\\&\quad\quad+2q^{17}+2q^{18}+q^{19}+q^{20}, \\
&f^2_{12,11}=q+q^2+q^3+q^4+q^5+q^6+q^7+q^8+q^9+q^{10}+q^{11}, \quad f^2_{12,12}=1
\end{aligned}
\]

$\underline{\boldsymbol{n=3.}}$
\[
\begin{aligned}
&\text{For } m=1 \text{ case:} \\
&f^3_{1,1}=1 \\
&\text{For } m=2 \text{ case:} \\
&f^3_{2,1}=q, \quad f^3_{2,2}=1 \\
&\text{For } m=3 \text{ case:} \\
&f^3_{3,1}=q^2, \quad f^3_{3,2}=q+q^2, \quad f^3_{3,3}=1 \\
&\text{For } m=4 \text{ case:} \\
&f^3_{4,1}=0, \quad f^3_{4,2}=q^2+q^3+q^4, \quad f^3_{4,3}=q+q^2+q^3, \quad f^3_{4,4}=1 \\
&\text{For } m=5 \text{ case:} \\
&f^3_{5,1}=0, \quad f^3_{5,2}=q^5+q^6, \quad f^3_{5,3}=q^2+q^3+2q^4+q^5+q^6, \\
&f^3_{5,4}=q+q^2+q^3+q^4, \quad f^3_{5,5}=1 \\
&\text{For } m=6 \text{ case:} \\
&f^3_{6,1}=0, \quad f^3_{6,2}=q^8, \\
&f^3_{6,3}=q^5+2q^6+2q^7+q^8+q^9, \\
&f^3_{6,4}=q^2+q^3+2q^4+2q^5+2q^6+q^7+q^8, \\
&f^3_{6,5}=q+q^2+q^3+q^4+q^5, \quad f^3_{6,6}=1 \\
&\text{For } m=7 \text{ case:} \\
&f^3_{7,1}=0, \quad f^3_{7,2}=0, \\
&f^3_{7,3}=q^8+q^9+2q^{10}+q^{11}+q^{12}, \\
&f^3_{7,4}=q^5+2q^6+3q^7+3q^8+3q^9+2q^{10}+q^{11}+q^{12}, \\
&f^3_{7,5}=q^2+q^3+2q^4+2q^5+3q^6+2q^7+2q^8+q^9+q^{10}, \\
&f^3_{7,6}=q+q^2+q^3+q^4+q^5+q^6, \quad f^3_{7,7}=1 \\
&\text{For } m=8 \text{ case:} \\
&f^3_{8,1}=0, \quad f^3_{8,2}=0, \\
&f^3_{8,3}=q^{13}+q^{14}+q^{15}, \\
&f^3_{8,4}=q^8+q^9+3q^{10}+3q^{11}+4q^{12}+3q^{13}+2q^{14}+q^{15}+q^{16}, \\
&f^3_{8,5}=q^5+2q^6+3q^7+4q^8+5q^9+4q^{10}+4q^{11}+3q^{12}+2q^{13}+q^{14}+q^{15}, \\
&f^3_{8,6}=q^2+q^3+2q^4+2q^5+3q^6+3q^7+3q^8+2q^9+2q^{10}+q^{11}+q^{12}, \\
&f^3_{8,7}=q+q^2+q^3+q^4+q^5+q^6+q^7, \quad f^3_{8,8}=1. 
\end{aligned}
\]

\[
\begin{aligned}
&\text{For } m=9 \text{ case:} \\
&f^3_{9,1}=0, \quad f^3_{9,2}=0, \quad f^3_{9,3}=q^{18}, \\
&f^3_{9,4}=q^{13}+2q^{14}+3q^{15}+3q^{16}+3q^{17}+2q^{18}+q^{19}+q^{20}, \\
&f^3_{9,5}=q^8+q^9+3q^{10}+4q^{11}+6q^{12}+6q^{13}+7q^{14}+5q^{15}+5q^{16}+3q^{17}+2q^{18}+q^{19}+q^{20}, \\
&f^3_{9,6}=q^5+2q^6+3q^7+4q^8+6q^9+6q^{10}+6q^{11}+6q^{12}+5q^{13}+4q^{14}+3q^{15}+2q^{16}+q^{17}+q^{18}, \\
&f^3_{9,7}=q^2+q^3+2q^4+2q^5+3q^6+3q^7+4q^8+3q^9+3q^{10}+2q^{11}+2q^{12}+q^{13}+q^{14}, \\
&f^3_{9,8}=q+q^2+q^3+q^4+q^5+q^6+q^7+q^8, \quad f^3_{9,9}=1 \\
&\text{For } m=10 \text{ case:} \\
&f^3_{10,1}=0, \quad f^3_{10,2}=0, \quad f^3_{10,3}=0, \\
&f^3_{10,4}=q^{18}+q^{19}+2q^{20}+2q^{21}+2q^{22}+q^{23}+q^{24}, \\
&f^3_{10,5}=q^{13}+2q^{14}+4q^{15}+5q^{16}+7q^{17}+7q^{18}+7q^{19}+6q^{20}+5q^{21}+3q^{22}+2q^{23}+q^{24}+q^{25}, \\
&f^3_{10,6}=q^8+q^9+3q^{10}+4q^{11}+7q^{12}+8q^{13}+10q^{14}+10q^{15}+11q^{16}+9q^{17}+8q^{18}+6q^{19}+5q^{20}
\\&\quad\quad+3q^{21}+2q^{22}+q^{23}+q^{24}, \\
&f^3_{10,7}=q^5+2q^6+3q^7+4q^8+6q^9+7q^{10}+8q^{11}+8q^{12}+8q^{13}+7q^{14}+7q^{15}+5q^{16}+4q^{17}
\\&\quad\quad+3q^{18}+2q^{19}+q^{20}+q^{21}, \\
&f^3_{10,8}=q^2+q^3+2q^4+2q^5+3q^6+3q^7+4q^8+4q^9+4q^{10}+3q^{11}+3q^{12}+2q^{13}+2q^{14}+q^{15}+q^{16}, \\
&f^3_{10,9}=q+q^2+q^3+q^4+q^5+q^6+q^7+q^8+q^9, \quad f^3_{10,10}=1 \\
&\text{For } m=11 \text{ case:} \\
&f^3_{11,1}=0, \quad f^3_{11,2}=0, \quad f^3_{11,3}=0, \\
&f^3_{11,4}=q^{25}+q^{26}+q^{27}+q^{28}, \\
&f^3_{11,5}=q^{18}+q^{19}+3q^{20}+4q^{21}+6q^{22}+6q^{23}+7q^{24}+5q^{25}+5q^{26}+3q^{27}+2q^{28}+q^{29}+q^{30}, \\
&f^3_{11,6}=q^{13}+2q^{14}+4q^{15}+6q^{16}+9q^{17}+11q^{18}+13q^{19}+14q^{20}+14q^{21}+13q^{22}+11q^{23}\\&\quad\quad+9q^{24}+7q^{25}+5q^{26}+3q^{27}+2q^{28}+q^{29}+q^{30}, \\
&f^3_{11,7}=q^8+q^9+3q^{10}+4q^{11}+7q^{12}+9q^{13}+12q^{14}+13q^{15}+16q^{16}+15q^{17}+16q^{18}+14q^{19}\\&\quad\quad+13q^{20}+10q^{21}+9q^{22}+6q^{23}+5q^{24}+3q^{25}+2q^{26}+q^{27}+q^{28}, \\
&f^3_{11,8}=q^5+2q^6+3q^7+4q^8+6q^9+7q^{10}+9q^{11}+10q^{12}+10q^{13}+10q^{14}+10q^{15}+9q^{16}\\&\quad\quad
+8q^{17}+7q^{18}+5q^{19}+4q^{20}+3q^{21}+2q^{22}+q^{23}+q^{24}, \\
&f^3_{11,9}=q^2+q^3+2q^4+2q^5+3q^6+3q^7+4q^8+4q^9+5q^{10}+4q^{11}+4q^{12}+3q^{13}+3q^{14}\\&\quad\quad+2q^{15}+2q^{16}+q^{17}+q^{18}, \\
&f^3_{11,10}=q+q^2+q^3+q^4+q^5+q^6+q^7+q^8+q^9+q^{10}, \quad f^3_{11,11}=1 
\end{aligned}\]
\[
\begin{aligned}
&\text{For } m=12 \text{ case:} \\
&f^3_{12,1}=0, \quad f^3_{12,2}=0, \quad f^3_{12,3}=0, \\
&f^3_{12,4}=q^{32}, \\
&f^3_{12,5}=q^{25}+2q^{26}+3q^{27}+4q^{28}+5q^{29}+4q^{30}+4q^{31}+3q^{32}+2q^{33}+q^{34}+q^{35}, \\
&f^3_{12,6}=q^{18}+q^{19}+3q^{20}+5q^{21}+8q^{22}+10q^{23}+14q^{24}+14q^{25}+16q^{26}+15q^{27}+14q^{28}\\&\quad\quad+11q^{29}+10q^{30}+7q^{31}+5q^{32}+3q^{33}+2q^{34}+q^{35}+q^{36}, \\
&f^3_{12,7}=q^{13}+2q^{14}+4q^{15}+6q^{16}+10q^{17}+13q^{18}+17q^{19}+20q^{20}+23q^{21}+24q^{22}\\&\quad\quad
+25q^{23}+23q^{24}+22q^{25}+19q^{26}+16q^{27}+12q^{28}+10q^{29}+7q^{30}+5q^{31}+3q^{32}+2q^{33}+q^{34}+q^{35}, \\
&f^3_{12,8}=q^8+q^9+3q^{10}+4q^{11}+7q^{12}+9q^{13}+13q^{14}+15q^{15}+19q^{16}+20q^{17}+22q^{18}+22q^{19}\\&\quad\quad+23q^{20}+20q^{21}+19q^{22}+16q^{23}+14q^{24}+11q^{25}+9q^{26}+6q^{27}+5q^{28}+3q^{29}+2q^{30}+q^{31}+q^{32}, \\
&f^3_{12,9}=q^5+2q^6+3q^7+4q^8+6q^9+7q^{10}+9q^{11}+11q^{12}+12q^{13}+12q^{14}+13q^{15}+12q^{16}\\&\quad\quad+12q^{17}+11q^{18}+10q^{19}+8q^{20}+7q^{21}+5q^{22}+4q^{23}+3q^{24}+2q^{25}+q^{26}+q^{27}, \\
&f^3_{12,10}=q^2+q^3+2q^4+2q^5+3q^6+3q^7+4q^8+4q^9+5q^{10}+5q^{11}+5q^{12}+4q^{13}+4q^{14}+3q^{15}\\&\quad\quad+3q^{16}+2q^{17}+2q^{18}+q^{19}+q^{20}, \\
&f^3_{12,11}=q+q^2+q^3+q^4+q^5+q^6+q^7+q^8+q^9+q^{10}+q^{11}, \quad f^3_{12,12}=1
\end{aligned}
\]
